\documentclass[11pt,a4paper]{book}

% BEGIN preamble.tex
\usepackage[T1]{fontenc}
\usepackage{lmodern}
\usepackage{amsmath,amssymb,amsthm,mathtools}
\usepackage[a4paper,margin=27mm,headheight=14pt]{geometry}
\usepackage{microtype}
\usepackage{enumitem}
\usepackage{xcolor}
\usepackage{longtable,booktabs,array}
\usepackage[colorlinks=true,linkcolor=blue!45!black,urlcolor=blue!45!black]{hyperref}
\usepackage{fancyhdr}
\pagestyle{fancy}\fancyhf{}\fancyhead[LE,RO]{\thepage}
\fancyhead[LO]{\small D05 Unified Notes}\fancyhead[RE]{\small\nouppercase{\leftmark}}
\makeatletter
\renewcommand{\chaptermark}[1]{\markboth{\if@mainmatter\thechapter.\ \fi#1}{}}
\renewcommand{\l@section}{\@dottedtocline{1}{1.5em}{3.2em}}
\makeatother
\setlength{\parindent}{0pt}\setlength{\parskip}{5pt}
\setlength{\emergencystretch}{3em}
\allowdisplaybreaks[2]
\raggedbottom
\newtheorem{theorem}{Theorem}[chapter]
\newtheorem{lemma}[theorem]{Lemma}
\newtheorem{proposition}[theorem]{Proposition}
\newtheorem{corollary}[theorem]{Corollary}
\theoremstyle{definition}\newtheorem{definition}[theorem]{Definition}
\newtheorem{example}[theorem]{Example}
\theoremstyle{remark}\newtheorem{remark}[theorem]{Remark}
\newcommand{\K}{\mathbb K}\newcommand{\R}{\mathbb R}\newcommand{\C}{\mathbb C}
\newcommand{\N}{\mathbb N}\newcommand{\Z}{\mathbb Z}
\DeclareMathOperator{\Hom}{Hom}\DeclareMathOperator{\End}{End}
\DeclareMathOperator{\im}{im}\DeclareMathOperator{\rank}{rank}
\DeclareMathOperator{\tr}{tr}\DeclareMathOperator{\Span}{span}
\DeclareMathOperator{\diag}{diag}\DeclareMathOperator{\id}{id}
\DeclareMathOperator{\Sym}{Sym}\DeclareMathOperator{\Ann}{Ann}
\DeclareMathOperator{\Pf}{pf}\DeclareMathOperator{\codim}{codim}
\DeclareMathOperator{\GL}{GL}
\newcommand{\ip}[2]{\langle #1,#2\rangle}
\newcommand{\norm}[1]{\lVert #1\rVert}
\newcommand{\abs}[1]{\lvert #1\rvert}
\newcommand{\dual}{^{\vee}}
\newcommand{\adj}{^{\dagger}}
\newcommand{\pinv}{^{+}}
\newcommand{\kref}[1]{\hyperref[k:#1]{\texttt{K05.\detokenize{#1}}}}
\newcommand{\Kentry}[2]{\section{#2}\label{k:#1}\noindent\textsf{\small K05.\detokenize{#1}}\par
  \begingroup\small\color{black!65}\textit{Source route: }\csname source@#1\endcsname\par\endgroup}

% BEGIN source-metadata.tex
% Generated source routes from the preserved D05 registry.
\expandafter\def\csname source@01.system\endcsname{F04:1.2 (print p.~3; PDF p.~12)}
\expandafter\def\csname source@01.row_ops\endcsname{F04:1.3 (print p.~6; PDF p.~15)}
\expandafter\def\csname source@01.solution_invariance\endcsname{F04:1.3 (print p.~6; PDF p.~15)}
\expandafter\def\csname source@01.echelon\endcsname{F04:1.4 (print p.~11; PDF p.~20)}
\expandafter\def\csname source@01.elimination\endcsname{F04:1.3 (print p.~6; PDF p.~15);\allowbreak\ F04:1.4 (print p.~11; PDF p.~20)}
\expandafter\def\csname source@01.rref_unique\endcsname{F04:2.5 (print p.~55; PDF p.~64)}
\expandafter\def\csname source@01.consistency\endcsname{F04:1.4 (print p.~11; PDF p.~20);\allowbreak\ F04:2.6 (print p.~58; PDF p.~67)}
\expandafter\def\csname source@01.parametrization\endcsname{F04:1.4 (print p.~11; PDF p.~20);\allowbreak\ F04:2.6 (print p.~58; PDF p.~67)}
\expandafter\def\csname source@01.affine_solutions\endcsname{F04:1.2 (print p.~3; PDF p.~12);\allowbreak\ F05:3B (print p.~59; PDF p.~73)}
\expandafter\def\csname source@01.pivot_rank\endcsname{F04:2.5 (print p.~55; PDF p.~64);\allowbreak\ F04:2.6 (print p.~58; PDF p.~67)}
\expandafter\def\csname source@01.inverse_algorithm\endcsname{F04:1.6 (print p.~21; PDF p.~30)}
\expandafter\def\csname source@02.vector_space\endcsname{F04:2.1 (print p.~28; PDF p.~37);\allowbreak\ F05:1B (print p.~12; PDF p.~26)}
\expandafter\def\csname source@02.subspace\endcsname{F04:2.2 (print p.~34; PDF p.~43);\allowbreak\ F05:1C (print p.~18; PDF p.~32)}
\expandafter\def\csname source@02.span\endcsname{F04:2.2 (print p.~34; PDF p.~43);\allowbreak\ F05:2A (print p.~28; PDF p.~42)}
\expandafter\def\csname source@02.span_properties\endcsname{F04:2.2 (print p.~34; PDF p.~43);\allowbreak\ F05:2A (print p.~28; PDF p.~42)}
\expandafter\def\csname source@02.independence\endcsname{F04:2.3 (print p.~40; PDF p.~49);\allowbreak\ F05:2A (print p.~28; PDF p.~42)}
\expandafter\def\csname source@02.exchange\endcsname{F04:2.3 (print p.~40; PDF p.~49);\allowbreak\ F05:2A (print p.~28; PDF p.~42)}
\expandafter\def\csname source@02.basis\endcsname{F04:2.3 (print p.~40; PDF p.~49);\allowbreak\ F05:2B (print p.~39; PDF p.~53)}
\expandafter\def\csname source@02.basis_extension\endcsname{F04:2.3 (print p.~40; PDF p.~49);\allowbreak\ F05:2B (print p.~39; PDF p.~53)}
\expandafter\def\csname source@02.dimension\endcsname{F04:2.3 (print p.~40; PDF p.~49);\allowbreak\ F05:2C (print p.~44; PDF p.~58)}
\expandafter\def\csname source@02.subspace_dimension\endcsname{F04:2.3 (print p.~40; PDF p.~49);\allowbreak\ F05:2C (print p.~44; PDF p.~58)}
\expandafter\def\csname source@02.sum_intersection\endcsname{F04:2.2 (print p.~34; PDF p.~43);\allowbreak\ F05:1C (print p.~18; PDF p.~32)}
\expandafter\def\csname source@02.grassmann\endcsname{F04:2.3 (print p.~40; PDF p.~49);\allowbreak\ F05:2C (print p.~44; PDF p.~58)}
\expandafter\def\csname source@02.examples\endcsname{F04:2.1 (print p.~28; PDF p.~37);\allowbreak\ F04:2.3 (print p.~40; PDF p.~49);\allowbreak\ F05:2B (print p.~39; PDF p.~53)}
\expandafter\def\csname source@03.linear_map\endcsname{F04:3.1 (print p.~67; PDF p.~76);\allowbreak\ F05:3A (print p.~52; PDF p.~66)}
\expandafter\def\csname source@03.map_algebra\endcsname{F04:3.2 (print p.~74; PDF p.~83);\allowbreak\ F05:3A (print p.~52; PDF p.~66)}
\expandafter\def\csname source@03.basis_extension\endcsname{F04:3.1 (print p.~67; PDF p.~76);\allowbreak\ F05:3A (print p.~52; PDF p.~66)}
\expandafter\def\csname source@03.kernel\endcsname{F04:3.1 (print p.~67; PDF p.~76);\allowbreak\ F05:3B (print p.~59; PDF p.~73)}
\expandafter\def\csname source@03.image\endcsname{F04:3.1 (print p.~67; PDF p.~76);\allowbreak\ F05:3B (print p.~59; PDF p.~73)}
\expandafter\def\csname source@03.injective_surjective\endcsname{F04:3.1 (print p.~67; PDF p.~76);\allowbreak\ F05:3B (print p.~59; PDF p.~73)}
\expandafter\def\csname source@03.rank_nullity\endcsname{F04:3.1 (print p.~67; PDF p.~76);\allowbreak\ F05:3B (print p.~59; PDF p.~73)}
\expandafter\def\csname source@03.isomorphism\endcsname{F04:3.3 (print p.~84; PDF p.~93);\allowbreak\ F05:3D (print p.~82; PDF p.~96)}
\expandafter\def\csname source@03.dimension_classification\endcsname{F04:3.3 (print p.~84; PDF p.~93);\allowbreak\ F05:3D (print p.~82; PDF p.~96)}
\expandafter\def\csname source@03.equal_dimension\endcsname{F04:3.3 (print p.~84; PDF p.~93);\allowbreak\ F05:3D (print p.~82; PDF p.~96)}
\expandafter\def\csname source@03.composition_rank\endcsname{F04:3.2 (print p.~74; PDF p.~83);\allowbreak\ F05:3B (print p.~59; PDF p.~73)}
\expandafter\def\csname source@03.first_iso\endcsname{F04:A.4 (print p.~394; PDF p.~403);\allowbreak\ F05:3E (print p.~96; PDF p.~110)}
\expandafter\def\csname source@03.hom_dimension\endcsname{F04:3.2 (print p.~74; PDF p.~83);\allowbreak\ F05:3D (print p.~82; PDF p.~96)}
\expandafter\def\csname source@04.coordinates\endcsname{F04:2.4 (print p.~49; PDF p.~58);\allowbreak\ F05:3D (print p.~82; PDF p.~96)}
\expandafter\def\csname source@04.matrix_of_map\endcsname{F04:3.4 (print p.~86; PDF p.~95);\allowbreak\ F05:3C (print p.~69; PDF p.~83)}
\expandafter\def\csname source@04.matrix_product\endcsname{F04:1.5 (print p.~16; PDF p.~25);\allowbreak\ F05:3C (print p.~69; PDF p.~83)}
\expandafter\def\csname source@04.composition\endcsname{F04:3.4 (print p.~86; PDF p.~95);\allowbreak\ F05:3C (print p.~69; PDF p.~83)}
\expandafter\def\csname source@04.inverse_matrix\endcsname{F04:1.6 (print p.~21; PDF p.~30);\allowbreak\ F05:3D (print p.~82; PDF p.~96)}
\expandafter\def\csname source@04.change_basis\endcsname{F04:3.4 (print p.~86; PDF p.~95);\allowbreak\ F05:3D (print p.~82; PDF p.~96)}
\expandafter\def\csname source@04.similarity\endcsname{F04:3.4 (print p.~86; PDF p.~95);\allowbreak\ F04:6.3 (print p.~190; PDF p.~199);\allowbreak\ F05:3D (print p.~82; PDF p.~96)}
\expandafter\def\csname source@04.equivalence\endcsname{F04:2.6 (print p.~58; PDF p.~67);\allowbreak\ F04:3.4 (print p.~86; PDF p.~95)}
\expandafter\def\csname source@04.blocks\endcsname{F04:6.7 (print p.~213; PDF p.~222);\allowbreak\ F05:8B (print p.~308; PDF p.~322)}
\expandafter\def\csname source@04.rank_bridge\endcsname{F04:3.7 (print p.~111; PDF p.~120);\allowbreak\ F05:3C (print p.~69; PDF p.~83);\allowbreak\ F05:3F (print p.~105; PDF p.~119)}
\expandafter\def\csname source@04.trace\endcsname{F05:8D (print p.~326; PDF p.~340);\allowbreak\ F04:3.4 (print p.~86; PDF p.~95)}
\expandafter\def\csname source@04.rank_factorization\endcsname{F05:3C (print p.~69; PDF p.~83)}
\expandafter\def\csname source@05.alternating\endcsname{F04:5.2 (print p.~141; PDF p.~150);\allowbreak\ F05:9B (print p.~346; PDF p.~360)}
\expandafter\def\csname source@05.det_definition\endcsname{F04:5.2 (print p.~141; PDF p.~150);\allowbreak\ F05:9C (print p.~354; PDF p.~368)}
\expandafter\def\csname source@05.det_uniqueness\endcsname{F04:5.3 (print p.~150; PDF p.~159);\allowbreak\ F05:9B (print p.~346; PDF p.~360);\allowbreak\ F05:9C (print p.~354; PDF p.~368)}
\expandafter\def\csname source@05.det_row_rules\endcsname{F04:5.2 (print p.~141; PDF p.~150);\allowbreak\ F04:5.4 (print p.~156; PDF p.~165);\allowbreak\ F05:9C (print p.~354; PDF p.~368)}
\expandafter\def\csname source@05.det_product\endcsname{F04:5.4 (print p.~156; PDF p.~165);\allowbreak\ F05:9C (print p.~354; PDF p.~368)}
\expandafter\def\csname source@05.det_invertible\endcsname{F04:5.4 (print p.~156; PDF p.~165);\allowbreak\ F05:9C (print p.~354; PDF p.~368)}
\expandafter\def\csname source@05.cofactor\endcsname{F04:5.4 (print p.~156; PDF p.~165);\allowbreak\ F05:9C (print p.~354; PDF p.~368)}
\expandafter\def\csname source@05.adjugate\endcsname{F04:5.4 (print p.~156; PDF p.~165)}
\expandafter\def\csname source@05.cramer\endcsname{F04:5.4 (print p.~156; PDF p.~165)}
\expandafter\def\csname source@05.triangular_block\endcsname{F04:5.4 (print p.~156; PDF p.~165);\allowbreak\ F05:9C (print p.~354; PDF p.~368)}
\expandafter\def\csname source@05.operator_det\endcsname{F05:9C (print p.~354; PDF p.~368);\allowbreak\ F06:5 (print p.~96; PDF p.~103)}
\expandafter\def\csname source@05.orientation\endcsname{F05:9C (print p.~354; PDF p.~368)}
\expandafter\def\csname source@05.volume\endcsname{F05:7F (print p.~280; PDF p.~294);\allowbreak\ F05:9C (print p.~354; PDF p.~368)}
\expandafter\def\csname source@05.hadamard\endcsname{F05:9C (print p.~354; PDF p.~368)}
\expandafter\def\csname source@05.vandermonde\endcsname{F05:9C (print p.~354; PDF p.~368);\allowbreak\ F04:4.3 (print p.~124; PDF p.~133)}
\expandafter\def\csname source@06.quotient\endcsname{F04:A.4 (print p.~394; PDF p.~403);\allowbreak\ F05:3E (print p.~96; PDF p.~110)}
\expandafter\def\csname source@06.quotient_map\endcsname{F04:A.4 (print p.~394; PDF p.~403);\allowbreak\ F05:3E (print p.~96; PDF p.~110)}
\expandafter\def\csname source@06.quotient_lift\endcsname{F04:A.4 (print p.~394; PDF p.~403);\allowbreak\ F05:3E (print p.~96; PDF p.~110)}
\expandafter\def\csname source@06.quotient_dimension\endcsname{F04:A.4 (print p.~394; PDF p.~403);\allowbreak\ F05:3E (print p.~96; PDF p.~110)}
\expandafter\def\csname source@06.direct_sum\endcsname{F04:6.6 (print p.~209; PDF p.~218);\allowbreak\ F05:1C (print p.~18; PDF p.~32)}
\expandafter\def\csname source@06.external_sum\endcsname{F05:3E (print p.~96; PDF p.~110);\allowbreak\ F04:6.6 (print p.~209; PDF p.~218);\allowbreak\ F06:3 (print p.~60; PDF p.~67)}
\expandafter\def\csname source@06.complement\endcsname{F04:A.4 (print p.~394; PDF p.~403);\allowbreak\ F05:3E (print p.~96; PDF p.~110)}
\expandafter\def\csname source@06.projection\endcsname{F04:6.6 (print p.~209; PDF p.~218);\allowbreak\ F05:3E (print p.~96; PDF p.~110)}
\expandafter\def\csname source@06.sum_dimension\endcsname{F04:6.6 (print p.~209; PDF p.~218);\allowbreak\ F05:3E (print p.~96; PDF p.~110)}
\expandafter\def\csname source@06.invariant\endcsname{F04:6.4 (print p.~198; PDF p.~207);\allowbreak\ F05:5A (print p.~133; PDF p.~147)}
\expandafter\def\csname source@06.invariant_decomposition\endcsname{F04:6.7 (print p.~213; PDF p.~222);\allowbreak\ F05:8B (print p.~308; PDF p.~322)}
\expandafter\def\csname source@06.noninvariant_complement\endcsname{F04:6.4 (print p.~198; PDF p.~207);\allowbreak\ F04:6.7 (print p.~213; PDF p.~222)}
\expandafter\def\csname source@07.dual\endcsname{F04:3.5 (print p.~97; PDF p.~106);\allowbreak\ F05:3F (print p.~105; PDF p.~119)}
\expandafter\def\csname source@07.dual_basis\endcsname{F04:3.5 (print p.~97; PDF p.~106);\allowbreak\ F05:3F (print p.~105; PDF p.~119)}
\expandafter\def\csname source@07.annihilator\endcsname{F04:3.5 (print p.~97; PDF p.~106);\allowbreak\ F05:3F (print p.~105; PDF p.~119)}
\expandafter\def\csname source@07.annihilator_dimension\endcsname{F04:3.5 (print p.~97; PDF p.~106);\allowbreak\ F05:3F (print p.~105; PDF p.~119)}
\expandafter\def\csname source@07.dual_map\endcsname{F04:3.7 (print p.~111; PDF p.~120);\allowbreak\ F05:3F (print p.~105; PDF p.~119)}
\expandafter\def\csname source@07.dual_kernel_image\endcsname{F04:3.7 (print p.~111; PDF p.~120);\allowbreak\ F05:3F (print p.~105; PDF p.~119)}
\expandafter\def\csname source@07.dual_rank\endcsname{F04:3.7 (print p.~111; PDF p.~120);\allowbreak\ F05:3F (print p.~105; PDF p.~119)}
\expandafter\def\csname source@07.biddual\endcsname{F04:3.6 (print p.~107; PDF p.~116);\allowbreak\ F05:3F (print p.~105; PDF p.~119)}
\expandafter\def\csname source@07.double_annihilator\endcsname{F04:3.5 (print p.~97; PDF p.~106);\allowbreak\ F04:3.6 (print p.~107; PDF p.~116);\allowbreak\ F05:3F (print p.~105; PDF p.~119)}
\expandafter\def\csname source@07.dual_matrix\endcsname{F04:3.7 (print p.~111; PDF p.~120);\allowbreak\ F05:3F (print p.~105; PDF p.~119)}
\expandafter\def\csname source@07.quotient_dual\endcsname{F04:A.4 (print p.~394; PDF p.~403);\allowbreak\ F04:3.5 (print p.~97; PDF p.~106);\allowbreak\ F05:3E (print p.~96; PDF p.~110);\allowbreak\ F05:3F (print p.~105; PDF p.~119)}
\expandafter\def\csname source@07.hyperplanes\endcsname{F04:3.6 (print p.~107; PDF p.~116)}
\expandafter\def\csname source@07.functional_constraints\endcsname{F04:3.6 (print p.~107; PDF p.~116)}
\expandafter\def\csname source@08.evaluation\endcsname{F04:4.2 (print p.~119; PDF p.~128);\allowbreak\ F04:6.3 (print p.~190; PDF p.~199);\allowbreak\ F05:5A (print p.~133; PDF p.~147)}
\expandafter\def\csname source@08.annihilating\endcsname{F04:6.3 (print p.~190; PDF p.~199);\allowbreak\ F05:5B (print p.~143; PDF p.~157)}
\expandafter\def\csname source@08.minpoly\endcsname{F04:6.3 (print p.~190; PDF p.~199);\allowbreak\ F05:5B (print p.~143; PDF p.~157)}
\expandafter\def\csname source@08.charpoly\endcsname{F04:6.2 (print p.~182; PDF p.~191);\allowbreak\ F05:9C (print p.~354; PDF p.~368)}
\expandafter\def\csname source@08.cayley_hamilton\endcsname{F04:6.3 (print p.~190; PDF p.~199);\allowbreak\ F05:9C (print p.~354; PDF p.~368)}
\expandafter\def\csname source@08.eigen_roots\endcsname{F04:6.2 (print p.~182; PDF p.~191);\allowbreak\ F04:6.3 (print p.~190; PDF p.~199);\allowbreak\ F05:5B (print p.~143; PDF p.~157)}
\expandafter\def\csname source@08.bezout_decomposition\endcsname{F04:6.8 (print p.~219; PDF p.~228)}
\expandafter\def\csname source@08.primary\endcsname{F04:6.8 (print p.~219; PDF p.~228)}
\expandafter\def\csname source@08.block_polynomials\endcsname{F04:6.7 (print p.~213; PDF p.~222);\allowbreak\ F04:6.8 (print p.~219; PDF p.~228)}
\expandafter\def\csname source@08.polynomial_inverse\endcsname{F04:6.3 (print p.~190; PDF p.~199);\allowbreak\ F05:5B (print p.~143; PDF p.~157)}
\expandafter\def\csname source@08.coefficients\endcsname{F05:8D (print p.~326; PDF p.~340);\allowbreak\ F05:9C (print p.~354; PDF p.~368)}
\expandafter\def\csname source@08.interpolation\endcsname{F04:4.3 (print p.~124; PDF p.~133)}
\expandafter\def\csname source@08.restriction_minpoly\endcsname{F05:5B (print p.~143; PDF p.~157);\allowbreak\ F05:5D (print p.~163; PDF p.~177);\allowbreak\ F04:6.4 (print p.~198; PDF p.~207)}
\expandafter\def\csname source@08.newton_identities\endcsname{F06:7.7 (print p.~186; PDF p.~193);\allowbreak\ F06:7.8 (print p.~187; PDF p.~194)}
\expandafter\def\csname source@09.eigenspace\endcsname{F04:6.2 (print p.~182; PDF p.~191);\allowbreak\ F05:5A (print p.~133; PDF p.~147)}
\expandafter\def\csname source@09.distinct_independence\endcsname{F04:6.2 (print p.~182; PDF p.~191);\allowbreak\ F05:5A (print p.~133; PDF p.~147)}
\expandafter\def\csname source@09.diagonalizable\endcsname{F04:6.2 (print p.~182; PDF p.~191);\allowbreak\ F05:5D (print p.~163; PDF p.~177)}
\expandafter\def\csname source@09.eigenspace_criterion\endcsname{F04:6.2 (print p.~182; PDF p.~191);\allowbreak\ F05:5D (print p.~163; PDF p.~177)}
\expandafter\def\csname source@09.multiplicities\endcsname{F04:6.2 (print p.~182; PDF p.~191);\allowbreak\ F05:8B (print p.~308; PDF p.~322)}
\expandafter\def\csname source@09.minpoly_criterion\endcsname{F04:6.4 (print p.~198; PDF p.~207);\allowbreak\ F05:5D (print p.~163; PDF p.~177)}
\expandafter\def\csname source@09.distinct_roots\endcsname{F04:6.2 (print p.~182; PDF p.~191);\allowbreak\ F05:5D (print p.~163; PDF p.~177)}
\expandafter\def\csname source@09.commutant_invariance\endcsname{F04:6.5 (print p.~206; PDF p.~215);\allowbreak\ F05:5E (print p.~175; PDF p.~189)}
\expandafter\def\csname source@09.simultaneous_diagonal\endcsname{F04:6.5 (print p.~206; PDF p.~215);\allowbreak\ F05:5E (print p.~175; PDF p.~189)}
\expandafter\def\csname source@09.field_dependence\endcsname{F04:6.2 (print p.~182; PDF p.~191);\allowbreak\ F05:5D (print p.~163; PDF p.~177)}
\expandafter\def\csname source@09.existence\endcsname{F05:5B (print p.~143; PDF p.~157)}
\expandafter\def\csname source@09.gershgorin\endcsname{F05:5D (print p.~163; PDF p.~177)}
\expandafter\def\csname source@09.commuting_spectrum\endcsname{F05:5E (print p.~175; PDF p.~189)}
\expandafter\def\csname source@09.diagonal_dominance\endcsname{F05:5D (print p.~163; PDF p.~177)}
\expandafter\def\csname source@10.triangularization\endcsname{F04:6.4 (print p.~198; PDF p.~207);\allowbreak\ F05:5C (print p.~154; PDF p.~168)}
\expandafter\def\csname source@10.simultaneous_triangular\endcsname{F04:6.5 (print p.~206; PDF p.~215);\allowbreak\ F05:5E (print p.~175; PDF p.~189)}
\expandafter\def\csname source@10.nilpotent\endcsname{F04:7.3 (print p.~244; PDF p.~253);\allowbreak\ F05:8A (print p.~298; PDF p.~312)}
\expandafter\def\csname source@10.kernel_stabilization\endcsname{F05:8A (print p.~298; PDF p.~312);\allowbreak\ F04:6.8 (print p.~219; PDF p.~228)}
\expandafter\def\csname source@10.generalized\endcsname{F04:7.3 (print p.~244; PDF p.~253);\allowbreak\ F05:8A (print p.~298; PDF p.~312);\allowbreak\ F05:8B (print p.~308; PDF p.~322)}
\expandafter\def\csname source@10.generalized_decomposition\endcsname{F04:7.3 (print p.~244; PDF p.~253);\allowbreak\ F05:8B (print p.~308; PDF p.~322)}
\expandafter\def\csname source@10.jordan_chains\endcsname{F04:7.3 (print p.~244; PDF p.~253);\allowbreak\ F05:8C (print p.~319; PDF p.~333)}
\expandafter\def\csname source@10.jordan_existence\endcsname{F04:7.3 (print p.~244; PDF p.~253);\allowbreak\ F05:8C (print p.~319; PDF p.~333)}
\expandafter\def\csname source@10.jordan_uniqueness\endcsname{F04:7.3 (print p.~244; PDF p.~253);\allowbreak\ F05:8C (print p.~319; PDF p.~333)}
\expandafter\def\csname source@10.cyclic\endcsname{F04:7.1 (print p.~227; PDF p.~236)}
\expandafter\def\csname source@10.companion\endcsname{F04:7.1 (print p.~227; PDF p.~236);\allowbreak\ F04:7.2 (print p.~231; PDF p.~240)}
\expandafter\def\csname source@10.rational\endcsname{F04:7.2 (print p.~231; PDF p.~240)}
\expandafter\def\csname source@10.invariant_factors\endcsname{F04:7.2 (print p.~231; PDF p.~240);\allowbreak\ F04:7.4 (print p.~251; PDF p.~260)}
\expandafter\def\csname source@10.cyclic_test\endcsname{F04:7.2 (print p.~231; PDF p.~240)}
\expandafter\def\csname source@10.semisimple\endcsname{F04:7.5 (print p.~262; PDF p.~271)}
\expandafter\def\csname source@10.additive_decomposition\endcsname{F04:6.8 (print p.~219; PDF p.~228);\allowbreak\ F05:8B (print p.~308; PDF p.~322)}
\expandafter\def\csname source@10.semisimple_decomposition\endcsname{F04:7.5 (print p.~262; PDF p.~271)}
\expandafter\def\csname source@10.invariant_factor_computation\endcsname{F04:7.4 (print p.~251; PDF p.~260)}
\expandafter\def\csname source@10.complex_square_roots\endcsname{F05:8C (print p.~319; PDF p.~333);\allowbreak\ F04:7.3 (print p.~244; PDF p.~253)}
\expandafter\def\csname source@10.transpose_similarity\endcsname{F04:7.3 (print p.~244; PDF p.~253)}
\expandafter\def\csname source@10.fitting\endcsname{F05:8A (print p.~298; PDF p.~312);\allowbreak\ F04:6.8 (print p.~219; PDF p.~228)}
\expandafter\def\csname source@11.bilinear\endcsname{F04:10.1 (print p.~359; PDF p.~368);\allowbreak\ F05:9A (print p.~333; PDF p.~347)}
\expandafter\def\csname source@11.form_matrix\endcsname{F04:10.1 (print p.~359; PDF p.~368);\allowbreak\ F05:9A (print p.~333; PDF p.~347)}
\expandafter\def\csname source@11.radicals\endcsname{F04:10.1 (print p.~359; PDF p.~368)}
\expandafter\def\csname source@11.nondegenerate\endcsname{F04:10.1 (print p.~359; PDF p.~368)}
\expandafter\def\csname source@11.symmetric_quadratic\endcsname{F04:10.2 (print p.~367; PDF p.~376);\allowbreak\ F05:9A (print p.~333; PDF p.~347)}
\expandafter\def\csname source@11.polarization\endcsname{F04:10.2 (print p.~367; PDF p.~376);\allowbreak\ F05:9A (print p.~333; PDF p.~347)}
\expandafter\def\csname source@11.diagonal_congruence\endcsname{F04:10.2 (print p.~367; PDF p.~376);\allowbreak\ F05:9A (print p.~333; PDF p.~347)}
\expandafter\def\csname source@11.inertia\endcsname{F04:10.2 (print p.~367; PDF p.~376);\allowbreak\ F05:9A (print p.~333; PDF p.~347)}
\expandafter\def\csname source@11.complex_classification\endcsname{F04:10.2 (print p.~367; PDF p.~376)}
\expandafter\def\csname source@11.alternating\endcsname{F04:10.3 (print p.~375; PDF p.~384);\allowbreak\ F05:9B (print p.~346; PDF p.~360)}
\expandafter\def\csname source@11.symplectic_blocks\endcsname{F04:10.3 (print p.~375; PDF p.~384)}
\expandafter\def\csname source@11.hermitian\endcsname{F04:9.2 (print p.~320; PDF p.~329);\allowbreak\ F04:9.4 (print p.~332; PDF p.~341)}
\expandafter\def\csname source@11.hermitian_inertia\endcsname{F04:9.4 (print p.~332; PDF p.~341)}
\expandafter\def\csname source@11.definiteness\endcsname{F04:9.3 (print p.~325; PDF p.~334);\allowbreak\ F04:10.2 (print p.~367; PDF p.~376)}
\expandafter\def\csname source@11.principal_minors\endcsname{F04:9.3 (print p.~325; PDF p.~334)}
\expandafter\def\csname source@11.isometry_groups\endcsname{F04:10.4 (print p.~379; PDF p.~388);\allowbreak\ F04:8.4 (print p.~299; PDF p.~308)}
\expandafter\def\csname source@11.form_operator\endcsname{F04:9.2 (print p.~320; PDF p.~329)}
\expandafter\def\csname source@11.hermitian_diagonal\endcsname{F04:9.2 (print p.~320; PDF p.~329)}
\expandafter\def\csname source@11.pfaffian\endcsname{F06:8.4 (print p.~200; PDF p.~207);\allowbreak\ F06:8.5 (print p.~202; PDF p.~209)}
\expandafter\def\csname source@12.inner_product\endcsname{F04:8.1 (print p.~270; PDF p.~279);\allowbreak\ F05:6A (print p.~182; PDF p.~196)}
\expandafter\def\csname source@12.convention_bridge\endcsname{F04:8.1 (print p.~270; PDF p.~279);\allowbreak\ F05:6A (print p.~182; PDF p.~196)}
\expandafter\def\csname source@12.norm\endcsname{F04:8.2 (print p.~277; PDF p.~286);\allowbreak\ F05:6A (print p.~182; PDF p.~196)}
\expandafter\def\csname source@12.cauchy_schwarz\endcsname{F04:8.2 (print p.~277; PDF p.~286);\allowbreak\ F05:6A (print p.~182; PDF p.~196)}
\expandafter\def\csname source@12.norm_laws\endcsname{F04:8.2 (print p.~277; PDF p.~286);\allowbreak\ F05:6A (print p.~182; PDF p.~196)}
\expandafter\def\csname source@12.orthonormal\endcsname{F04:8.2 (print p.~277; PDF p.~286);\allowbreak\ F05:6B (print p.~197; PDF p.~211)}
\expandafter\def\csname source@12.gram_schmidt\endcsname{F04:8.2 (print p.~277; PDF p.~286);\allowbreak\ F05:6B (print p.~197; PDF p.~211)}
\expandafter\def\csname source@12.on_basis\endcsname{F04:8.2 (print p.~277; PDF p.~286);\allowbreak\ F05:6B (print p.~197; PDF p.~211)}
\expandafter\def\csname source@12.fourier_parseval\endcsname{F04:8.2 (print p.~277; PDF p.~286);\allowbreak\ F05:6B (print p.~197; PDF p.~211);\allowbreak\ F05:6C (print p.~211; PDF p.~225)}
\expandafter\def\csname source@12.orthogonal_complement\endcsname{F04:8.2 (print p.~277; PDF p.~286);\allowbreak\ F05:6C (print p.~211; PDF p.~225)}
\expandafter\def\csname source@12.orthogonal_decomposition\endcsname{F04:8.2 (print p.~277; PDF p.~286);\allowbreak\ F05:6C (print p.~211; PDF p.~225)}
\expandafter\def\csname source@12.projection\endcsname{F04:8.2 (print p.~277; PDF p.~286);\allowbreak\ F05:6C (print p.~211; PDF p.~225)}
\expandafter\def\csname source@12.best_approximation\endcsname{F04:8.2 (print p.~277; PDF p.~286);\allowbreak\ F05:6C (print p.~211; PDF p.~225)}
\expandafter\def\csname source@12.riesz\endcsname{F04:8.3 (print p.~290; PDF p.~299);\allowbreak\ F05:6B (print p.~197; PDF p.~211)}
\expandafter\def\csname source@12.parallelogram\endcsname{F05:6A (print p.~182; PDF p.~196);\allowbreak\ F04:8.2 (print p.~277; PDF p.~286)}
\expandafter\def\csname source@13.adjoint\endcsname{F04:8.3 (print p.~290; PDF p.~299);\allowbreak\ F05:7A (print p.~228; PDF p.~242)}
\expandafter\def\csname source@13.adjoint_laws\endcsname{F04:8.3 (print p.~290; PDF p.~299);\allowbreak\ F05:7A (print p.~228; PDF p.~242)}
\expandafter\def\csname source@13.adjoint_kernel\endcsname{F04:8.3 (print p.~290; PDF p.~299);\allowbreak\ F05:7A (print p.~228; PDF p.~242);\allowbreak\ F05:7E (print p.~270; PDF p.~284)}
\expandafter\def\csname source@13.self_adjoint\endcsname{F04:8.5 (print p.~311; PDF p.~320);\allowbreak\ F05:7A (print p.~228; PDF p.~242)}
\expandafter\def\csname source@13.unitary\endcsname{F04:8.4 (print p.~299; PDF p.~308);\allowbreak\ F05:7D (print p.~258; PDF p.~272)}
\expandafter\def\csname source@13.self_adjoint_eigenvalues\endcsname{F04:8.5 (print p.~311; PDF p.~320);\allowbreak\ F05:7A (print p.~228; PDF p.~242)}
\expandafter\def\csname source@13.real_spectral\endcsname{F04:8.5 (print p.~311; PDF p.~320);\allowbreak\ F05:7B (print p.~243; PDF p.~257)}
\expandafter\def\csname source@13.complex_spectral\endcsname{F04:8.5 (print p.~311; PDF p.~320);\allowbreak\ F05:7B (print p.~243; PDF p.~257)}
\expandafter\def\csname source@13.spectral_projections\endcsname{F04:9.5 (print p.~335; PDF p.~344);\allowbreak\ F05:7B (print p.~243; PDF p.~257)}
\expandafter\def\csname source@13.algebraic_projections\endcsname{F04:6.7 (print p.~213; PDF p.~222);\allowbreak\ F04:9.5 (print p.~335; PDF p.~344)}
\expandafter\def\csname source@13.functional_calculus\endcsname{F04:9.5 (print p.~335; PDF p.~344)}
\expandafter\def\csname source@13.simultaneous_spectral\endcsname{F04:9.5 (print p.~335; PDF p.~344);\allowbreak\ F05:7B (print p.~243; PDF p.~257)}
\expandafter\def\csname source@13.real_normal_blocks\endcsname{F04:9.6 (print p.~349; PDF p.~358)}
\expandafter\def\csname source@13.unitary_spectrum\endcsname{F04:9.5 (print p.~335; PDF p.~344)}
\expandafter\def\csname source@13.schur\endcsname{F05:6B (print p.~197; PDF p.~211);\allowbreak\ F04:8.5 (print p.~311; PDF p.~320)}
\expandafter\def\csname source@13.normal_reducing\endcsname{F04:9.6 (print p.~349; PDF p.~358)}
\expandafter\def\csname source@13.normal_classification\endcsname{F04:9.6 (print p.~349; PDF p.~358)}
\expandafter\def\csname source@13.normal\endcsname{F05:7A (print p.~228; PDF p.~242);\allowbreak\ F04:8.5 (print p.~311; PDF p.~320)}
\expandafter\def\csname source@13.normal_basic\endcsname{F05:7A (print p.~228; PDF p.~242);\allowbreak\ F04:8.5 (print p.~311; PDF p.~320)}
\expandafter\def\csname source@13.simultaneous_schur\endcsname{F05:6B (print p.~197; PDF p.~211);\allowbreak\ F04:6.5 (print p.~206; PDF p.~215)}
\expandafter\def\csname source@14.positive\endcsname{F04:9.3 (print p.~325; PDF p.~334);\allowbreak\ F05:7C (print p.~251; PDF p.~265)}
\expandafter\def\csname source@14.positive_spectrum\endcsname{F04:9.3 (print p.~325; PDF p.~334);\allowbreak\ F05:7C (print p.~251; PDF p.~265)}
\expandafter\def\csname source@14.positive_root\endcsname{F04:9.5 (print p.~335; PDF p.~344);\allowbreak\ F05:7C (print p.~251; PDF p.~265)}
\expandafter\def\csname source@14.singular_values\endcsname{F05:7E (print p.~270; PDF p.~284)}
\expandafter\def\csname source@14.svd\endcsname{F05:7E (print p.~270; PDF p.~284)}
\expandafter\def\csname source@14.pseudoinverse\endcsname{F05:6C (print p.~211; PDF p.~225)}
\expandafter\def\csname source@14.penrose\endcsname{F05:6C (print p.~211; PDF p.~225);\allowbreak\ F05:7E (print p.~270; PDF p.~284)}
\expandafter\def\csname source@14.least_squares\endcsname{F05:6C (print p.~211; PDF p.~225)}
\expandafter\def\csname source@14.normal_equations\endcsname{F04:8.3 (print p.~290; PDF p.~299);\allowbreak\ F05:6C (print p.~211; PDF p.~225);\allowbreak\ F05:7E (print p.~270; PDF p.~284)}
\expandafter\def\csname source@14.operator_norm\endcsname{F05:7F (print p.~280; PDF p.~294)}
\expandafter\def\csname source@14.low_rank\endcsname{F05:7F (print p.~280; PDF p.~294)}
\expandafter\def\csname source@14.frobenius\endcsname{F05:7E (print p.~270; PDF p.~284);\allowbreak\ F05:7F (print p.~280; PDF p.~294);\allowbreak\ F05:8D (print p.~326; PDF p.~340)}
\expandafter\def\csname source@14.polar\endcsname{F04:9.5 (print p.~335; PDF p.~344);\allowbreak\ F05:7F (print p.~280; PDF p.~294)}
\expandafter\def\csname source@14.qr_cholesky\endcsname{F05:7D (print p.~258; PDF p.~272)}
\expandafter\def\csname source@14.svd_pseudoinverse\endcsname{F05:7E (print p.~270; PDF p.~284)}
\expandafter\def\csname source@14.ellipsoids\endcsname{F05:7F (print p.~280; PDF p.~294)}
\expandafter\def\csname source@14.eigenvalue_frobenius_bound\endcsname{F05:8B (print p.~308; PDF p.~322)}
\expandafter\def\csname source@14.positive_det_bound\endcsname{F05:9C (print p.~354; PDF p.~368)}
\expandafter\def\csname source@15.multilinear\endcsname{F04:5.6 (print p.~166; PDF p.~175);\allowbreak\ F05:9B (print p.~346; PDF p.~360);\allowbreak\ F06:1 (print p.~1; PDF p.~8)}
\expandafter\def\csname source@15.tensor_universal\endcsname{F06:1 (print p.~1; PDF p.~8);\allowbreak\ F05:9D (print p.~370; PDF p.~384)}
\expandafter\def\csname source@15.tensor_existence\endcsname{F06:1 (print p.~1; PDF p.~8)}
\expandafter\def\csname source@15.tensor_basis\endcsname{F06:1 (print p.~1; PDF p.~8);\allowbreak\ F05:9D (print p.~370; PDF p.~384)}
\expandafter\def\csname source@15.tensor_maps\endcsname{F06:1 (print p.~1; PDF p.~8);\allowbreak\ F05:9D (print p.~370; PDF p.~384)}
\expandafter\def\csname source@15.tensor_coordinates\endcsname{F06:1 (print p.~1; PDF p.~8);\allowbreak\ F05:9D (print p.~370; PDF p.~384)}
\expandafter\def\csname source@15.associativity_symmetry\endcsname{F06:1 (print p.~1; PDF p.~8);\allowbreak\ F05:9D (print p.~370; PDF p.~384)}
\expandafter\def\csname source@15.tensor_direct_sum\endcsname{F06:1 (print p.~1; PDF p.~8)}
\expandafter\def\csname source@15.tensor_quotient\endcsname{F06:1 (print p.~1; PDF p.~8)}
\expandafter\def\csname source@15.tensor_dual_hom\endcsname{F06:1 (print p.~1; PDF p.~8);\allowbreak\ F06:2 (print p.~41; PDF p.~48)}
\expandafter\def\csname source@15.tensor_algebra\endcsname{F06:3 (print p.~60; PDF p.~67)}
\expandafter\def\csname source@15.permutation_actions\endcsname{F06:4 (print p.~84; PDF p.~91);\allowbreak\ F05:9B (print p.~346; PDF p.~360)}
\expandafter\def\csname source@15.symmetrizers\endcsname{F06:4 (print p.~84; PDF p.~91);\allowbreak\ F06:5 (print p.~96; PDF p.~103);\allowbreak\ F06:9 (print p.~209; PDF p.~216)}
\expandafter\def\csname source@15.exterior\endcsname{F04:5.7 (print p.~173; PDF p.~182);\allowbreak\ F06:5 (print p.~96; PDF p.~103)}
\expandafter\def\csname source@15.exterior_basis\endcsname{F06:5 (print p.~96; PDF p.~103);\allowbreak\ F04:5.7 (print p.~173; PDF p.~182)}
\expandafter\def\csname source@15.wedge_sign\endcsname{F06:5 (print p.~96; PDF p.~103)}
\expandafter\def\csname source@15.determinant_line\endcsname{F06:5 (print p.~96; PDF p.~103);\allowbreak\ F05:9C (print p.~354; PDF p.~368)}
\expandafter\def\csname source@15.symmetric\endcsname{F06:9 (print p.~209; PDF p.~216)}
\expandafter\def\csname source@15.symmetric_basis\endcsname{F06:9 (print p.~209; PDF p.~216)}
\expandafter\def\csname source@15.polynomial_bridge\endcsname{F06:9 (print p.~209; PDF p.~216)}
\expandafter\def\csname source@15.normalization\endcsname{F06:5 (print p.~96; PDF p.~103);\allowbreak\ F06:9 (print p.~209; PDF p.~216)}
\expandafter\def\csname source@15.inner_tensor\endcsname{F05:9D (print p.~370; PDF p.~384);\allowbreak\ F06:2 (print p.~41; PDF p.~48)}
\expandafter\def\csname source@15.tensor_kernel\endcsname{F06:1.19 (print p.~24; PDF p.~31)}
\expandafter\def\csname source@15.tensor_intersections\endcsname{F06:1.15 (print p.~19; PDF p.~26)}
\expandafter\def\csname source@15.tensor_pairing\endcsname{F06:1.22 (print p.~30; PDF p.~37);\allowbreak\ F06:1.23 (print p.~32; PDF p.~39);\allowbreak\ F05:9D (print p.~370; PDF p.~384)}
\expandafter\def\csname source@15.multilinear_dimension\endcsname{F04:5.6 (print p.~166; PDF p.~175);\allowbreak\ F05:9B (print p.~346; PDF p.~360);\allowbreak\ F05:9D (print p.~370; PDF p.~384)}
\expandafter\def\csname source@15.contraction\endcsname{F06:3.11 (print p.~71; PDF p.~78);\allowbreak\ F06:3.13 (print p.~72; PDF p.~79)}
\expandafter\def\csname source@15.tensor_derivation\endcsname{F06:3.6 (print p.~66; PDF p.~73)}
\expandafter\def\csname source@15.exterior_pairing\endcsname{F06:5.6 (print p.~106; PDF p.~113);\allowbreak\ F06:5.9 (print p.~109; PDF p.~116)}
\expandafter\def\csname source@15.exterior_direct_sum\endcsname{F06:5.15 (print p.~120; PDF p.~127);\allowbreak\ F06:5.20 (print p.~125; PDF p.~132)}
\expandafter\def\csname source@15.interior_product\endcsname{F06:5.13 (print p.~117; PDF p.~124);\allowbreak\ F06:5.14 (print p.~118; PDF p.~125);\allowbreak\ F06:5.32 (print p.~143; PDF p.~150)}
\expandafter\def\csname source@15.exterior_derivation\endcsname{F06:5.10 (print p.~110; PDF p.~117);\allowbreak\ F06:5.17 (print p.~123; PDF p.~130)}
\expandafter\def\csname source@15.top_pairing\endcsname{F06:6.8 (print p.~159; PDF p.~166);\allowbreak\ F06:6.9 (print p.~161; PDF p.~168)}
\expandafter\def\csname source@15.charcoefficients\endcsname{F06:7.1 (print p.~174; PDF p.~181);\allowbreak\ F06:7.5 (print p.~182; PDF p.~189)}
\expandafter\def\csname source@15.symmetric_direct_sum\endcsname{F06:9.8 (print p.~217; PDF p.~224)}
\expandafter\def\csname source@15.symmetric_derivation\endcsname{F06:9.5 (print p.~213; PDF p.~220);\allowbreak\ F06:9.6 (print p.~214; PDF p.~221);\allowbreak\ F06:9.16 (print p.~224; PDF p.~231)}
\expandafter\def\csname source@15.hilbert_series\endcsname{F06:2.5 (print p.~44; PDF p.~51);\allowbreak\ F06:9.11 (print p.~219; PDF p.~226)}
\expandafter\def\csname source@15.polynomial_polarization\endcsname{F06:9.12 (print p.~219; PDF p.~226);\allowbreak\ F06:9.15 (print p.~223; PDF p.~230);\allowbreak\ F06:9.18 (print p.~226; PDF p.~233)}
\expandafter\def\csname source@09.real_complex_bridge\endcsname{F04:8.1 (print p.~270; PDF p.~279);\allowbreak\ F06:7.9 (print p.~188; PDF p.~195)}

% END source-metadata.tex


% END preamble.tex

\begin{document}
\frontmatter
\title{Linear Algebra and Foundational Higher Algebra\\Unified Mathematical Notes}
\author{D05 Mathematics Library}
\date{27 October 2026}
\maketitle

% BEGIN frontmatter.tex
\chapter*{Scope, conventions, and foundations}
\addcontentsline{toc}{chapter}{Scope, conventions, and foundations}
These notes cover all fifteen agreed D05 modules, from linear systems to multilinear algebra. The coverage unit is the 249 stable entries of the current manuscript (248 from the original D05 registry and one added by the source-led final review), not every exercise in the three source books. Proofs are newly written in a common language. Source routes allow a reader to compare hypotheses and constructions; a source citation never replaces the proof given here.

The scalars are a named field $\K$. Inner-product arguments use $\R$ or $\C$ and the convention that $\ip{u}{v}$ is linear in $u$ and conjugate-linear in $v$. Coordinates are columns. The algebraic dual map is $T^\vee$, the inner-product adjoint is $T^\dagger$, and the pseudoinverse is $T^+$. In particular, Axler's dagger notation for the pseudoinverse is translated to a plus here. Similarity and congruence are different operations. A Jordan chain is ordered so that its matrix has ones above the diagonal. Invariant factors are ordered by increasing divisibility.

A statement about a finite-dimensional space includes the zero space unless expressly excluded. The determinant of an empty matrix is $1$, an empty sum is $0$, and an empty product is $1$. Characteristic restrictions and splitting hypotheses are stated where used. Algebraic tensor, exterior, and symmetric constructions are separated from averaging formulas requiring factorials to be invertible.

\section*{Meaning of proof closure}
The notes prove the registered D05 conclusions relative to explicitly stated shared foundations. They do not reconstruct set theory, the real numbers, or Lebesgue measure from axioms. The foundational chapter identifies those interfaces and proves the particular algebraic and analytic consequences needed here. An external interface is not an unstated missing D05 lemma. Finite-dimensional proofs use finite basis extraction; the axiom of choice is invoked explicitly only for unrestricted basis extension.

This is a mathematical manuscript, not a Lean-checked library. The verification register distinguishes source organization, mathematical review, structural checks, and compilation. Successful typesetting or the presence of an entry label is not a formal correctness certificate. The modular sources and the flattened input are generated from the same mathematical text.

\section*{Reading order}
Vector-space and linear-map foundations precede elimination. Some later constructions are mentioned earlier and proved at their stable entry; this is a reading reference, not an additional assumption. In particular, determinant is first built by alternating multilinearity, and its top-exterior-power interpretation is proved later. Spectral theory does not supply the algebraic classification of forms or the construction of determinant.

% END frontmatter.tex

\tableofcontents
\mainmatter

% BEGIN foundations.tex
\chapter{Explicit shared foundations and interface lemmas}
\label{ch:foundations}

\section{Logical and scalar foundations}
We use classical sets, functions, finite Cartesian products, quotients by equivalence relations, the natural numbers with induction, and the integers. Finite sums may be regrouped and reindexed by bijections. A field is a commutative ring with $1\ne0$ in which every nonzero element has an inverse. A vector space and a linear map will be defined in the D05 chapters. Finite-dimensional assertions do not require the general axiom of choice. For unrestricted vector spaces, whenever a basis or a complement is chosen, we use Zorn's lemma: a partially ordered set in which every chain has an upper bound has a maximal element. Its use in basis extension is stated explicitly.

An associative unital $\K$-algebra means a $\K$-vector space with a bilinear associative multiplication and an identity element; an algebra homomorphism is a linear map preserving multiplication and the identity. Commutativity is an additional condition and is not imposed on tensor algebras or endomorphism algebras. A module over a ring has the vector-space addition and scalar-action identities, with scalars from that ring; no inverse-scalar axiom is part of its definition. These are definitions of the structures used below, not imported decomposition or classification theorems.

The shared real/complex interface comprises a complete ordered field $\R$, its usual finite-dimensional topology, $\C=\R[i]$, conjugation, absolute value, nonnegative real square roots, and the elementary trigonometric identities and polar representation of a nonzero complex number. We also use the following real-analysis facts: continuity of finite sums and products, the intermediate value theorem, and attainment of extrema by continuous functions on compact subsets of $\R^n$; closed bounded subsets of $\R^n$ are compact. These are D02/D03 foundations, not conclusions asserted to have been proved by this algebra manuscript. The fundamental theorem of algebra itself is proved below from this interface.

For geometric volume only, we use complete Lebesgue measure on $\R^n$, its product-measure/Tonelli property for nonnegative measurable functions, the ordinary one-dimensional length rule under translations and nonzero dilations, and that Borel sets and their null-set completions generate the Lebesgue-measurable sets. The linear change rule needed in D05 is proved below. General measure construction belongs to D04/D11.

\section{Free vectors and finite combinatorics}
For a set $S$, let $\K^{(S)}$ be the finitely supported functions $S\to\K$. Pointwise addition and scalar multiplication preserve finite support. For $s\in S$, let $\delta_s$ be $1$ at $s$ and $0$ elsewhere. Every $a\in\K^{(S)}$ has the unique expression $a=\sum_{s\in\operatorname{supp}a}a(s)\delta_s$: evaluate at each $s$ to prove uniqueness. Thus any function $f:S\to W$ has a unique linear extension $\sum a(s)\delta_s\mapsto\sum a(s)f(s)$ into a vector space $W$. This construction uses no basis choice for $S$.

An $n$-element set has $\binom np$ subsets of size $p$. Indeed each such subset has $p!$ orderings, whereas ordered lists of $p$ distinct elements number $n(n-1)\cdots(n-p+1)$; dividing these positive integers gives the formula. The value is $0$ for $p>n$ and $1$ for $p=0$. For $n\ge1$, nonnegative integer solutions of $a_1+\cdots+a_n=p$ correspond to strings of $p$ stars and $n-1$ separators, hence number $\binom{n+p-1}{p}$. For $n=0$ there is one empty tuple when $p=0$ and none otherwise. These are integer counts, not divisions inside an arbitrary field.

For $\sigma\in S_n$, define $\operatorname{sgn}\sigma=(-1)^{I(\sigma)}$, where $I(\sigma)$ counts inverted pairs. To justify the sign laws, work first in $\Z[x_1,\ldots,x_n]$ with $D=\prod_{i<j}(x_j-x_i)$. Permuting variables gives $D(x_{\sigma(1)},\ldots,x_{\sigma(n)})=\operatorname{sgn}(\sigma)D(x_1,\ldots,x_n)$, since each unordered pair occurs once and is reversed precisely for an inversion. Applying two permutations and cancelling the nonzero polynomial $D$ proves multiplicativity. Swapping two variables changes the sign: the factor for that pair changes sign and all other paired factors merely exchange. Thus a transposition has sign $-1$. These integer identities remain true after mapping $1,-1$ into any field, including characteristic two.

\section{Polynomial division, principal ideals, and factorization}
For a commutative ring $R$, a formal polynomial in $m$ variables is a finitely supported coefficient family $(a_\alpha)_{\alpha\in\N^m}$. Addition is coefficientwise and multiplication is convolution:
$(ab)_\gamma=\sum_{\alpha+\beta=\gamma}a_\alpha b_\beta$.
Only finitely many terms are nonzero. Regrouping the finite sums proves associativity and distributivity; swapping the two indices proves commutativity. The coefficient family supported at $0$ with value $1$ is the unit. Thus the notation $R[x_1,\ldots,x_m]$ denotes an explicit ring, not an identification with polynomial functions. In one variable over a field, the highest nonzero coefficient of a product is the product of the highest coefficients, so nonzero degrees add. We use the degree only for a nonzero polynomial, or explicitly separate the zero case. The finite free module $R^n$ consists of coordinate columns with componentwise operations; statements about determinants over $R[x]$ use only these operations and finite sums, not a classification theorem for arbitrary modules.

\begin{lemma}[Division algorithm]\label{aux:polynomial-division}
For $a,b\in\K[x]$ with $b\ne0$, there exist unique $q,r$ with $a=qb+r$ and $r=0$ or $\deg r<\deg b$.
\end{lemma}
\begin{proof}
If $a=0$, take $q=r=0$; if $a\ne0$ and $\deg a<\deg b$, take $q=0,r=a$. Otherwise subtract the unique monomial multiple of $b$ cancelling the leading term of $a$; its coefficient exists because the leading coefficient of $b$ is invertible. The new remainder is zero or has strictly smaller degree. Repeating terminates since nonnegative integers cannot decrease indefinitely. Summing the subtracted monomials gives $q$. If two expressions existed, $(q-q')b=r'-r$. A nonzero left side has degree at least $\deg b$, whereas a nonzero right side has smaller degree; hence both are zero.
\end{proof}

Every nonzero ideal $I\subseteq\K[x]$ has a monic element $d$ of least degree. Division of $f\in I$ by $d$ leaves a remainder in $I$ that is zero or has smaller degree, necessarily zero. Therefore $I=(d)$. Applied to $(a,b)$ with at least one of $a,b$ nonzero, this proves that a monic greatest common divisor $d$ exists, divides $a,b$, and satisfies $d=ua+vb$ for some polynomials $u,v$. Conversely every common divisor divides this expression. For $a=b=0$, set $\gcd(0,0)=0$; it is not a monic polynomial and the generated ideal is the zero ideal. For a nonzero pair the Euclidean algorithm computes the monic gcd: repeatedly divide, and express each remainder as a combination of the preceding two; backward substitution expresses the last nonzero remainder as a combination of $a,b$, followed by normalization of its leading coefficient.

\begin{lemma}[Irreducibles, unique factorization, and Chinese remainders]\label{aux:polynomial-crt}
Every nonzero polynomial is a unit times a product of monic irreducibles, uniquely up to their order. If $q_1,\ldots,q_s$ are pairwise relatively prime monic polynomials, then
\[
 \K[x]/(q_1\cdots q_s)\longrightarrow\bigoplus_i\K[x]/(q_i),\qquad [f]\longmapsto([f]_i)_i
\]
is an algebra isomorphism.
\end{lemma}
\begin{proof}
Factor a nonconstant reducible polynomial into smaller positive degrees and induct on the degree to obtain existence. If an irreducible $p$ does not divide $a$, its gcd with $a$ is $1$, so $up+va=1$. If $p\mid ab$, multiplying by $b$ gives $p\mid b$. Thus irreducibles are prime. Cancelling one prime factor at a time proves uniqueness; a degree-zero nonzero polynomial is a unit.

Put $Q=\prod_iq_i$ and $Q_i=Q/q_i$. The gcd of $Q_i,q_i$ is $1$, so choose $u_iQ_i+v_iq_i=1$. The polynomial $e_i=u_iQ_i$ is $1$ modulo $q_i$ and $0$ modulo every other $q_j$. Thus $\sum_i f_i e_i$ has any prescribed residue tuple, proving surjectivity. If all residues of $f$ vanish, all $q_i$ divide $f$; successive primality of their irreducible factors implies $Q\mid f$. This is exactly the kernel needed for injectivity. For an empty family, $Q=1$ and both sides are zero algebras in the corresponding quotient convention.
\end{proof}

The quotient $\K[x]/(q)$ for monic $q$ of degree $d$ has basis $1,x,\ldots,x^{d-1}$, since division gives a unique representative of degree below $d$. An element $[f]$ is invertible precisely when $\gcd(f,q)=1$: Bezout gives its inverse, and an inverse identity modulo $q$ gives a Bezout identity. In particular the quotient by an irreducible polynomial is a field.

For later nilpotent calculations, formal differentiation satisfies $(fg)'=f'g+fg'$ by expanding monomials. The binomial expansion gives
\[
 f(x+h)=f(x)+hf'(x)+h^2R(x,h)
\]
for a polynomial $R$ over any field. In characteristic zero it also gives the finite Taylor formula $f(x+h)=\sum_{j=0}^{\deg f} f^{(j)}(x)h^j/j!$. These are polynomial identities, so commuting elements of any algebra may be substituted. No analytic convergence is involved.

\section{Roots over the real and complex fields}
\begin{lemma}[Complex roots]\label{aux:complex-roots}
Every complex number has an $m$th root for each positive integer $m$.
\end{lemma}
\begin{proof}
Zero has root zero. If $z=r(\cos\theta+i\sin\theta)$ with $r>0$, the intermediate value theorem applied to $t\mapsto t^m$ gives a positive real $r^{1/m}$. The multiplication formula for sine and cosine, iterated $m$ times, shows that $r^{1/m}(\cos(\theta/m)+i\sin(\theta/m))$ has $m$th power $z$.
\end{proof}

\begin{theorem}[Fundamental theorem of algebra]\label{aux:fta}
Every nonconstant polynomial in $\C[x]$ has a root and factors into linear factors.
\end{theorem}
\begin{proof}
For $p(z)=a_nz^n+\cdots+a_0$, $a_n\ne0$, the inequality
\[
 |p(z)|\ge |a_n||z|^n-\sum_{j<n}|a_j||z|^j
\]
shows $|p(z)|\to\infty$ as $|z|\to\infty$: for $|z|\ge1$ the lower-degree sum is at most $C|z|^{n-1}$. Choose a disk outside which $|p(z)|>|p(0)|$. Continuity and compactness give a global minimizer $z_0$ of $|p|$ inside that disk.

Suppose $p(z_0)\ne0$. Expand $p(z_0+w)/p(z_0)=1+aw^m+\sum_{j>m}b_jw^j$ with $a\ne0$ and $m\ge1$ the first nonconstant exponent. Choose $u$ with $au^m=-1$ by the preceding lemma. For real $0<t\le1$,
\[
 \left|\frac{p(z_0+tu)}{p(z_0)}\right|
 \le 1-t^m+Ct^{m+1},\qquad C=\sum_{j>m}|b_ju^j|.
\]
Choose $t<1$ and, if $C>0$, also $Ct<1$. The right side is strictly below $1$, contradicting minimality. Hence $p(z_0)=0$. Division by $x-z_0$ has zero remainder by evaluation at $z_0$. Induction on degree factors the quotient and proves the factorization assertion.
\end{proof}

For a real polynomial, conjugating its coefficients shows that nonreal roots occur with their conjugates and with equal multiplicity. Pairing the factors gives real quadratics $(x-z)(x-\bar z)=x^2-2\Re(z)x+|z|^2$ with negative discriminant. The remaining factors are real linear factors. Every odd-degree real polynomial has a real root: its leading term has opposite signs at sufficiently large positive and negative real inputs, and the intermediate value theorem applies. These facts distinguish real triangularizability from complex triangularizability.

\section{The exact volume interface}
\begin{theorem}[Linear change of volume]\label{aux:volume-change}
For an invertible real $n$-by-$n$ matrix $A$ and Lebesgue-measurable $E\subseteq\R^n$,
\[
 \mu_n(AE)=|\det A|\,\mu_n(E).
\]
For singular $A$, $AE$ is contained in a null linear subspace and has measure zero. The zero-dimensional measure gives mass $1$ to its one point.
\end{theorem}
\begin{proof}
First take Borel $E$ and an elementary coordinate transformation. A coordinate permutation preserves product measure by reordering the nonnegative iterated integral of $1_E$. A dilation of one coordinate by $c\ne0$ multiplies its inner one-dimensional integral by $|c|$, hence multiplies the whole integral by $|c|$. A shear $x_i\mapsto x_i+cx_j$ with $j\ne i$ preserves the integral: when all coordinates except $x_i$ are fixed, this is a translation by a constant, and one-dimensional translation invariance applies. Tonelli permits this order of integration even when the measure is infinite. Elementary transformations are homeomorphisms, so the images of Borel sets are Borel.

Elimination (\kref{01.elimination}) expresses an invertible $A$ as a product of such elementary matrices: reducing $A$ to $I$ gives $E_r\cdots E_1A=I$, hence $A=E_1^{-1}\cdots E_r^{-1}$. The measure multipliers of the factors are their absolute determinants by the determinant row rules. Multiplicativity of determinant (\kref{05.det_product}) proves the formula for Borel $E$. If $E$ is Lebesgue measurable, write it as a Borel set modified by a subset of a Borel null set. The Borel formula sends that containing null set to a null set; completeness therefore proves measurability of $AE$ and the same formula.

If $A$ is singular, extend a basis of its proper image to a basis of $\R^n$. An invertible linear map then sends that image into $\R^{n-1}\times\{0\}$. This coordinate hyperplane has measure zero: intersect with $[-m,m]^n$, use product measure and the zero length of $\{0\}$, then take the countable union over $m$. The already proved invertible rule shows the original proper subspace is null. Every subset of it is measurable and null by completeness, including $AE$. We state this separately to avoid an undefined expression $0\cdot\infty$.
\end{proof}

% END foundations.tex


% BEGIN modules/M05.02.tex
\chapter{Vector spaces, bases, and dimension}
Throughout this chapter $\K$ is a field. A finite family is indexed, so two occurrences of the same vector are not silently identified.

\Kentry{02.vector_space}{Vector spaces and the scalar field}
\begin{definition}
A $\K$-vector space is an abelian group $(V,+,0)$ with a map $\K\times V\to V$, $(a,v)\mapsto av$, satisfying
$a(u+v)=au+av$, $(a+b)v=av+bv$, $(ab)v=a(bv)$ and $1v=v$.
\end{definition}
These axioms imply $0v=0$ by cancelling $0v$ in $0v=(0+0)v$; similarly $a0=0$. Also $(-a)v=-(av)$ because their sum is $0v$. If $av=0$ and $a\ne0$, multiplication by $a^{-1}$ gives $v=0$. A complex space can be viewed as a real space by restricting scalars, but these are different structures and can have different dimensions.

\Kentry{02.subspace}{Subspaces}
\begin{proposition}
A subset $U\subseteq V$ is a vector subspace with the inherited operations if and only if it contains $0$ and is closed under addition and scalar multiplication.
\end{proposition}
\begin{proof}
The conditions are necessary because a vector space has these operations. Conversely they make the inherited operations maps into $U$. Additive inverses belong to $U$ since $-u=(-1)u$. All axioms are equalities inherited from $V$, including the group laws. Thus $U$, with these operations, is a vector space. The inclusion $U\to V$ is linear; this structured inclusion must be distinguished from an arbitrary choice of coordinates on $U$.
\end{proof}

\Kentry{02.span}{Finite linear combinations and span}
For a subset $S\subseteq V$, define
\[
 \Span(S)=\left\{\sum_{s\in F}a_s s:F\subseteq S\text{ finite},\ a_s\in\K\right\}.
\]
The empty sum is $0$. Adding two such sums uses the finite union of their supports, and scalar multiplication multiplies their coefficients. Thus the span is a subspace. This definition uses finite support even when $S$ is infinite.

\Kentry{02.span_properties}{The universal property of span}
\begin{proposition}
For every subspace $U$, $\Span(S)\subseteq U$ if and only if $S\subseteq U$. Consequently span is monotone and idempotent, and $\Span(\varnothing)=\{0\}$.
\end{proposition}
\begin{proof}
Each generator is a one-term sum, so $S\subseteq\Span(S)$. A subspace containing $S$ contains each finite linear combination by repeated closure under addition and scalar multiplication. This proves the equivalence. If $S\subseteq T$, apply it to the subspace $\Span(T)$. Applying it to $S'=\Span(S)$ proves idempotence. With no generators the only possible sum is the empty sum.
\end{proof}

\Kentry{02.independence}{Linear independence}
An indexed family $(v_i)_{i\in I}$ is linearly independent if $\sum_i a_i v_i=0$ with finitely supported $(a_i)$ implies $a_i=0$ for every $i$. The empty family is independent. If $v_j=0$, the relation with coefficient $1$ at $j$ disproves independence. If $v_i=v_j$ for distinct indices, the coefficients $1,-1$ give a nontrivial relation, also in characteristic two. Every subfamily of an independent family is independent, by extending its coefficients by zero. A finite list is dependent exactly when one member is a combination of the others: solve a nontrivial relation for a member with nonzero coefficient; conversely move that combination to one side.

\Kentry{02.exchange}{Exchange and the length inequality}
\begin{theorem}
If $v_1,\ldots,v_r$ are independent and $w_1,\ldots,w_s$ span $V$, then $r\le s$. One can replace $r$ members of the spanning list by $v_1,\ldots,v_r$ and still have a spanning list.
\end{theorem}
\begin{proof}
Suppose $v_1,\ldots,v_{j-1}$ have already replaced distinct members of the spanning list. Express $v_j$ in the current list. Some coefficient of an unreplaced $w$ must be nonzero: otherwise $v_j$ belongs to the span of its independent predecessors. Solve the expression for that $w$. The resulting list with $v_j$ in its place spans every member of the preceding list and hence all of $V$. This constructs the replacement at stage $j$. If $j>s$, there are no unreplaced members, which is impossible. Starting at $j=1$ therefore permits precisely the claimed replacements and yields $r\le s$.
\end{proof}

\Kentry{02.basis}{Bases and unique expansions}
\begin{proposition}
An indexed family is a basis if it is independent and spans $V$. Equivalently, each vector has exactly one finitely supported coordinate family in it.
\end{proposition}
\begin{proof}
Spanning gives existence. Subtracting two expansions produces a finite relation, so independence gives uniqueness. Conversely, existence implies spanning. A relation represents $0$; uniqueness compared with its all-zero coordinates implies independence.
\end{proof}

\Kentry{02.basis_extension}{Extraction and extension of bases}
\begin{theorem}
A finite spanning list contains a basis. Any independent list in a finitely generated space extends to a finite basis.
\end{theorem}
\begin{proof}
If a finite spanning list is dependent, solve a relation for one member with nonzero coefficient and discard that member. Spanning is preserved and the length decreases; termination leaves an independent spanning list. For extension, begin with the given independent list and process a finite spanning list of $V$. Append its next member if it is not in the span of the current list. Appending preserves independence: a relation with nonzero coefficient on the new vector would put it in the old span. At the end every spanning-list member lies in the span, so the resulting independent list is a basis. The length inequality shows that an independent family in a finitely generated space is finite: a family with more than $s$ indices would contain an independent subfamily of length $s+1$.
\end{proof}
For completeness, the genuinely infinite-dimensional variant uses the explicitly declared choice principle Zorn's lemma. Partially order independent families of vectors, viewed as subsets of $V$, containing a prescribed independent subset, by inclusion. The union of a chain is independent because every finite relation lies in one member of the chain. A maximal member therefore exists. If it did not span, one could append a vector outside its span and preserve independence, contradicting maximality. This proves basis extension under that choice principle; finite-dimensional arguments above do not use it.

\Kentry{02.dimension}{Finite dimension}
A space is finite-dimensional when it has a finite spanning family. By extraction it then has a finite basis. Two finite bases have equal lengths: apply the length inequality in each direction. Define $\dim_\K V$ to be this common length. The zero space has the empty basis and dimension zero. Conversely a basis of length zero spans only $0$, so dimension zero characterizes the zero space.

\Kentry{02.subspace_dimension}{Dimensions of subspaces}
\begin{theorem}
If $U\subseteq V$ and $\dim V=n<\infty$, then $U$ has a finite basis, $\dim U\le n$, and equality holds precisely when $U=V$.
\end{theorem}
\begin{proof}
Choose independent vectors in $U$ successively while their span is smaller than $U$. More than $n$ choices would violate exchange with a basis of $V$, so this process stops with a basis of $U$. Extend it to a basis of $V$. The extension adds $n-\dim U$ vectors. If dimensions agree, it adds none, so $U$ already spans $V$. The reverse implication is immediate from the definition of dimension.
\end{proof}

\Kentry{02.sum_intersection}{Sums and intersections}
For subspaces $U,W\subseteq V$, their intersection contains zero and is closed under both operations. The set $U+W=\{u+w:u\in U,w\in W\}$ has the same closure properties. It contains $U\cup W$, and every subspace containing that union contains each $u+w$; hence $U+W=\Span(U\cup W)$. Arbitrary intersections of subspaces are subspaces by the same closure test. A union need not be: in $\K^2$ the two coordinate lines contain $e_1,e_2$ but their union does not contain $e_1+e_2$.

\Kentry{02.grassmann}{The dimension formula for a sum}
\begin{theorem}
For finite-dimensional subspaces $U,W$,
\[\dim(U+W)+\dim(U\cap W)=\dim U+\dim W.\]
\end{theorem}
\begin{proof}
Let $c_1,\ldots,c_t$ be a basis of $U\cap W$. Extend it to $c_1,\ldots,c_t,u_1,\ldots,u_a$ in $U$ and to $c_1,\ldots,c_t,w_1,\ldots,w_b$ in $W$. Their combined list $c,u,w$ spans $U+W$. In a zero relation, the sum of the $w$ terms lies in $U\cap W$ and thus in the span of $c$. Independence of the basis $c,w$ forces all $w$ coefficients to vanish. Independence of $c,u$ then forces all remaining coefficients to vanish. Thus the combined list is a basis of length $t+a+b$, and the asserted identity is $(t+a+b)+t=(t+a)+(t+b)$.
\end{proof}

\Kentry{02.examples}{Standard coordinate spaces}
In $\K^n$, the standard vectors $e_j$ form a basis, since $x=\sum_j x_j e_j$ and the $j$th coordinate recovers its coefficient. In $m$-by-$n$ matrices, the matrices $E_{ij}$ with a single $1$ in position $(i,j)$ form a basis by the same entrywise argument. Thus the dimensions are $n$ and $mn$, including empty shapes. For $d\ge0$, formal polynomials of degree at most $d$ have basis $1,x,\ldots,x^d$, because equality of formal polynomials is equality of coefficients. The zero polynomial is included. All monomials $1,x,x^2,\ldots$ are independent, so $\K[x]$ is not finite-dimensional. The standard finitely supported sequences are independent in the space of all sequences $\K^{\N}$, so that space is also infinite-dimensional. Over finite fields, formal polynomials must not be identified with their value functions.

% END modules/M05.02.tex


% BEGIN modules/M05.03.tex
\chapter{Linear maps and structural descriptions}
All spaces and maps in this chapter use a common scalar field $\K$.

\Kentry{03.linear_map}{Linear maps}
A map $T:V\to W$ is linear if $T(u+v)=T(u)+T(v)$ and $T(av)=aT(v)$. It then preserves every finite linear combination, by iterating additivity and homogeneity. In particular $T(0)=0$ and $T(-v)=-T(v)$. A map involving complex conjugation is not complex-linear unless this can be checked separately; changing the scalar field changes the assertion.

\Kentry{03.map_algebra}{Spaces of maps and composition}
Define $(S+T)(v)=S(v)+T(v)$ and $(aT)(v)=aT(v)$. Substitution in the linearity identities shows that these maps are linear. The vector-space axioms hold pointwise, so $\Hom_\K(V,W)$ is a vector space. Composition of linear maps is linear and associative because function composition is associative. For composable maps, $R(S+T)=RS+RT$ and $(R+S)T=RT+ST$, with the analogous scalar identities, by evaluation on a vector. Thus $\End(V)$ is a unital algebra with identity $\id_V$. It need not be commutative: on $\K^2$, let $S(e_1)=e_2,S(e_2)=0$ and $T(e_1)=e_1,T(e_2)=0$. Then $ST(e_1)=e_2$ while $TS(e_1)=0$.

\Kentry{03.basis_extension}{Specifying a linear map on a basis}
\begin{theorem}
For a basis $(v_i)_{i\in I}$ of $V$ and arbitrary vectors $(w_i)_{i\in I}$ in $W$, there is a unique linear $T:V\to W$ such that $T(v_i)=w_i$.
\end{theorem}
\begin{proof}
Set $T(\sum_i a_i v_i)=\sum_i a_i w_i$, using finite-support expansions. Uniqueness of basis coordinates makes this independent of all choices. Coordinates of a sum or scalar multiple are the sums or scalar multiples of coordinates; hence the formula defines a linear map. Every linear map with the prescribed values must satisfy this formula, proving uniqueness.
\end{proof}

\Kentry{03.kernel}{Kernels}
Define $\ker T=\{v\in V:T(v)=0\}$. The zero vector belongs to it, and for $u,v$ in it, $T(au+bv)=aT(u)+bT(v)=0$. Thus the kernel is a subspace of the domain $V$.

\Kentry{03.image}{Images and rank}
Define $\im T=\{T(v):v\in V\}\subseteq W$. It contains zero and $aT(u)+bT(v)=T(au+bv)$, so it is a subspace. If $v_i$ span $V$, then $T(v_i)$ span the image by linearity. When the image is finite-dimensional, define $\rank T=\dim(\im T)$; this does not assert that $T$ is onto all of $W$.

\Kentry{03.injective_surjective}{Injectivity and surjectivity}
\begin{proposition}
$T$ is injective if and only if $\ker T=\{0\}$; it is surjective if and only if $\im T=W$.
\end{proposition}
\begin{proof}
Injectivity forces the only preimage of $T(0)=0$ to be $0$. Conversely, if $T(u)=T(v)$, then $u-v\in\ker T$, and a zero kernel gives $u=v$. The second assertion is the definition of surjectivity. No dimension assumptions are used.
\end{proof}

\Kentry{03.rank_nullity}{Rank and nullity}
\begin{theorem}
If $\dim V=n<\infty$ and $T:V\to W$ is linear, then
\[\dim V=\dim\ker T+\dim\im T.\]
The space $W$ itself need not be finite-dimensional.
\end{theorem}
\begin{proof}
Choose a basis $u_1,\ldots,u_k$ of the kernel and extend it by $v_1,\ldots,v_r$ to a basis of $V$. The vectors $T(v_j)$ span the image, since the $u_i$ map to zero. If $\sum_j a_jT(v_j)=0$, then $\sum_j a_jv_j$ belongs to the kernel and is a combination of the $u_i$. Independence of the combined basis forces all $a_j=0$. Thus $T(v_1),\ldots,T(v_r)$ are an image basis. Since $n=k+r$, the identity follows.
\end{proof}

\Kentry{03.isomorphism}{Linear isomorphisms}
A linear bijection is called an isomorphism. Its inverse is linear: if $y=T(u)$ and $z=T(v)$, then $T(au+bv)=ay+bz$, and uniqueness of preimages gives $T^{-1}(ay+bz)=au+bv$. Identities are isomorphisms, composites of isomorphisms are isomorphisms, and their inverses obey $(ST)^{-1}=T^{-1}S^{-1}$ by direct composition. An isomorphism carries a basis to a basis: spanning follows from surjectivity, and a relation in the images pulls back to a relation by injectivity.

\Kentry{03.dimension_classification}{Classification by dimension}
\begin{theorem}
Finite-dimensional spaces over the same field are isomorphic exactly when their dimensions agree.
\end{theorem}
\begin{proof}
An isomorphism transports a basis, so dimensions agree. Conversely, choose bases $v_1,\ldots,v_n$ and $w_1,\ldots,w_n$ and use the basis-extension theorem to define maps taking each basis to the other. Their composites fix every basis vector and hence every vector. They are inverse isomorphisms. Different choices of bases can produce different isomorphisms; existence does not make this correspondence canonical.
\end{proof}

\Kentry{03.equal_dimension}{Equal-dimensional injectivity and surjectivity}
\begin{theorem}
For $T:V\to W$ with $\dim V=\dim W=n<\infty$, injectivity, surjectivity and invertibility are equivalent.
\end{theorem}
\begin{proof}
Rank-nullity says $\rank T=n$ exactly when its kernel has dimension zero, equivalently when it is injective. Since $\im T\subseteq W$, dimension $n$ is also equivalent to $\im T=W$. A linear bijection has a linear inverse. This includes $n=0$, when the unique map between the two zero-dimensional spaces is their unique linear isomorphism.
\end{proof}

\Kentry{03.composition_rank}{Rank under composition}
\begin{proposition}
For $S:U\to V$ and $T:V\to W$ with finite-dimensional images where the ranks below are used,
\[\rank(TS)\le\min(\rank S,\rank T).\]
If $T$ is injective, $\rank(TS)=\rank S$. If $S$ is surjective, $\rank(TS)=\rank T$.
\end{proposition}
\begin{proof}
The image of $TS$ equals $T(\im S)$ and is contained in $\im T$. This gives the second bound by subspace dimension. Restricting $T$ to the finite-dimensional space $\im S$ and applying rank-nullity gives the first bound. When that restriction is injective its kernel is zero, so equality holds. If $S$ is onto, $T(\im S)=T(V)=\im T$, giving the last claim. In particular multiplication on either side by an isomorphism preserves rank.
\end{proof}

\Kentry{03.first_iso}{The first isomorphism theorem}
\begin{theorem}
The map $\widetilde T:V/\ker T\to\im T$, $v+\ker T\mapsto T(v)$, is a linear isomorphism, without a finite-dimensional hypothesis.
\end{theorem}
\begin{proof}
Cosets and their operations are specified in \kref{06.quotient}. If $v-w\in\ker T$, then $T(v)=T(w)$, so the formula is well-defined. It is linear by the representative formulas for sum and scalar multiplication. Every image vector has a preimage coset. If $\widetilde T(v+\ker T)=0$, then $v\in\ker T$, so this is the zero coset. The map is therefore bijective and its inverse is linear. These arguments also directly verify the representative-independence needed here; they do not rely on a dimension count.
\end{proof}

\Kentry{03.hom_dimension}{The dimension of a space of linear maps}
\begin{theorem}
If $v_1,\ldots,v_n$ and $w_1,\ldots,w_m$ are bases of $V,W$, define $E_{ij}(v_k)=\delta_{jk}w_i$. Then the $mn$ maps $E_{ij}$ form a basis of $\Hom(V,W)$.
\end{theorem}
\begin{proof}
Existence and uniqueness of each $E_{ij}$ follow from specifying values on a basis. Write $T(v_j)=\sum_i a_{ij}w_i$. The map $\sum_{i,j}a_{ij}E_{ij}$ has the same value as $T$ on every basis vector and hence equals $T$. If $\sum_{i,j}a_{ij}E_{ij}=0$, evaluate at $v_j$: independence of the $w_i$ forces $a_{ij}=0$ for every $i,j$. This proves both spanning and independence. If $m=0$ or $n=0$, the map space is the zero space and the empty family is its basis.
\end{proof}

% END modules/M05.03.tex


% BEGIN modules/M05.04.tex
\chapter{Matrices and changes of representation}

\Kentry{04.coordinates}{Coordinate isomorphisms}
For an ordered basis $B=(b_1,\ldots,b_n)$, let $c_B:V\to\K^n$ send $\sum_j x_jb_j$ to the column $(x_j)$. Unique expansion makes this well-defined. Coordinates add and scale componentwise, so it is linear. Its inverse sends $(x_j)$ to $\sum_jx_jb_j$, which is linear and inverse in both directions. Reordering $B$ permutes coordinates, so the order is part of the data.

\Kentry{04.matrix_of_map}{The matrix of a linear map}
For $T:V\to W$, bases $B=(b_j)_{j=1}^n$ and $C=(c_i)_{i=1}^m$, define $A=[T]_{C\leftarrow B}$ by $T(b_j)=\sum_i A_{ij}c_i$. If $v=\sum_jx_jb_j$, linearity gives
\[T(v)=\sum_i\left(\sum_jA_{ij}x_j\right)c_i.\]
Thus $c_C(Tv)=A c_B(v)$, and a matrix is uniquely determined by this identity. Conversely every $m$-by-$n$ matrix defines a linear map $\K^n\to\K^m$ by this formula and, by the coordinate isomorphisms, a unique map $V\to W$. Rows index the codomain basis and columns index the domain basis.

\Kentry{04.matrix_product}{Matrix arithmetic}
Addition and scalar multiplication are entrywise. For an $\ell$-by-$m$ matrix $A$ and an $m$-by-$n$ matrix $B$, set $(AB)_{ij}=\sum_{k=1}^m A_{ik}B_{kj}$. Associativity follows by expanding both sides:
\[((AB)C)_{ij}=\sum_{k,h}A_{ih}B_{hk}C_{kj}=(A(BC))_{ij}.\]
Interchanging finite sums is valid in the scalar ring. Distributivity follows by distributing the summands, and the Kronecker-delta formula for $I$ gives $AI=A$ and $IA=A$. Empty inner sums are zero. These rules also hold over any commutative ring, and do not describe entrywise multiplication.

\Kentry{04.composition}{Composition and identity matrices}
\begin{proposition}
If $T:U\to V$ and $S:V\to W$ use bases $B,C,D$ respectively, then
\[[ST]_{D\leftarrow B}=[S]_{D\leftarrow C}[T]_{C\leftarrow B}.\]
\end{proposition}
\begin{proof}
Write $T(b_j)=\sum_k B_{kj}c_k$ and $S(c_k)=\sum_i A_{ik}d_i$. Substitution yields $ST(b_j)=\sum_i(\sum_kA_{ik}B_{kj})d_i$. Uniqueness of coordinates proves the formula. The identity map sends each basis vector to itself, so its same-basis matrix is $I$. If its input and output bases differ, its matrix is a change-of-coordinate matrix, not necessarily $I$.
\end{proof}

\Kentry{04.inverse_matrix}{Inverse matrices and inverse maps}
\begin{proposition}
A square matrix $A$ has a two-sided inverse precisely when $x\mapsto Ax$ is bijective. Its inverse matrix represents the inverse linear map.
\end{proposition}
\begin{proof}
If $AB=BA=I$, the represented maps are mutual inverses. Conversely a bijective represented map has a linear inverse; let its matrix be $B$. The composition formula gives both identities. Inverse matrices are unique, because if $B,C$ are inverses then $B=B(AC)=(BA)C=C$. In finite dimension a one-sided inverse is also two-sided: $BA=I$ makes $A$ injective, hence bijective, and its inverse must be $B$; the other case is analogous by surjectivity.
\end{proof}

\Kentry{04.change_basis}{Independent changes of domain and codomain basis}
Let $P=c_Bc_{B'}^{-1}$ and $Q=c_Cc_{C'}^{-1}$ convert new coordinates to old coordinates. They are invertible. If $A=[T]_{C\leftarrow B}$ and $A'=[T]_{C'\leftarrow B'}$, then for every $x$,
\[QA'x=APx.\]
It follows that $A'=Q^{-1}AP$. Conversely this formula is forced by the two coordinate conversions, so it does not identify an arbitrary pair of basis changes with a single conjugation.

\Kentry{04.similarity}{Similarity}
When $T$ is an endomorphism and the same basis is used on both sides, the preceding formula becomes $A'=P^{-1}AP$. This defines similarity. It is an equivalence relation: use $I$, $P^{-1}$ and the product of the two change matrices for reflexivity, symmetry and transitivity. For every polynomial $p$, cancellation of adjacent $PP^{-1}$ gives $p(A')=P^{-1}p(A)P$, first for powers and then by linearity. Hence $p(A')=0$ exactly when $p(A)=0$, so the monic minimal polynomials defined in \kref{08.minpoly} agree. Also
\[tI-A'=P^{-1}(tI-A)P.\]
Determinant multiplicativity over $\K[t]$ from \kref{05.det_product} makes their characteristic polynomials equal. This is a forward use of the independently proved determinant identity, not a proof of that identity from similarity invariance.

\Kentry{04.equivalence}{Rectangular rank normal form}
\begin{theorem}
Every map $T:V\to W$ of rank $r$, with finite-dimensional domain and codomain, has bases in which its matrix has $I_r$ in the upper-left block and zeros elsewhere. Two matrices of the same shape are equivalent under independent invertible row and column changes exactly when their ranks agree.
\end{theorem}
\begin{proof}
Take a basis $u_1,\ldots,u_k$ of $\ker T$ and extend by $v_1,\ldots,v_r$ to a domain basis. Rank-nullity proves that $T(v_1),\ldots,T(v_r)$ form an image basis. Extend this image basis to one of $W$, and order the domain basis as $v_1,\ldots,v_r,u_1,\ldots,u_k$. The matrix is the asserted block matrix. Basis changes preserve rank because the coordinate maps are isomorphisms. Conversely matrices of equal rank are each equivalent to this same block matrix, and composing the equivalences gives an equivalence between them. The argument includes $r=0$.
\end{proof}

\Kentry{04.blocks}{Block matrices and invariant subspaces}
Relative to finite direct-sum decompositions of domain and codomain, the block $A_{ij}$ is the map from the $j$th input summand to the $i$th output summand obtained by inclusion, the map, and coordinate projection. Applying two block maps to a tuple gives $(AB)_{ij}=\sum_kA_{ik}B_{kj}$. This proves the block multiplication law without assuming equal block sizes. For an endomorphism with $V=U\oplus W$, invariance of $U$ is exactly the vanishing of the lower-left block, since vectors of $U$ have coordinates $(u,0)$. If both summands are invariant, both off-diagonal blocks vanish. The same component argument gives block diagonality for any finite invariant direct sum.

\Kentry{04.rank_bridge}{Row rank, column rank, and map rank}
\begin{theorem}
The image of $x\mapsto Ax$ is the column span of $A$. Column rank and row rank agree; transposition preserves rank.
\end{theorem}
\begin{proof}
The identity $Ax=\sum_jx_jA_{\cdot j}$ proves the image assertion in both directions. Let $s$ be the dimension of the row space, choose its basis as the rows of an $s$-by-$n$ matrix $R$, and express each row of $A$ as a combination of these rows. The coefficients form a matrix $C$ with $A=CR$. Since a map with intermediate space $\K^s$ has rank at most $s$, column rank of $A$ is at most its row rank. Apply the same reasoning to $A^t$. The column rank of $A^t$ is the row rank of $A$, and its row rank is the column rank of $A$, giving the reverse inequality. Thus the two ranks are equal, also for empty matrices.
\end{proof}

\Kentry{04.trace}{Trace and its invariance}
For a square matrix define $\tr A=\sum_i A_{ii}$. For $A$ of shape $m\times n$ and $B$ of shape $n\times m$,
\[\tr(AB)=\sum_{i=1}^m\sum_{j=1}^n A_{ij}B_{ji}
 =\sum_{j=1}^n\sum_{i=1}^m B_{ji}A_{ij}=\tr(BA).\]
Therefore $\tr(P^{-1}AP)=\tr(APP^{-1})=\tr A$. Define $\tr T$ by any basis matrix; the equality proves independence of the chosen basis. Entrywise addition shows $\tr(aA+bB)=a\tr A+b\tr B$. The trace of the zero-dimensional operator is the empty sum $0$.

\begin{theorem}[Trace functionals and commutators]
Let $n\ge1$ and let $\K$ be any field. A linear functional $f:\K^{n\times n}\to\K$ satisfies $f(AB)=f(BA)$ for all $A,B$ exactly when $f=c\tr$ for a unique scalar $c$. Moreover,
\[
\Span\{AB-BA:A,B\in\K^{n\times n}\}=\ker\tr.
\]
If additionally $f(I_n)=n$ and the element $n\cdot1_\K$ is nonzero, then $f=\tr$.
\end{theorem}
\begin{proof}
Write $E_{ij}$ for the matrix units, so $E_{ij}E_{k\ell}=\delta_{jk}E_{i\ell}$. Trace vanishes on every commutator by the cyclic trace identity, proving that their span lies in $\ker\tr$. For $i\ne j$,
\[
E_{ij}=[E_{ii},E_{ij}],\qquad E_{ii}-E_{jj}=[E_{ij},E_{ji}],
\]
where $[A,B]=AB-BA$. Every traceless matrix $A=(a_{ij})$ is a sum of its off-diagonal terms $a_{ij}E_{ij}$ and its diagonal part. The latter equals
$\sum_{i=1}^{n-1}a_{ii}(E_{ii}-E_{nn})$, because $a_{nn}=-\sum_{i<n}a_{ii}$. Thus it is a linear combination of commutators. When $n=1$, the kernel is zero and both sums are empty, so the same assertion holds.

The functional condition is exactly that $f$ vanish on all commutators, hence on $\ker\tr$. For every $A$, the matrix $A-(\tr A)E_{11}$ is traceless. It follows that $f(A)=f(E_{11})\tr A$, proving existence of $c$; evaluating at $E_{11}$ proves uniqueness. Conversely a scalar multiple of trace has the required cyclic identity. Finally $f(I_n)=cn$; if this equals $n$ and $n\ne0$ in the field, cancellation gives $c=1$. That hypothesis cannot be dropped: over $\mathbb F_2$ with $n=2$, the zero functional has $f(I_2)=0=2$ but is not trace. This theorem concerns a span of commutators, not an assertion that a single commutator represents every traceless matrix. For $n=0$, there is only the zero functional and the zero matrix; the span identity remains valid, but the scalar $c$ is not unique.
\end{proof}

\Kentry{04.rank_factorization}{Full-rank column-row factorization}
\begin{theorem}
If $A$ is $m$-by-$n$ of rank $r$, then $A=CR$ for an $m$-by-$r$ matrix $C$ with independent columns and an $r$-by-$n$ matrix $R$ with independent rows.
\end{theorem}
\begin{proof}
Choose a basis $c_1,\ldots,c_r$ of the column space and let these be the columns of $C$. Express each column of $A$ uniquely in this basis; its coordinate columns form $R$, and columnwise multiplication gives $A=CR$. The map $C:\K^r\to\K^m$ is injective. Consequently $r=\rank A=\rank R$ by rank preservation under injective postcomposition. The row-rank theorem then makes the $r$ rows of $R$ independent. If $r=0$, $A=0$, and the matrices of shapes $m\times0$ and $0\times n$ have product zero with vacuously independent columns and rows.
\end{proof}

% END modules/M05.04.tex


% BEGIN modules/M05.01.tex
\chapter{Linear systems and elimination}

\Kentry{01.system}{Linear systems and augmented matrices}
For $A\in\K^{m\times n}$ and $b\in\K^m$, the system $Ax=b$ asks for $x\in\K^n$ satisfying $\sum_{j=1}^n A_{ij}x_j=b_i$ for each row $i$. The augmented matrix is $[A\mid b]$. A system is homogeneous when $b=0$. If $m=0$ every $x$ is a solution. If $n=0$ there is one possible empty vector, and it is a solution exactly when $b=0$. These conventions are consequences of empty finite sums.

\Kentry{01.row_ops}{Elementary row operations}
The permitted operations are interchange of two distinct rows, multiplication of one row by $a\ne0$, and addition of $a$ times a different row to one row. Their inverses are respectively the same interchange, multiplication by $a^{-1}$, and addition of $-a$ times the other row. Applying an operation to $I_m$ gives an invertible elementary matrix $E$. Directly multiplying an arbitrary matrix by $E$ performs exactly that operation. For example, $E=I_m+aE_{ij}$, $i\ne j$, adds $a$ times row $j$ to row $i$, and $(I_m+aE_{ij})(I_m-aE_{ij})=I_m$ since $E_{ij}^2=0$.

\Kentry{01.solution_invariance}{Preservation of the full solution set}
\begin{proposition}
An elementary row operation on the entire augmented matrix preserves exactly the solutions of the system.
\end{proposition}
\begin{proof}
If $Ax=b$, multiplication by $E$ gives $(EA)x=Eb$. Conversely, applying $E^{-1}$ to that equality gives $Ax=b$. Iterating this equivalence proves the assertion for a finite sequence. Operating on $A$ without simultaneously operating on $b$ does not have this justification.
\end{proof}

\Kentry{01.echelon}{Echelon and reduced echelon form}
A matrix is in row-echelon form when its nonzero rows precede its zero rows and the first nonzero column positions of consecutive nonzero rows strictly increase. These first nonzero entries are pivots; all entries below a pivot are zero because later rows start farther right. The form is reduced when every pivot is $1$ and is the sole nonzero entry in its column. For an augmented matrix, the rightmost column is also included in these definitions; a pivot there can indicate inconsistency of the underlying system.

\Kentry{01.elimination}{Gaussian and Gauss--Jordan elimination}
\begin{theorem}
Every finite matrix over a field can be transformed by elementary row operations to echelon form and then to reduced echelon form.
\end{theorem}
\begin{proof}
Maintain an unprocessed submatrix consisting of rows $r,\ldots,m$ and columns $c,\ldots,n$. If all its entries are zero, stop. Otherwise choose the leftmost column with a nonzero entry, interchange its nonzero-entry row with row $r$, and divide row $r$ by the nonzero pivot. Subtract from each lower row its entry in the pivot column times row $r$. Previously processed pivot columns stay unchanged: the current and lower rows were already zero in them. Advance to row $r+1$ and to the column after the new pivot. At every stage the processed row count increases, and the pivot column strictly increases; therefore at most $\min(m,n)$ stages occur. Skipped columns are zero in all still-unprocessed rows. The result has the stated echelon properties. Finally work upward from the last pivot, subtracting multiples of its row from all rows above it. Later pivot columns in that row are already zero, so previously cleared pivot columns stay cleared. This finite backward process gives reduced form. Only nonzero pivots are inverted.
\end{proof}

\Kentry{01.rref_unique}{Uniqueness of reduced row-echelon form}
\begin{theorem}
Every matrix has a unique reduced row-echelon matrix in its row-equivalence class. Two matrices of the same shape are row-equivalent exactly when these reduced forms agree.
\end{theorem}
\begin{proof}
Row operations preserve the row space, since each new row is a combination of old rows and the inverse operations give the reverse containment. Let $R$ be reduced, with nonzero rows $r_1,\ldots,r_k$ and pivot columns $p_1<\cdots<p_k$. In any nonzero combination $\sum_i a_i r_i$, let $i_0$ be the first index with $a_{i_0}\ne0$. Entries before $p_{i_0}$ vanish, while its entry in column $p_{i_0}$ is $a_{i_0}$. Hence the possible leading positions of nonzero vectors in the row space are exactly the $p_i$. They are intrinsic to that space.

If $R,S$ are reduced and have the same row space, their nonzero rows therefore have the same ordered pivot columns. For each index $i$, $r_i-s_i$ lies in this common row space and has zero entries at every pivot column, since each reduced row has the corresponding Kronecker-delta pivot coordinates. A nonzero row-space vector cannot have all these coordinates zero by the preceding leading-position argument. Thus $r_i=s_i$. Remaining rows are zero, so $R=S$. Existence was proved by elimination. If two reduced forms coincide, reversing one elimination sequence and composing it with the other gives row equivalence of the original matrices; the converse follows from preservation of row space and uniqueness.
\end{proof}

\Kentry{01.consistency}{Consistency criteria}
\begin{theorem}
A system is consistent exactly when an echelon form of its augmented matrix has no row $(0,\ldots,0,c)$ with $c\ne0$. Equivalently $b$ belongs to the column span of $A$.
\end{theorem}
\begin{proof}
Such a row asserts $0=c$, so prevents a solution. If no such row exists, assign arbitrary values, for example zero, to the nonpivot unknowns. Starting at the last nonzero row and moving upward, solve for its pivot unknown using the already assigned later unknowns and the nonzero pivot coefficient. Every zero row then reads $0=0$. This constructs a solution to the echelon system, hence to the original by solution invariance. Finally $Ax=\sum_jx_jA_{\cdot j}$, so existence of $x$ is exactly membership in the column span.
\end{proof}

\Kentry{01.parametrization}{Free variables and all solutions}
Let a consistent augmented reduced form have coefficient pivots $p_1,\ldots,p_r$, right-hand entries $c_i$, and free-column set $F$. The complete parametrization is
\[x_j=t_j\quad(j\in F),\qquad x_{p_i}=c_i-\sum_{j\in F}R_{ij}t_j.\]
Substitution satisfies each pivot row and each zero row. Conversely every solution must satisfy these equations, and its free coordinates determine all the $t_j$ uniquely. Thus the solution set is in bijection with $\K^{n-r}$. In the homogeneous case the formula is linear; its inverse reads off the free coordinates. It is therefore a linear isomorphism, proving that the solution space has dimension $n-r$, including $r=0$ and $r=n$.

\Kentry{01.affine_solutions}{Affine structure of an inhomogeneous system}
If $x_0$ is one solution, then $Ax=b$ is equivalent to $A(x-x_0)=0$. Hence all solutions are exactly $x_0+\ker A$. If no solution exists the solution set is empty. A nonempty such coset is a subspace exactly when it contains zero, equivalently when $b=0$: a subspace must contain zero, and if $b=0$ it is precisely the kernel. This separates affine solution sets from vector spaces.

\Kentry{01.pivot_rank}{Pivots and rank}
\begin{theorem}
The number of coefficient pivots in an echelon form of $A$ equals its row rank and column rank. The corresponding pivot columns of the original $A$ form a basis of its column space.
\end{theorem}
\begin{proof}
Nonzero echelon rows are independent: in a zero combination the earliest pivot with nonzero coefficient could not cancel. They span the row space, so their number is its dimension. Row operations preserve row space, and row rank equals column rank by \kref{04.rank_bridge}. If $R=EA$ is reduced, the pivot columns of $R$ are the standard coordinate columns in its nonzero-row block; they are independent and every nonpivot column is a combination of them. Multiplication by the invertible $E^{-1}$ preserves these relations and independence, taking those columns back to the original pivot columns of $A$. This proves the original-column assertion, which would not follow from simply relabeling the altered columns as original ones.
\end{proof}

\Kentry{01.inverse_algorithm}{Computing inverses by elimination}
\begin{theorem}
For square $A$, row reduction of $[A\mid I]$ reaches $[I\mid B]$ exactly when $A$ is invertible, and then $B=A^{-1}$.
\end{theorem}
\begin{proof}
The accumulated row operations form an invertible $E$ and produce $[EA\mid E]$. If the left block is $I$, then $EA=I$; the equal-dimension result makes $A$ invertible and $E=A^{-1}$, so $B=E$. Conversely if $A$ is invertible it has rank $n$, so its reduced form has a pivot in every column and is $I$. Thus reducing the left block produces the required result. A missing pivot forces rank below $n$, and hence prevents invertibility. For $n=0$, both blocks are empty identities and the assertion remains valid.
\end{proof}

% END modules/M05.01.tex


% BEGIN modules/M05.05.tex
\chapter{Determinants and their interpretations}
The algebraic determinant identities are proved over a commutative ring $R$ with identity. Rank, division by an arbitrary nonzero scalar, positivity and volume are used only under the additional hypotheses explicitly stated.

\Kentry{05.alternating}{Alternating multilinear forms}
An $n$-linear map $f:M^n\to R$ is alternating if it vanishes whenever two input positions contain the same vector. For distinct positions occupied by $u+v,u+v$, expansion gives
\[0=f(\ldots,u,\ldots,v,\ldots)+f(\ldots,v,\ldots,u,\ldots),\]
because the two equal-input terms vanish. Thus an interchange negates $f$. Decomposing a permutation into interchanges gives its permutation-sign rule; independence of the decomposition is the sign identity from the foundational chapter. In a field of characteristic other than two, a sign-rule form is alternating, because interchanging equal arguments gives $f=-f$ and $2$ is invertible. This converse fails in characteristic two: $f(x,y)=xy$ on a one-dimensional space is symmetric, hence satisfies the sign rule there, but $f(1,1)=1$. Alternation, rather than just the sign rule, is the uniform definition.

\Kentry{05.det_definition}{The determinant}
For $A\in R^{n\times n}$ define
\[\det A=\sum_{\sigma\in S_n}\operatorname{sgn}(\sigma)
                    \prod_{j=1}^n A_{\sigma(j),j}.\]
For $n=0$, $S_0$ has one empty permutation, the empty product is $1$, and $\det(\varnothing)=1$. For each $n$, this function is $R$-linear in each column and takes value $1$ on $I_n$: the identity permutation is the only term not forced to zero there. If columns $j,k$ agree, pair each term indexed by $\sigma$ with that indexed by $\sigma\circ(j\ k)$. Their scalar products agree and their signs are opposite, so their sum is zero. The pairing has no fixed points. This proves alternation over every commutative ring, including characteristic two.

\Kentry{05.det_uniqueness}{Existence and uniqueness of the normalized form}
\begin{theorem}
The determinant is the unique alternating $n$-linear form on $n$ columns in $R^n$ with value $1$ on $(e_1,\ldots,e_n)$. More generally every such form $f$ satisfies
\[f(a_1,\ldots,a_n)=f(e_1,\ldots,e_n)\det[a_1\ \cdots\ a_n].\]
\end{theorem}
\begin{proof}
Expand each $a_j=\sum_i A_{ij}e_i$ using multilinearity. Each term with a repeated index vanishes by alternation. The remaining index tuples are exactly $(\sigma(1),\ldots,\sigma(n))$. Their form values are $\operatorname{sgn}(\sigma)f(e_1,\ldots,e_n)$ by the sign rule. Summing these terms gives the formula. Existence of the normalized form was verified directly above. When $n=0$, the form has one empty tuple as domain and is determined by its value there, so the same conclusion holds.
\end{proof}

\Kentry{05.det_product}{Multiplication and transposition}
\begin{theorem}
For square matrices of equal size over $R$,
\[\det(AB)=\det A\det B,\qquad \det(A^t)=\det A.\]
\end{theorem}
\begin{proof}
For fixed $A$, the function $(b_1,\ldots,b_n)\mapsto\det(Ab_1,\ldots,Ab_n)$ is alternating and multilinear, with value $\det A$ on the standard basis. The general uniqueness formula makes it $\det A\det B$. To prove the transpose identity, rewrite its Leibniz sum as
\[\sum_{\sigma}\operatorname{sgn}(\sigma)\prod_j A_{j,\sigma(j)}.
\]
Substitute $\tau=\sigma^{-1}$ and reorder the product. Since inverse permutations have equal signs, this is the original determinant sum. Both proofs include empty matrices.
\end{proof}

\Kentry{05.det_row_rules}{Row and column rules}
Interchanging columns negates determinant by alternation; scaling one column scales it by multilinearity. Adding $a$ times a different column to a column changes the determinant by $a$ times a determinant with two equal columns, which is zero. Apply the transpose identity to obtain exactly the same three row rules. If columns are linearly dependent over a field, solve a dependence relation for a column with nonzero coefficient. Multilinearity expresses the determinant as a sum of terms with repeated columns, all zero. The analogous row statement follows by transposition. No assertion that a negative sign is different from a positive sign is made in characteristic two.

\Kentry{05.cofactor}{Cofactor expansion}
For $n\ge1$, let $A^{\widehat i\widehat j}$ be obtained by deleting row $i$ and column $j$ and define $C_{ij}=(-1)^{i+j}\det A^{\widehat i\widehat j}$. Then
\[\det A=\sum_i A_{ij}C_{ij}\quad\text{for fixed }j,
 \qquad \det A=\sum_j A_{ij}C_{ij}\quad\text{for fixed }i.\]
Indeed, group the Leibniz terms by $\sigma(j)=i$. After deleting that assigned row and column, the other assignments are permutations on the ordered remaining sets. Moving row $i$ and column $j$ to the last positions uses $(n-i)+(n-j)$ interchanges, of parity $i+j$; therefore each remaining permutation sign gains precisely $(-1)^{i+j}$. The grouped sum is $A_{ij}C_{ij}$. The row formula follows by transposition. For $n=1$ the deleted minor is empty and its determinant is $1$, as required.

\Kentry{05.adjugate}{The adjugate identities}
For $n\ge1$, set $\operatorname{adj}(A)_{ij}=C_{ji}$. Then
\[A\operatorname{adj}(A)=\operatorname{adj}(A)A=(\det A)I_n.\]
The $(i,j)$ entry of the first product is $\sum_kA_{ik}C_{jk}$. If $i=j$, this is expansion along row $j$. If $i\ne j$, it is expansion along row $j$ of the matrix obtained by replacing that row by row $i$; its minors obtained by deleting row $j$ are unchanged, but its determinant is zero because two rows coincide. The other product follows by the same column argument. At $n=1$, the adjugate is $[1]$. Define the adjugate of the empty matrix to be the unique empty matrix, for which both identities are equalities of empty matrices. No invertibility was assumed.

\begin{lemma}[Rank and minors over a field]\label{aux:rank-minors}
For an $m$-by-$n$ matrix $A$ over a field, its rank is the largest size of a square submatrix with nonzero determinant. Size zero is allowed: its determinant is $1$.
\end{lemma}
\begin{proof}
Put $r=\rank A$. If $k>r$, any $k$ selected columns of $A$ are dependent. Restricting their dependence relation to any $k$ rows leaves a nontrivial relation among the square submatrix's columns, so its determinant is zero by \kref{05.det_row_rules}. To obtain an $r$-minor when $r>0$, select $r$ independent columns and let $C$ be the resulting $m$-by-$r$ matrix. Its column rank is $r$, hence its row rank is $r$ by \kref{04.rank_bridge}. Select $r$ independent rows of $C$. They form an invertible square matrix: its rows span the $r$-dimensional row-coordinate space, and the row-rank and column-rank equality implies full column rank, hence invertibility. The determinant criterion in \kref{05.det_invertible} makes this determinant nonzero. That criterion uses only the adjugate identities above, not this lemma, so this invocation is noncircular. If $r=0$, every entry is zero, all positive-size minors vanish, and the empty minor supplies the asserted maximum. This includes $m=0$ or $n=0$.
\end{proof}

\begin{corollary}[Ranks and iteration of the adjugate over a field]
For $A\in\K^{n\times n}$, $n\ge1$,
\[
\rank\operatorname{adj}(A)=
\begin{cases}
n,&\rank A=n,\\
1,&\rank A=n-1,\\
0,&\rank A\le n-2.
\end{cases}
\]
Also $\det\operatorname{adj}(A)=(\det A)^{n-1}$, and, for $n\ge2$,
\[
\operatorname{adj}(\operatorname{adj}(A))=(\det A)^{n-2}A.
\]
These rank assertions are field statements, distinct from the ring-valued adjugate identities above.
\end{corollary}
\begin{proof}
If $A$ is invertible, its adjugate equals $(\det A)A^{-1}$, so it has rank $n$. Multiplicativity of determinant and scaling all $n$ columns give
$\det\operatorname{adj}(A)=(\det A)^n\det(A^{-1})=(\det A)^{n-1}$.
If $\rank A=n-1$, Lemma~\ref{aux:rank-minors} supplies a nonzero cofactor, so $\operatorname{adj}(A)\ne0$. The identity $A\operatorname{adj}(A)=0$ places its image in the one-dimensional kernel of $A$, so its rank is exactly one. If $\rank A\le n-2$, every $(n-1)$-minor vanishes by that lemma, so the adjugate is zero. The middle case includes $n=1,A=0$, whose adjugate is $[1]$. For singular $A$ with $n\ge2$, the adjugate has rank at most one and is singular (at $n=2$ its rank is at most $1$); both sides of the determinant identity are zero. At $n=1$ the adjugate determinant and the zeroth power on the right are both one.

For invertible $A$ with $n\ge2$, applying the inverse formula to its adjugate yields
\[
\operatorname{adj}(\operatorname{adj}(A))
 = (\det A)^{n-1}\bigl((\det A)A^{-1}\bigr)^{-1}
 = (\det A)^{n-2}A.
\]
For singular $A$ and $n\ge3$, its adjugate has rank at most $1\le n-2$, so applying the rank cases to it makes its adjugate zero; the right side also vanishes since $n-2>0$. For $n=2$, direct substitution in
$\operatorname{adj}\begin{pmatrix}a&b\\c&d\end{pmatrix}=\begin{pmatrix}d&-b\\-c&a\end{pmatrix}$
shows that applying the operation twice returns $A$, including singular $A$. Here the power $(\det A)^0$ is one, even when $\det A=0$.
\end{proof}

\Kentry{05.det_invertible}{The determinant criterion for invertibility}
\begin{theorem}
Over a commutative ring, $A$ is invertible exactly when $\det A$ is a unit. In particular, over a field it is invertible exactly when $\det A\ne0$. In that case $\det(A^{-1})=(\det A)^{-1}$.
\end{theorem}
\begin{proof}
If $AB=BA=I$, multiplicativity gives $\det A\det B=1$, so the determinant is a unit. Conversely, if $d=\det A$ is a unit, the adjugate identities give a two-sided inverse $d^{-1}\operatorname{adj}(A)$. Multiplicativity applied to $AA^{-1}=I$ proves the final identity. The field specialization uses the fact that precisely the nonzero field elements are units.
\end{proof}

\Kentry{05.cramer}{Cramer's rule}
Let $A$ be square over a field and $\det A\ne0$. There is a unique solution $x=A^{-1}b$. Let $A[j\leftarrow b]$ denote replacement of column $j$ by $b$. Since $b=\sum_kx_k A_{\cdot k}$, multilinearity and repeated-column vanishing give
\[\det A[j\leftarrow b]=x_j\det A.\]
Division by the nonzero determinant yields $x_j=\det A[j\leftarrow b]/\det A$. For a singular matrix this division is unavailable, and neither existence nor uniqueness follows from this formula. The empty system has its unique empty solution and no coordinate formulas to check.

\Kentry{05.triangular_block}{Triangular and block-triangular determinants}
For an upper-triangular $A$, a nonzero Leibniz term requires $\sigma(j)\le j$ for every $j$. Since the sums of both sides over $j$ agree, all these inequalities are equalities; only the identity permutation contributes. Thus $\det A=\prod_jA_{jj}$. The lower-triangular case follows by transposition. More generally, for square diagonal blocks,
\[\det\begin{pmatrix}B&C\\0&D\end{pmatrix}=\det B\det D.\]
If $B$ has size $r$, each nonzero permutation term assigns the first $r$ columns to the first $r$ rows, because all entries below that block are zero. As these sets both have $r$ elements, the assignment is a permutation on that block and the remaining assignment is a permutation on the other block. Its sign is the product of their signs, so the sum factors into the two determinant sums. Repeated application proves the finite-block formula; zero-sized blocks have determinant $1$.

\Kentry{05.operator_det}{Determinant of an operator and the top exterior power}
For $T\in\End(V)$ with $\dim V=n$, define $\det T=\det[T]_B$. Under a basis change, $[T]_{B'}=P^{-1}[T]_BP$, and multiplicativity gives the same scalar. Thus the definition is basis-independent, and also gives $\det(ST)=\det S\det T$.

The independent construction and basis theorem for exterior powers are proved in \kref{15.exterior} and \kref{15.exterior_basis}. They give a one-dimensional space $\bigwedge^nV$ with basis $b_1\wedge\cdots\wedge b_n$, including $\bigwedge^0\{0\}=\K$. Expanding
$Tb_1\wedge\cdots\wedge Tb_n$ multilinearly eliminates repeated basis indices; the permutation-sign rule leaves precisely
\[(\bigwedge^nT)(b_1\wedge\cdots\wedge b_n)
 =(\det[T]_B)b_1\wedge\cdots\wedge b_n.\]
Hence this action is the same scalar determinant, rather than a second unrelated definition. This argument requires only the exterior basis and sign rules, not the determinant-line identification being proved.

\Kentry{05.orientation}{Orientation}
Two ordered bases of a finite-dimensional real space have the same orientation if the coordinate-conversion matrix between them has positive determinant. Identity has positive determinant; inversion preserves positivity; multiplication of two positive determinants is positive. Thus this is an equivalence relation. For positive dimension there are exactly two classes: interchanging two basis vectors when possible, or negating one vector in dimension one, changes the sign, while any pair of bases has either positive or negative conversion determinant. In dimension zero there is only the empty basis and one class. For invertible real $T$, the determinant of the conversion from $B$ to $T(B)$ is $\det T$, so its sign characterizes preservation or reversal of orientation in positive dimension.

\Kentry{05.volume}{Signed volume and ordinary volume}
Fix an oriented orthonormal basis of a real Euclidean $n$-space. The signed volume of the ordered parallelepiped $(v_1,\ldots,v_n)$ is the determinant of its coordinate columns; its ordinary parallelepiped volume is the absolute value. Replacing the oriented orthonormal reference basis changes coordinates by an orthogonal matrix of positive determinant. The identity $Q^tQ=I$ gives $(\det Q)^2=1$, and positivity gives $\det Q=1$; hence signed volume is independent of this choice. Multiplication of column matrices shows that $T$ multiplies signed volume by $\det T$ and ordinary volume by $|\det T|$.

For agreement with Lebesgue measure, the measurable parallelepiped is the image of $[0,1]^n$. The independently established measure bridge (Theorem~\ref{aux:volume-change}) gives its measure $|\det A|$ when $A$ is invertible. If $A$ is singular, its image lies in a proper linear subspace, which is a null set by that bridge, and both its measure and determinant vanish. The same bridge proves $\mu(AE)=|\det A|\mu(E)$ for invertible $A$ and measurable $E$. For singular $A$, state the null-image assertion directly rather than use an undefined expression $0\cdot\infty$.

\Kentry{05.hadamard}{Hadamard's inequality and equality cases}
\begin{theorem}
For real or complex square $A$ with columns $v_j$,
\[|\det A|\le\prod_j\norm{v_j}.\]
If every column is nonzero, equality holds exactly when the columns are pairwise orthogonal. If one column is zero, equality holds regardless of the angles of the others.
\end{theorem}
\begin{proof}
The zero-column assertion follows from multilinearity. Assume all columns are nonzero. If they are dependent, determinant vanishes while the norm product is positive, so equality is impossible. Otherwise apply the Gram--Schmidt construction proved in \kref{12.gram_schmidt}: write $v_j=\sum_{i<j}a_{ij}q_i+u_j$, with the $q_i$ orthonormal, $u_j\perp q_i$ and $q_j=u_j/\norm{u_j}$. Thus $A=QR$ where $Q$ is unitary and $R$ is triangular with positive diagonal $\norm{u_j}$. The determinant sum gives $\det Q^\dagger=\overline{\det Q}$; since $Q^\dagger Q=I$, $|\det Q|=1$. Therefore $|\det A|=\prod_j\norm{u_j}$. Pythagoras gives $\norm{v_j}^2=\sum_{i<j}|a_{ij}|^2+\norm{u_j}^2$, so each ratio $\norm{u_j}/\norm{v_j}$ is in $(0,1]$. Their product equals $1$ exactly when every ratio does, equivalently when every $a_{ij}=0$. Then $v_j$ is orthogonal to all its predecessors; conversely pairwise orthogonality makes all these coefficients vanish. The empty matrix gives the equality $1=1$.
\end{proof}

\Kentry{05.vandermonde}{Vandermonde determinants and evaluation}
\begin{theorem}
For scalars $b_1,\ldots,b_n$ in $\R$ or $\C$ and $A_{ij}=b_i^{j-1}$,
\[\det A=\prod_{1\le i<j\le n}(b_j-b_i).\]
Consequently evaluation at distinct nodes is an isomorphism from polynomials of degree less than $n$ to $\K^n$.
\end{theorem}
\begin{proof}
The formula is $1=1$ for $n=0,1$. Suppose it is established at size $n-1$. If the first $n-1$ nodes contain a repetition, both sides vanish, because the matrix has repeated rows. Otherwise regard the determinant as a polynomial $f(t)$ in its last node. Expansion along the last row shows degree at most $n-1$ and leading coefficient equal to the size-$n-1$ determinant. The values $t=b_1,\ldots,b_{n-1}$ make two rows equal, so all are roots. Repeated application of the polynomial factor theorem at these distinct roots gives
$f(t)=c\prod_{i<n}(t-b_i)$; comparing the leading coefficient gives $c=\det A_{n-1}$. The induction hypothesis proves the formula at size $n$. For distinct nodes the product is nonzero, so the evaluation matrix in the monomial basis is invertible. When $n=0$, the polynomial domain is the zero space and this is the unique isomorphism to $\K^0$.
\end{proof}

% END modules/M05.05.tex


% BEGIN modules/M05.06.tex
\chapter{Quotients, direct sums, and invariant structure}

\Kentry{06.quotient}{Quotient vector spaces}
For a subspace $U\subseteq V$, define $v\sim w$ when $v-w\in U$. Reflexivity uses $0\in U$, symmetry uses closure under negation, and transitivity uses closure under addition. Its equivalence class is $v+U$. Define
\[(v+U)+(w+U)=(v+w)+U,\qquad a(v+U)=av+U.\]
If $v'=v+u$ and $w'=w+u'$ with $u,u'\in U$, then $(v'+w')-(v+w)=u+u'\in U$ and $av'-av=au\in U$. Thus the operations do not depend on representatives. Their vector-space axioms are inherited by taking representatives and using the corresponding equalities in $V$. The resulting space is $V/U$, its zero vector is the coset $U$, and its additive inverse operation is $-(v+U)=-v+U$.

\Kentry{06.quotient_map}{The quotient map}
The map $q:V\to V/U$, $q(v)=v+U$, is linear by the defining operation formulas. It is surjective because every coset has a representative. Its kernel consists exactly of the $v$ with $v+U=U$, equivalently $v\in U$. This distinguishes the zero coset from a chosen complement to $U$, which is a different kind of object.

\Kentry{06.quotient_lift}{The universal property of a quotient}
\begin{theorem}
A linear $T:V\to W$ factors uniquely as $T=\overline Tq$ with linear $\overline T:V/U\to W$ if and only if $U\subseteq\ker T$.
\end{theorem}
\begin{proof}
If the containment holds, set $\overline T(v+U)=T(v)$. Changing the representative changes $v$ by an element of $U$, which maps to zero, so the definition is well-defined. The coset operation formulas and linearity of $T$ show that $\overline T$ is linear. It satisfies the required factorization, and surjectivity of $q$ forces uniqueness of every value of $\overline T$. Conversely, if $T=\overline Tq$ and $u\in U$, then $q(u)=0$, so $T(u)=\overline T(0)=0$.
\end{proof}

\Kentry{06.quotient_dimension}{Dimension of a quotient}
\begin{theorem}
If $V$ is finite-dimensional, then
\[\dim(V/U)+\dim U=\dim V.\]
\end{theorem}
\begin{proof}
Choose a basis $u_1,\ldots,u_k$ of $U$ and extend it by $v_1,\ldots,v_r$ to a basis of $V$. The cosets $v_j+U$ span the quotient because the $u_i$ disappear there. If a combination $\sum_j a_jv_j$ belongs to $U$, express it in the basis of $U$; independence of the full basis forces all $a_j=0$. Thus these cosets form a basis, and the dimensions are $r,k,k+r$. This proof also treats $U=0$, $U=V$ and $V=0$.
\end{proof}

\Kentry{06.direct_sum}{Internal direct sums}
For subspaces $U_1,\ldots,U_r$ of $V$, the assertion $V=\bigoplus_{i=1}^r U_i$ means every $v$ has a unique expression $v=\sum_i u_i$ with $u_i\in U_i$. Equivalently the sum spans $V$ and the only zero sum $\sum_i u_i=0$ has all components zero. Indeed, subtracting two expressions reduces uniqueness to this zero-sum condition. For two subspaces, the zero-sum condition is equivalent to $U\cap W=0$: from $u+w=0$ one has $u=-w\in U\cap W$; conversely $z\in U\cap W$ gives $z+(-z)=0$. Thus $V=U\oplus W$ exactly when $U+W=V$ and $U\cap W=0$. For $r=0$, the assertion means $V=0$.

\Kentry{06.external_sum}{External sums and products}
For any family $(V_i)_{i\in I}$, define $\bigoplus_i V_i$ to be the space of tuples $(v_i)$ with finite support, with componentwise operations. The support of a sum is contained in the union of the two supports and scalar multiplication does not enlarge support; hence these operations are well-defined. The vector-space laws hold in every component. If $I$ is finite, every tuple has finite support, so this equals the product $\prod_i V_i$.

For a finite family of subspaces $U_i\subseteq V$, the summation map $s:\bigoplus_i U_i\to V$ is linear. It is onto exactly when the subspaces span $V$, and its kernel consists of their zero-sum tuples. Therefore it is an isomorphism exactly when the subspaces form an internal direct sum. This is an explicit map between the external construction and the internal one, rather than equality of their underlying sets.

\Kentry{06.complement}{Complements of finite-dimensional subspaces}
\begin{theorem}
Every subspace $U$ of a finite-dimensional space $V$ has a complement $W$ with $V=U\oplus W$.
\end{theorem}
\begin{proof}
Extend a basis of $U$ to a basis of $V$, and let $W$ be the span of the added vectors. The full basis gives $U+W=V$. If a vector belongs to both, subtract its two basis expansions; independence makes every coefficient zero. Hence $U\cap W=0$. The construction depends on the extension. For example, in $\K^2$ both $\Span(e_2)$ and $\Span(e_1+e_2)$ complement $\Span(e_1)$, so no preferred complement has been specified by vector-space structure alone.
\end{proof}

\Kentry{06.projection}{Idempotents and projections}
\begin{theorem}
If $P\in\End(V)$ satisfies $P^2=P$, then $V=\im P\oplus\ker P$. Conversely each decomposition $V=U\oplus W$ determines a unique idempotent with image $U$ and kernel $W$.
\end{theorem}
\begin{proof}
Write $v=Pv+(v-Pv)$. The first term is in the image, and $P(v-Pv)=Pv-P^2v=0$, so the second is in the kernel. If $x=Py$ also satisfies $Px=0$, then $x=Py=P^2y=Px=0$, establishing directness. Conversely, set $P(u+w)=u$ using the unique decomposition. Sums and scalar multiples decompose componentwise, so $P$ is linear, with exactly the stated image and kernel, and $P^2=P$. Any idempotent acts as identity on its image and zero on its kernel, so these two spaces force the same formula and uniqueness. No inner product or orthogonality is assumed.
\end{proof}

\Kentry{06.sum_dimension}{Dimensions of direct sums}
If $V=\bigoplus_{i=1}^r U_i$ and each $U_i$ is finite-dimensional, concatenate chosen bases of the summands. Decomposing a vector and expanding each component proves spanning. In a zero relation, group terms by their summands. Directness first makes each component zero, and independence within its basis then makes each coefficient zero. Hence the concatenated list is a basis and $\dim V=\sum_i\dim U_i$. The same proof with component injections gives a basis of the external sum; equivalently use its summation isomorphism. For the empty sum the dimension and empty sum of dimensions are both zero.

\Kentry{06.invariant}{Invariant subspaces, restrictions, and induced maps}
A subspace $U$ is invariant under $T\in\End(V)$ when $T(U)\subseteq U$. This permits the restriction $T|_U:U\to U$; the underlying formula is unchanged, but the domain and codomain are now $U$. It also permits the induced map $\widehat T:V/U\to V/U$, $v+U\mapsto T(v)+U$. If $v-w\in U$, invariance gives $T(v)-T(w)\in U$, so this is well-defined; the operation formulas prove linearity. These maps obey $\iota T|_U=T\iota$ and $\widehat Tq=qT$, where $\iota:U\to V$ is inclusion. Restriction and passage to a quotient are distinct constructions even when their dimensions coincide.

\Kentry{06.invariant_decomposition}{Invariant summands and commuting projections}
\begin{theorem}
For $V=\bigoplus_i U_i$, let $P_i$ be the component projections. All summands are $T$-invariant exactly when every $P_i$ commutes with $T$. In that case $T$ is the direct sum of its restrictions to the summands.
\end{theorem}
\begin{proof}
If each $U_i$ is invariant and $v=\sum_i u_i$, then $Tv=\sum_i Tu_i$ is its component decomposition. Thus $P_iTv=Tu_i=TP_iv$. Conversely, if $P_iT=TP_i$ and $u\in U_i$, then $P_i(Tu)=T(P_iu)=Tu$. A vector fixed by $P_i$ belongs to its image $U_i$, so $Tu\in U_i$. Under the external-to-internal summation isomorphism, the map therefore acts componentwise by $T|_{U_i}$, which is the block-diagonal statement.
\end{proof}

\Kentry{06.noninvariant_complement}{A complement need not be invariant}
On $\K^2$ define $N(e_1)=0$ and $N(e_2)=e_1$. The subspace $U=\Span(e_1)$ is invariant. Suppose it had an invariant complement $W$. The dimension formula gives $\dim W=1$, and a nonzero generator has form $w=ae_1+be_2$ with $b\ne0$, since otherwise $W\subseteq U$. Invariance would give $Nw=be_1\in W\cap U=0$, contradicting $b\ne0$. Thus an invariant complement does not exist, despite the existence of ordinary vector-space complements. The example works over every field.

% END modules/M05.06.tex


% BEGIN modules/M05.07.tex
\chapter{Algebraic duality and annihilators}

\Kentry{07.dual}{The algebraic dual}
Define $V^*=\Hom_\K(V,\K)$, with pointwise addition and scalar multiplication. These operations form a vector space by \kref{03.map_algebra}. Its members are linear functionals. No topology is supplied, and no condition of continuity is imposed. This is the algebraic dual, independently of any later inner product used to represent its elements.

\Kentry{07.dual_basis}{Dual coordinates and the dual basis}
\begin{theorem}
For a finite basis $v_1,\ldots,v_n$ of $V$, there exist unique $\phi_i\in V^*$ with $\phi_i(v_j)=\delta_{ij}$, and $\phi_1,\ldots,\phi_n$ form a basis of $V^*$.
\end{theorem}
\begin{proof}
The map-specification theorem supplies the functionals. For $v=\sum_j a_jv_j$, one has $\phi_i(v)=a_i$. For any $f\in V^*$, the functional $\sum_i f(v_i)\phi_i$ agrees with $f$ on every basis vector and hence on $V$, proving spanning. If $\sum_i c_i\phi_i=0$, evaluating at $v_j$ gives $c_j=0$, proving independence. In particular $\dim V^*=\dim V$, and $f(v)=\sum_i f(v_i)\phi_i(v)$. When $V=0$, its dual is also zero and the basis is empty.
\end{proof}

\Kentry{07.annihilator}{Annihilators and preannihilators}
For a subspace $U\subseteq V$, put
\[U^0=\{f\in V^*: f(u)=0\text{ for every }u\in U\}.
\]
It is a subspace: linear combinations of vanishing functionals vanish on the same $U$. For a family $F\subseteq V^*$, put ${}^0F=\{v\in V:f(v)=0\text{ for all }f\in F\}=\bigcap_{f\in F}\ker f$, a subspace of $V$. These operations reverse containments. They do not have the same ambient space: $U^0$ is in $V^*$, while ${}^0F$ is in $V$.

\Kentry{07.annihilator_dimension}{Dimension of an annihilator}
\begin{theorem}
If $\dim V=n$ and $\dim U=k$, then $\dim U^0=n-k$.
\end{theorem}
\begin{proof}
Extend a basis $v_1,\ldots,v_k$ of $U$ to $v_1,\ldots,v_n$ of $V$ and take its dual basis. Since every functional is $\sum_i f(v_i)\phi_i$, it vanishes on $U$ exactly when its first $k$ coefficients vanish. Thus $\phi_{k+1},\ldots,\phi_n$ form a basis of $U^0$. This also shows $V^0=0$ and $\{0\}^0=V^*$.
\end{proof}

\Kentry{07.dual_map}{The dual map and reversed composition}
For $T:V\to W$, define $T^\vee:W^*\to V^*$ by $T^\vee(f)=f\circ T$. Composition preserves linearity, so its values are functionals. Pointwise operations show $T^\vee$ itself is linear. If $S:W\to X$, then
\[(ST)^\vee(g)=g\circ S\circ T=T^\vee(S^\vee(g)),\]
so $(ST)^\vee=T^\vee S^\vee$, and $(\id_V)^\vee=\id_{V^*}$. Also $(aS+bT)^\vee=aS^\vee+bT^\vee$ for maps with common source and target. This algebraic operation uses no conjugation, including over $\C$; it is distinct from the inner-product adjoint $T^\dagger$.

\Kentry{07.dual_kernel_image}{Kernel and image of a dual map}
\begin{theorem}
For every linear $T:V\to W$,
\[\ker T^\vee=(\im T)^0\subseteq W^*.
\]
If $V,W$ are finite-dimensional, then
\[\im T^\vee=(\ker T)^0\subseteq V^*.
\]
\end{theorem}
\begin{proof}
A functional $f\in W^*$ satisfies $T^\vee f=0$ exactly when $f(Tv)=0$ for every $v$, which is exactly annihilation of $\im T$. For the second identity, every $f\circ T$ vanishes on $\ker T$, proving one containment. Conversely let $g\in V^*$ vanish on $\ker T$. Define $h:\im T\to\K$ by $h(Tv)=g(v)$. If $Tv=Tw$, then $v-w\in\ker T$ and $g(v)=g(w)$, so $h$ is well-defined; linearity follows by representatives. Extend a basis of $\im T$ to one of $W$, prescribe $h$ on the former and zero on the added basis vectors, and extend linearly to $f\in W^*$. Then $g=f\circ T=T^\vee f$, proving the reverse containment. This proof does not infer equality merely from a dimension count.
\end{proof}

\Kentry{07.dual_rank}{Rank and injectivity-surjectivity duality}
\begin{theorem}
For finite-dimensional $V,W$, $\rank T^\vee=\rank T$. Moreover $T$ is injective exactly when $T^\vee$ is surjective, and $T$ is surjective exactly when $T^\vee$ is injective.
\end{theorem}
\begin{proof}
Rank-nullity for $T^\vee$ and the preceding kernel formula give
\[\rank T^\vee=\dim W^*-\dim(\im T)^0
 =\dim W-(\dim W-\dim\im T)=\rank T.\]
Annihilator dimension shows $(\ker T)^0=V^*$ exactly when $\ker T=0$. Together with the image formula this proves the first equivalence. Similarly $(\im T)^0=0$ exactly when $\dim\im T=\dim W$, equivalently $\im T=W$. Together with the kernel formula this proves the second equivalence. The rank argument above uses finite dimension. The two equivalences for the algebraic dual hold in arbitrary dimension: a functional on $\im T$ extends to $W$ by extending a basis, and a proper subspace of $W$ is annihilated by a nonzero functional obtained in the same way. Thus $T^\vee$ is onto exactly when $T$ is injective, and $T^\vee$ is injective exactly when $T$ is onto.
\end{proof}

\Kentry{07.biddual}{The canonical evaluation isomorphism}
\begin{theorem}
For finite-dimensional $V$, the map $J_V:V\to V^{**}$ given by $J_V(v)(f)=f(v)$ is a canonical linear isomorphism.
\end{theorem}
\begin{proof}
For fixed $v$, evaluation is linear in $f$, so it belongs to $V^{**}$. Evaluation is also linear in $v$, giving linearity of $J_V$. If $J_V(v)=0$, every coordinate functional of a chosen finite basis vanishes on $v$, so $v=0$. Thus it is injective, and the twice-applied dual dimension theorem gives $\dim V^{**}=\dim V$, hence surjectivity. More explicitly, if $F\in V^{**}$, set $v=\sum_iF(\phi_i)v_i$. Writing an arbitrary $f=\sum_i f(v_i)\phi_i$ gives $f(v)=\sum_i f(v_i)F(\phi_i)=F(f)$, so $J_V(v)=F$. The formula defining $J_V$ contains no choice of basis, proving its canonical character. For $T:V\to W$, evaluation gives $T^{\vee\vee}J_V=J_WT$ on every functional.
\end{proof}
If $V$ is infinite-dimensional, injectivity can instead be proved using the optional choice-based basis extension: extend a nonzero vector to a basis and define a functional taking value $1$ there and $0$ on the other basis elements. Surjectivity is not asserted for that variant.

\Kentry{07.double_annihilator}{Double annihilation}
\begin{theorem}
For $U\subseteq V$ with $V$ finite-dimensional, ${}^0(U^0)=U$. Equivalently the annihilator of $U^0$ in $V^{**}$ is $J_V(U)$.
\end{theorem}
\begin{proof}
Every $u\in U$ is killed by every member of $U^0$, giving one containment. Extend a basis of $U$ to $v_1,\ldots,v_n$ of $V$. If $v\notin U$, some coordinate after the initial $U$ basis is nonzero. Its coordinate functional belongs to $U^0$ but does not vanish on $v$, proving the reverse containment. By definition, $J_V(v)$ annihilates $U^0$ exactly when every $f\in U^0$ has $f(v)=0$. Since $J_V$ is onto, this proves the equivalent bidual formulation. It is an identification through $J_V$, not literal equality of subsets in different spaces.
\end{proof}

\Kentry{07.dual_matrix}{Matrices of dual maps}
Let $B=(v_j)$ and $C=(w_i)$ be bases, with dual bases $(\phi_j)$ and $(\psi_i)$. If $T(v_j)=\sum_iA_{ij}w_i$, then
\[(T^\vee\psi_i)(v_j)=\psi_i(Tv_j)=A_{ij}.
\]
Expansion in the dual basis gives $T^\vee\psi_i=\sum_jA_{ij}\phi_j$. Thus the matrix of $T^\vee$ has row $j$, column $i$ entry $A_{ij}$ and equals $A^t$. No complex conjugates occur. This proves transposition as a coordinate representation of the algebraic dual, independently of the conjugate transpose used for inner-product adjoints.

\Kentry{07.quotient_dual}{The dual of a quotient}
\begin{theorem}
Pullback by $q:V\to V/U$ gives a canonical isomorphism $(V/U)^*\to U^0$, $h\mapsto h\circ q$, without a finite-dimensional assumption.
\end{theorem}
\begin{proof}
The composite vanishes on $U$, and the pullback is linear. It is injective because $q$ is onto: if $h\circ q=0$, every coset value of $h$ is zero. For $f\in U^0$, set $h(v+U)=f(v)$. Difference of representatives lies in $U$, so this is well-defined, and the quotient operation formulas prove linearity. It is inverse to pullback in both directions by evaluating on representatives. Consequently the isomorphism is canonical and is not obtained by choosing a complement or comparing dimensions.
\end{proof}

\Kentry{07.hyperplanes}{Hyperplanes and kernels of functionals}
\begin{theorem}
For an arbitrary vector space $V$ over a field, the following are equivalent for a subspace $U$: it is maximal among proper subspaces; $\dim(V/U)=1$; it is the kernel of a nonzero linear functional. Such a subspace is called a hyperplane. Two nonzero functionals have the same kernel precisely when one is a nonzero scalar multiple of the other.
\end{theorem}
\begin{proof}
If $U$ is maximal proper, choose $v\notin U$. The subspace $U+\K v$ strictly contains $U$, hence equals $V$. Also $U\cap\K v=0$, because a nonzero multiple of $v$ in $U$ would put $v$ in $U$. Thus every vector has a unique expression $u+av$, and $f(u+av)=a$ is a well-defined nonzero linear functional with kernel $U$; the coset $v+U$ is a basis of $V/U$.

Conversely, if $\dim(V/U)=1$, choose $v+U$ as a basis of the quotient. Every $x\in V$ has a unique expression $u+av$, by the same quotient-coordinate argument, so the displayed functional again has kernel $U$. If $U=\ker f$ for $f\ne0$, choose $v$ with $f(v)=1$ by scaling a vector with nonzero value. Then $x=(x-f(x)v)+f(x)v$ is a decomposition in $U\oplus\K v$, proving $\dim(V/U)=1$. Any subspace strictly containing $U$ contains some $x$ with $f(x)\ne0$; its decomposition shows that it contains $v$, hence all of $V$. Thus $U$ is maximal proper. No infinite basis extension or subtraction of infinite dimensions is used.

If $\ker f=\ker g$ and $f(v)=1$, then $c=g(v)\ne0$. Since $x-f(x)v$ lies in both kernels, $g(x)=c f(x)$ for every $x$. Conversely, multiplication by a nonzero scalar leaves the kernel unchanged. The zero-dimensional space has neither hyperplanes nor nonzero functionals.
\end{proof}

\Kentry{07.functional_constraints}{Redundant linear constraints}
\begin{theorem}
For a finite list of functionals $f_1,\ldots,f_r,g$ on an arbitrary vector space $V$,
\[g\in\Span(f_1,\ldots,f_r)\quad\Longleftrightarrow\quad
 \bigcap_{i=1}^r\ker f_i\subseteq\ker g.\]
\end{theorem}
\begin{proof}
If $g=\sum_i a_if_i$, simultaneous vanishing of the $f_i$ gives vanishing of $g$. Conversely define $F:V\to\K^r$ by $F(v)=(f_i(v))$. Its kernel is the displayed intersection. The assumed containment lets us define $h:\im F\to\K$ by $h(Fv)=g(v)$: equal images differ by a vector in $\ker F$, hence have equal $g$ values. The resulting $h$ is linear. Extend an image basis to a basis of $\K^r$ and extend $h$ by zero on the added basis vectors. Every linear functional on $\K^r$ is $y\mapsto\sum_i a_i y_i$, by its values on the standard basis. Therefore $g(v)=h(Fv)=\sum_i a_i f_i(v)$ for all $v$. If $r=0$, the intersection is $V$, the span is zero, and both conditions say $g=0$. This is the precise reason that a homogeneous linear constraint is redundant exactly when its row functional lies in the span of the other rows.
\end{proof}

% END modules/M05.07.tex


% BEGIN modules/M05.08.tex
\chapter{Polynomials in operators}
Throughout this chapter, $V$ is a finite-dimensional vector space over the stated field $\K$, $n=\dim V$, and $T\in\End_\K(V)$. The identity of the zero space is its zero map. A scalar polynomial and its value at an operator are different objects.

\Kentry{08.evaluation}{Evaluation at an operator}
For $p(x)=\sum_{j=0}^d a_jx^j$, set $p(T)=\sum_{j=0}^d a_jT^j$, with $T^0=I$. This definition is unchanged by adjoining zero coefficients. Addition and scalar multiplication are preserved term by term. For multiplication, distributivity and $T^iT^j=T^{i+j}$ give
\[
 \Bigl(\sum_i a_iT^i\Bigr)\Bigl(\sum_j b_jT^j\Bigr)
 =\sum_k\Bigl(\sum_{i+j=k}a_ib_j\Bigr)T^k.
\]
Thus evaluation is a unital $\K$-algebra homomorphism into $\End_\K(V)$. In particular all polynomials in the same operator commute. This does not assert multiplicativity for arbitrary evaluations of noncommuting operators.

\Kentry{08.annihilating}{The annihilator ideal}
Define $\mathcal I_T=\{p\in\K[x]:p(T)=0\}$. It is an ideal: sums and negatives remain in the set, and $(qp)(T)=q(T)p(T)=0$. Since $\dim\End_\K(V)=n^2$, the $n^2+1$ operators $I,T,\ldots,T^{n^2}$ are linearly dependent. Their nontrivial relation gives a nonzero member of $\mathcal I_T$. If $n=0$, every polynomial annihilates $T$, and $\mathcal I_T=\K[x]$.

\Kentry{08.minpoly}{The minimal polynomial}
\begin{theorem}
There is a unique monic polynomial $m_T$ such that $\mathcal I_T=(m_T)$. It has least degree among nonzero annihilators, and divides every annihilator. For $V=0$ it is $1$.
\end{theorem}
\begin{proof}
Choose an annihilator of least degree and divide it by its nonzero leading coefficient to obtain $m$. For any $p\in\mathcal I_T$, Euclidean division gives $p=qm+r$ with $r=0$ or $\deg r<\deg m$. Evaluation gives $r(T)=0$, so minimality forces $r=0$. Conversely every multiple of $m$ annihilates $T$. Two monic generators divide one another, hence have equal degree and differ by a unit; monicity makes that unit $1$. In dimension zero the ideal contains $1$. In positive dimension a nonzero constant evaluates to a nonzero scalar multiple of $I$, so $\deg m_T\ge1$.
\end{proof}

\Kentry{08.charpoly}{The characteristic polynomial}
Choose a basis and let $A=[T]$. Define $\chi_T(x)=\det(xI-A)\in\K[x]$. If $A'=P^{-1}AP$, multiplicativity of the determinant over the commutative ring $\K[x]$ gives
\[
 \det(xI-A')=\det(P^{-1})\det(xI-A)\det(P)=\det(xI-A).
\]
Thus the definition is independent of basis. In the permutation formula, the only degree-$n$ term is the product of the $n$ diagonal $x$ terms, with coefficient $1$. Hence $\chi_T$ is monic of degree $n$. The empty determinant is $1$, giving $\chi_T=1$ on the zero space.

\Kentry{08.cayley_hamilton}{Cayley--Hamilton without a splitting assumption}
\begin{theorem}
$\chi_T(T)=0$ and consequently $m_T\mid\chi_T$.
\end{theorem}
\begin{proof}
For $n>0$, write $\chi_T(x)=\sum_{j=0}^n c_jx^j$ and $\operatorname{adj}(xI-A)=\sum_{j=0}^{n-1}B_jx^j$. The adjugate identity gives
\[
 (xI-A)\operatorname{adj}(xI-A)=\chi_T(x)I.
\]
Set $B_{-1}=B_n=0$. Comparison of coefficients gives $B_{j-1}-AB_j=c_jI$ for $0\le j\le n$. Multiply the $j$th equality on the left by $A^j$ and add. The two sums involving the $B_j$ cancel after shifting their index, so $\sum_{j=0}^n c_jA^j=0$. This argument does not substitute a matrix into a polynomial identity with noncommuting coefficients. In dimension zero, $\chi_T(T)=I=0$. Divisibility follows from the minimal-polynomial property.
\end{proof}

\Kentry{08.eigen_roots}{Scalar roots and eigenvalues}
\begin{theorem}
For $\lambda\in\K$, the following are equivalent: $\ker(T-\lambda I)\ne0$; $\chi_T(\lambda)=0$; $m_T(\lambda)=0$.
\end{theorem}
\begin{proof}
The first two conditions are equivalent because a square matrix has zero determinant exactly when it is not injective. If $Tv=\lambda v$ with $v\ne0$, induction on $j$ gives $T^jv=\lambda^jv$, whence $0=m_T(T)v=m_T(\lambda)v$ and $m_T(\lambda)=0$. Conversely, write $m_T=(x-\lambda)q$ if $m_T(\lambda)=0$. If $T-\lambda I$ were invertible, $(T-\lambda I)q(T)=0$ would imply $q(T)=0$, contradicting the minimal degree of $m_T$. Thus $T-\lambda I$ has a nonzero kernel. In dimension zero both polynomials are $1$ and all three conditions are false. No root outside $\K$ is being counted as a $\K$-eigenvalue.
\end{proof}

\Kentry{08.bezout_decomposition}{Coprime kernel decomposition}
\begin{theorem}
If $p,q\in\K[x]$ are coprime and $(pq)(T)=0$, then
\[V=\ker p(T)\oplus\ker q(T).\]
\end{theorem}
\begin{proof}
Choose $a,b$ with $ap+bq=1$. Every $v$ has the decomposition
\[v=b(T)q(T)v+a(T)p(T)v.\]
The first term is killed by $p(T)$ and the second by $q(T)$. If $v$ lies in both kernels, the same identity gives $v=0$. These formulas also supply the two projections. Their images and kernels show directly that they are idempotent and mutually annihilate one another. All maps involved commute with $T$, so both summands are invariant.
\end{proof}

\Kentry{08.primary}{Primary decomposition}
\begin{theorem}
If $m_T=\prod_{i=1}^s p_i^{a_i}$, where the $p_i$ are distinct monic irreducibles and $a_i\ge1$, then
\[V=\bigoplus_{i=1}^s V_i,\qquad V_i=\ker p_i(T)^{a_i}.\]
Each summand is invariant, and each projection onto a summand is a polynomial in $T$.
\end{theorem}
\begin{proof}
For $s=1$ the annihilating equation gives the claim. For $s>1$, $p_1^{a_1}$ is coprime to $q=\prod_{i>1}p_i^{a_i}$, so the previous theorem splits off $V_1$ and $\ker q(T)$. On the latter summand, apply the induction hypothesis to the pairwise coprime factors of $q$; it only requires that their product annihilate the restriction. The kernels obtained are precisely $V_i$, since a vector killed by $p_i(T)^{a_i}$ is also killed by $q(T)$. Each projection is a composition of the polynomial projections in the binary splits, hence polynomial in $T$. The empty factorization occurs only for $m_T=1$, namely $V=0$, and gives the empty direct sum.
\end{proof}

\Kentry{08.block_polynomials}{Invariant direct sums and their polynomials}
Suppose $V=\bigoplus_{i=1}^s V_i$, with all summands invariant, and let $T_i=T|_{V_i}$. Concatenating bases gives a block diagonal matrix; therefore
\[\chi_T=\prod_i\chi_{T_i}.\]
Furthermore $p(T)=0$ if and only if $p(T_i)=0$ for every $i$. This is equivalent to divisibility of $p$ by every $m_{T_i}$, hence by their monic least common multiple. Uniqueness of the monic generator gives
\[m_T=\operatorname{lcm}_i m_{T_i}.\]
Empty products and the least common multiple of an empty family are $1$.

\Kentry{08.polynomial_inverse}{Inverses represented by polynomials}
\begin{theorem}
$T$ is invertible if and only if $m_T(0)\ne0$. In that case $T^{-1}$ is a polynomial in $T$.
\end{theorem}
\begin{proof}
If $m_T(0)\ne0$, then $x$ and $m_T$ are coprime. Choose $a(x)x+b(x)m_T(x)=1$. Evaluation gives $a(T)T=Ta(T)=I$, proving invertibility and the formula. If $m_T(0)=0$, write $m_T=xq$; invertibility would imply $q(T)=0$, contradicting minimal degree. In dimension zero the condition is $1\ne0$ and the unique endomorphism is its own inverse.
\end{proof}

\Kentry{08.coefficients}{The trace and determinant coefficients}
For $n\ge1$, the coefficient of $x^{n-1}$ in $\det(xI-A)$ is $-\sum_i A_{ii}=-\tr A$. Indeed, a term of that degree chooses one matrix entry and $n-1$ diagonal $x$ entries. Any permutation fixing $n-1$ indices fixes the last one, so these terms all come from the identity permutation. Setting $x=0$ gives constant coefficient $\det(-A)=(-1)^n\det A$. All these equalities hold inside $\K$, including in positive characteristic. The degree-$n-1$ assertion is not made for $n=0$.

\paragraph{All coefficients of the adjugate.}
For $n\ge1$, write
\[\det(tI+A)=\sum_{p=0}^n e_p t^{n-p},\qquad e_0=1,
\qquad \operatorname{adj}(tI+A)=\sum_{p=0}^{n-1} Q_p t^{n-1-p}.
\]
Then
\[Q_p=\sum_{j=0}^p(-1)^j e_{p-j}A^j,
\qquad \tr Q_p=(n-p)e_p \quad(0\le p<n).
\]
Indeed, comparing coefficients in $(tI+A)\operatorname{adj}(tI+A)=\det(tI+A)I$ gives $Q_0=I$ and $Q_p+AQ_{p-1}=e_p I$ for $1\le p<n$. Induction on $p$ gives the displayed polynomial formula: insert the formula for $Q_{p-1}$ into $Q_p=e_p I-AQ_{p-1}$ and shift the summation index. Differentiating the determinant's permutation expansion with respect to $t$ differentiates one diagonal entry at a time; its cofactor is the corresponding diagonal adjugate entry. Thus
\[\frac{d}{dt}\det(tI+A)=\tr\operatorname{adj}(tI+A).
\]
Comparison of coefficients of $t^{n-p-1}$ proves the trace formula. In particular
\[\operatorname{adj}(A)=\sum_{j=0}^{n-1}(-1)^j e_{n-1-j}A^j.
\]
The $e_p$ here are the same elementary characteristic coefficients as in $\chi_A(x)=\sum_p(-1)^p e_p x^{n-p}$, because $\det(tI+A)=(-1)^n\chi_A(-t)$. This explicitly converts Greub's plus-sign expansion to our characteristic-polynomial convention. The argument uses no Newton recurrence and therefore creates no dependence on the later proof of Newton identities. It works in any characteristic; the integer $n-p$ is interpreted in the field, and is never divided out.

\Kentry{08.interpolation}{Evaluation and Lagrange interpolation}
Let $a_1,\ldots,a_r$ be distinct elements of $\K$. For $r>0$ put
\[L_i(x)=\prod_{j\ne i}\frac{x-a_j}{a_i-a_j}.\]
All denominators are nonzero; $\deg L_i<r$ and $L_i(a_j)=\delta_{ij}$. Thus the linear map from polynomials of degree less than $r$ to $\K^r$, $p\mapsto(p(a_1),\ldots,p(a_r))$, has right inverse $(c_i)\mapsto\sum_i c_iL_i$. It is injective because a nonzero polynomial of degree less than $r$ has fewer than $r$ distinct roots, by repeated division by linear factors. These maps are inverse linear isomorphisms. For $r=0$, interpret the domain as the zero polynomial space and the codomain as the zero coordinate space; the empty sum is the inverse map.

\Kentry{08.restriction_minpoly}{Restrictions and quotients}
Let $U$ be invariant under $T$, with restriction $T_U$ and induced quotient operator $\overline T(v+U)=Tv+U$. The latter is well defined because $T(U)\subset U$. Induction gives $T_U^j=T^j|_U$ and $\overline T^{\,j}(v+U)=T^jv+U$, so every annihilator of $T$ annihilates both induced operators. Consequently
\[m_{T_U}\mid m_T,\qquad m_{\overline T}\mid m_T.\]
A divisor of a split polynomial is split; a divisor of a product of distinct linear factors again has distinct linear factors. The criteria proved in \kref{10.triangularization} and \kref{09.minpoly_criterion} therefore show that triangularizability and diagonalizability pass to restriction and quotient. Those criteria are proved independently of this corollary; their proofs use only the preceding annihilator-divisibility argument.

\Kentry{08.newton_identities}{Newton identities without an analytic series}
\begin{theorem}
Write $\chi_T(x)=\sum_{k=0}^n(-1)^k e_kx^{n-k}$, where $e_0=1$. Over a characteristic-zero field, for $1\le k\le n$,
\[k e_k=\sum_{i=1}^k(-1)^{i-1}e_{k-i}\tr(T^i).\]
\end{theorem}
\begin{proof}
All series below are formal: each coefficient is a finite sum. Set $E(t)=\det(I-tA)=\sum_{k=0}^n(-1)^ke_kt^k$. The identity $E(t)=t^n\chi_T(t^{-1})$ follows by factoring $t^{-1}$ from each row in the Laurent-polynomial ring. The matrix series $G(t)=\sum_{j\ge0}t^jA^j$ satisfies $(I-tA)G=G(I-tA)=I$ by cancellation of consecutive coefficients. Differentiating the permutation formula for a determinant, one factor at a time, gives
\[E'(t)=-\tr\bigl(\operatorname{adj}(I-tA)A\bigr).
\]
Multiplication of the adjugate identity by $G$ gives $\operatorname{adj}(I-tA)=E(t)G(t)$. Therefore
\[E'(t)=-E(t)\sum_{i\ge1}\tr(A^i)t^{i-1}.\]
The coefficient of $t^{k-1}$ is
\[(-1)^k k e_k=-\sum_{i=1}^k(-1)^{k-i}e_{k-i}\tr(A^i),\]
which is the required formula. Characteristic zero ensures that division by $k$ recovers $e_k$ recursively. The identity before division makes sense in any characteristic, but a recursion dividing by a vanishing scalar does not. When $n=0$ there are no equations to check.
\end{proof}

% END modules/M05.08.tex


% BEGIN modules/M05.09.tex
\chapter{Eigenvalues and diagonalization}
Here $T$ is an endomorphism of a finite-dimensional $\K$-space $V$. Every assertion about diagonalization refers to the stated field $\K$.

\Kentry{09.eigenspace}{Eigenvalues and eigenspaces}
For $\lambda\in\K$, define $E_\lambda(T)=\ker(T-\lambda I)$. It is a subspace because it is a kernel. The scalar $\lambda$ is an eigenvalue precisely when this kernel contains a nonzero vector; an eigenvector is such a nonzero vector. On $V=0$ there are no eigenvalues. On a nonzero space the identity has just the eigenvalue $1$.

\Kentry{09.distinct_independence}{Independence of distinct eigenspaces}
\begin{theorem}
If $\lambda_1,\ldots,\lambda_s$ are distinct, then $\sum_i E_{\lambda_i}(T)$ is a direct sum.
\end{theorem}
\begin{proof}
Suppose $\sum_i v_i=0$ with $Tv_i=\lambda_i v_i$. Apply $\prod_{j\ne i}(T-\lambda_jI)$ to this equality. All terms except the $i$th vanish, leaving $\bigl(\prod_{j\ne i}(\lambda_i-\lambda_j)\bigr)v_i=0$. The scalar is a product of nonzero field elements, hence $v_i=0$. This holds for each $i$. In particular any list of nonzero eigenvectors for distinct eigenvalues is independent, so there are at most $\dim V$ distinct eigenvalues.
\end{proof}

\Kentry{09.diagonalizable}{Diagonalization as an eigenbasis}
An operator is diagonalizable over $\K$ if $V$ has a basis $(v_i)$ with $Tv_i=\lambda_i v_i$ and $\lambda_i\in\K$. The coordinate column of $Tv_i$ is then $\lambda_i$ times the $i$th coordinate vector, so the matrix is diagonal. Conversely a diagonal matrix in a basis gives these eigenvector equations for every basis vector. The empty basis diagonalizes the unique operator on the zero space.

\Kentry{09.eigenspace_criterion}{The eigenspace dimension criterion}
\begin{theorem}
$T$ is diagonalizable if and only if the sum of the dimensions of its distinct eigenspaces is $\dim V$.
\end{theorem}
\begin{proof}
There are at most $\dim V$ nonzero eigenspaces, and their sum is direct by the preceding theorem. Its dimension is therefore the sum of their dimensions. If this is $\dim V$, the sum equals $V$, and the union of bases of the summands is an eigenbasis. Conversely an eigenbasis is contained in this sum, so the sum is $V$ and the dimension equality follows. Empty sums give the zero-dimensional case.
\end{proof}

\Kentry{09.multiplicities}{Algebraic and geometric multiplicities}
For an eigenvalue $\lambda$, its algebraic multiplicity is the exponent of $x-\lambda$ in $\chi_T$; its geometric multiplicity is $r=\dim E_\lambda(T)$. Extend a basis of $E_\lambda(T)$ to a basis of $V$. In this basis the matrix of $T$ has the form
\[\begin{pmatrix}\lambda I_r&B\\0&C\end{pmatrix},\]
so $\chi_T(x)=(x-\lambda)^r\det(xI-C)$. Thus the algebraic multiplicity is at least $r$, and $r\ge1$ by definition of eigenvalue. This proof does not require that the remaining factor split.

\Kentry{09.minpoly_criterion}{The minimal-polynomial criterion}
\begin{theorem}
$T$ is diagonalizable over $\K$ if and only if $m_T$ is a product of distinct linear factors over $\K$, including the empty product $1$.
\end{theorem}
\begin{proof}
In an eigenbasis, $p(T)$ is diagonal with diagonal entries $p(\lambda_i)$. It vanishes exactly when $p$ is divisible by each of the distinct factors $x-\lambda_i$, so the monic minimal polynomial is their product. Conversely, if $m_T=\prod_{i=1}^s(x-\lambda_i)$ with distinct $\lambda_i$, the coprime-kernel decomposition of \kref{08.primary} gives $V=\bigoplus_i\ker(T-\lambda_iI)$. Choose a basis in each summand. Their union is an eigenbasis. If $m_T=1$, then $I=0$ and $V=0$.
\end{proof}

\Kentry{09.distinct_roots}{Simple spectrum suffices}
If $\chi_T$ has $n=\dim V$ distinct roots in $\K$, each is an eigenvalue by \kref{08.eigen_roots}. Choose one nonzero eigenvector for each root. The resulting $n$ vectors are independent, hence a basis, and $T$ is diagonalizable. The condition is sufficient but not necessary: the identity on a space of dimension at least two has a repeated root in its characteristic polynomial and is diagonalizable.

\Kentry{09.commutant_invariance}{The subspaces preserved by commuting maps}
If $ST=TS$, induction gives $ST^j=T^jS$ for all $j\ge0$, and hence $Sp(T)=p(T)S$. If $p(T)v=0$, then $p(T)Sv=Sp(T)v=0$. Thus $S$ preserves every kernel of a polynomial in $T$. Taking $p=x-\lambda$ gives preservation of eigenspaces, and taking a power of this polynomial gives preservation of generalized eigenspaces.
It also preserves every polynomial image: if $w=p(T)v$, then $Sw=Sp(T)v=p(T)Sv\in\im p(T)$. Setting $S=T$ gives invariance of both $\ker p(T)$ and $\im p(T)$ under $T$ itself.

\Kentry{09.simultaneous_diagonal}{Simultaneous diagonalization}
\begin{theorem}
A possibly infinite family of pairwise commuting operators, each diagonalizable over $\K$, has a common eigenbasis.
\end{theorem}
\begin{proof}
First treat a finite family by induction on its number $s$ of operators. With no operators any basis works, and with one operator this is the hypothesis. Decompose $V$ into the eigenspaces of the first operator. Commutation makes each summand invariant under the other operators. Each of their restrictions is annihilated by the original minimal polynomial, a product of distinct linear factors. Its own minimal polynomial divides that polynomial by direct restriction of the annihilating identity. The criterion just proved makes each restriction diagonalizable. The restrictions still commute, so the induction hypothesis supplies a common eigenbasis in each summand for the last $s-1$ operators. Concatenating these bases gives a basis for the whole finite family, because the first operator acts as a scalar on each summand.

For an arbitrary family $\mathcal F$, its linear span in the finite-dimensional space $\End_\K(V)$ is spanned by a finite subfamily $T_1,\ldots,T_r$ of $\mathcal F$: repeatedly add a member outside the current span, and dimension bounds the number of additions by $n^2$. Apply the finite result to this subfamily. Every member of $\mathcal F$ is a linear combination of these diagonal operators, so its matrix in the same basis is diagonal. This argument does not assume, before constructing the common basis, that all linear combinations are diagonalizable. The empty family and the zero space are included.
\end{proof}
Conversely, a family with a common eigenbasis consists of diagonalizable operators that commute pairwise: diagonal matrix products multiply corresponding scalar entries, and scalar multiplication is commutative. Thus, for a family of individually diagonalizable operators, pairwise commutation is equivalent to simultaneous diagonalizability.

\Kentry{09.field_dependence}{Field and multiplicity counterexamples}
For $R=\begin{pmatrix}0&-1\\1&0\end{pmatrix}$, $\chi_R=x^2+1$ has no real roots. It therefore has no real eigenbasis. Over $\C$, the vectors $(1,-i)^t$ and $(1,i)^t$ have eigenvalues $i$ and $-i$; they are independent, and diagonalize $R$. For $J=\begin{pmatrix}1&1\\0&1\end{pmatrix}$ the only eigenspace is $\ker(J-I)=\Span\{(1,0)^t\}$, of dimension one, so $J$ is not diagonalizable even over $\C$. In contrast $I_2$ has the same characteristic polynomial $(x-1)^2$ and is already diagonal. Thus repeated roots alone do not decide diagonalizability.

\Kentry{09.existence}{Existence of eigenvalues}
On a nonzero finite-dimensional complex space, $\chi_T$ has positive degree, so the fundamental theorem of algebra supplies a root. By \kref{08.eigen_roots} it is an eigenvalue. For a real space of odd dimension, $\chi_T$ is a real monic polynomial of odd degree. Its leading term dominates the sum of the lower terms for sufficiently large absolute value of the argument, so its values have opposite signs at a sufficiently large positive and negative argument. Continuity and the intermediate value theorem give a real root. In even positive dimension, take a direct sum of the real rotation matrices in the previous entry; its characteristic polynomial $(x^2+1)^{n/2}$ has no real root. These examples establish the precise boundary of the universal existence statement.

\paragraph{The degree-divisibility obstruction.}
If $p\in\K[x]$ is irreducible of degree $d>0$, then
\[d\mid\dim_\K\ker p(T).
\]
To prove this, let $W=\ker p(T)$. Since $W$ is preserved by every polynomial in $T$, the field $F=\K[x]/(p)$ acts on $W$ by $[q]\cdot w=q(T)w$. The action is well defined: changing $q$ by a multiple of $p$ changes its value by an operator zero on $W$. The field and module identities follow from polynomial evaluation; a nonzero class has an inverse by Bezout. A finite $\K$-spanning list also spans over $F$, so choose an $F$-basis $w_1,\ldots,w_r$. The vectors $T^jw_i$, $1\le i\le r$ and $0\le j<d$, form a $\K$-basis of $W$. They span by division modulo $p$. For independence, group a proposed relation by $i$ into $\sum_i[q_i]\cdot w_i=0$, with $\deg q_i<d$; $F$-independence forces $[q_i]=0$, and division then forces every $q_i=0$. Thus $\dim_\K W=dr$, including $r=0$.

In particular, over $\R$, if $b^2<4c$, the polynomial $x^2+bx+c=(x+b/2)^2+c-b^2/4$ has no real root and is irreducible. Hence $\dim_\R\ker(T^2+bT+cI)$ is even. This is a statement about a particular kernel, not only about the existence of an eigenvalue on an odd-dimensional whole space.

\Kentry{09.gershgorin}{Gershgorin disks}
\begin{theorem}
Every eigenvalue $\lambda$ of a complex $n\times n$ matrix $A$ lies in at least one disk
\[\abs{\lambda-A_{ii}}\le\sum_{j\ne i}\abs{A_{ij}}.\]
The analogous disks with column sums also contain the spectrum.
\end{theorem}
\begin{proof}
Take an eigenvector $v\ne0$ and choose $i$ with $\abs{v_i}=\max_j\abs{v_j}>0$. The $i$th row of $Av=\lambda v$ gives
\[\abs{\lambda-A_{ii}}\abs{v_i}
 =\Bigl|\sum_{j\ne i}A_{ij}v_j\Bigr|
 \le\sum_{j\ne i}\abs{A_{ij}}\abs{v_j}
 \le\Bigl(\sum_{j\ne i}\abs{A_{ij}}\Bigr)\abs{v_i}.
\]
Divide by $\abs{v_i}$. For the column statement, $\det(xI-A^t)=\det((xI-A)^t)=\det(xI-A)$, so $A^t$ has the same eigenvalues. Apply the row statement to $A^t$. For $n=0$ the spectrum is empty. No claim about counting eigenvalues in unions of disks is part of this result.
\end{proof}

\Kentry{09.diagonal_dominance}{Strict diagonal dominance}
If $\abs{A_{ii}}>\sum_{j\ne i}\abs{A_{ij}}$ for every row, zero lies in none of the row disks. Thus zero is not an eigenvalue and $A$ is invertible. For a real matrix apply the same proof over $\C$; a nonzero real determinant is equivalent to invertibility over $\R$. Column dominance follows using the column disks. The strict inequality is essential to this argument: the all-ones $2\times2$ matrix satisfies the non-strict row inequalities and is singular. The empty matrix is invertible and the condition is vacuous.

\Kentry{09.commuting_spectrum}{Spectra of sums and products}
For commuting complex operators $S,T$, simultaneous triangularization, proved independently in \kref{10.simultaneous_triangular}, gives a basis in which both are upper triangular. Write the paired diagonal entries as $(a_i,b_i)$. Matrix addition has diagonal entries $a_i+b_i$, and the diagonal of the product is $a_ib_i$, because all other summands in its diagonal entries contain a below-diagonal factor. The characteristic polynomial of an upper triangular matrix is the product of $x$ minus its diagonal entries. Hence, with algebraic multiplicity, the eigenvalues of $S+T$ and $ST$ are respectively $a_i+b_i$ and $a_ib_i$. Pairing is with respect to a common triangular basis: arbitrary choices of one eigenvalue of $S$ and one of $T$ need not occur together. For example $S=T=\diag(0,1)$ has no eigenvalue $1$ in $S+T$, even though $0+1=1$.

\paragraph{Polynomial spectral mapping.}
If $\chi_T$ splits over $\K$ and $p\in\K[x]$, then
\[\operatorname{spec}_\K p(T)=\{p(\lambda):\lambda\in\operatorname{spec}_\K T\}.
\]
More precisely, if the diagonal entries in a triangular matrix for $T$ are $\lambda_1,\ldots,\lambda_n$, then
\[\chi_{p(T)}(x)=\prod_{j=1}^n(x-p(\lambda_j)).\]
Indeed, the diagonal of a product of upper triangular matrices is the product of their diagonals. Induction makes the diagonal entries of $T^k$ equal to $\lambda_j^k$, and linearity makes those of $p(T)$ equal to $p(\lambda_j)$. The triangular determinant formula proves the assertion, including algebraic multiplicities and the empty product at $n=0$. Over $\C$ splitting is automatic. Without splitting only the inclusion from right to left always follows: $Tv=\lambda v$ implies $p(T)v=p(\lambda)v$. Equality can fail over $\R$: the rotation $R=\begin{pmatrix}0&-1\\1&0\end{pmatrix}$ has empty real spectrum, whereas for $p(x)=x^2$ one has $p(R)=-I$ and real spectrum $\{-1\}$.

\Kentry{09.real_complex_bridge}{Realification, real forms, and complex orientation}
These are two different operations. Realification forgets complex scalar multiplication while retaining the underlying vectors. Complexification starts with a real space and adds an imaginary copy.

\paragraph{Complexification and real forms.}
For a real vector space $W$, form $W_\C=W\oplus W$ and write $(x,y)$ as $x+iy$. Set
\[(a+ib)(x+iy)=(ax-by)+i(bx+ay).
\]
Addition is componentwise. Distributing the two sides of scalar associativity gives the same real and imaginary coefficients using $(a+ib)(c+id)=(ac-bd)+i(ad+bc)$; distributivity and the identity scalar follow by the same two-coordinate calculation. Thus this is a complex vector space. If $U$ is another real vector space, a real-linear map $L:W\to U$ has the unique complex-linear extension $L_\C:W_\C\to U_\C$, given by $L_\C(x+iy)=Lx+iLy$: the displayed scalar formula proves linearity, and any extension must have these values by complex linearity. If $(w_j)$ is a real basis of $W$, it is a complex basis of $W_\C$. For spanning, expand $x,y$ in this basis; for independence, separate the real and imaginary parts of any complex relation.

More intrinsically, a conjugation on a complex space $V$ is an additive map $J$ with $J(zv)=\overline zJv$ and $J^2=I$. Its fixed space $W=\{v:Jv=v\}$ is a real subspace. Every $v\in V$ has the decomposition
\[v=x+iy,\qquad x=\frac{v+Jv}{2},\qquad y=\frac{v-Jv}{2i},\qquad x,y\in W.
\]
Applying $J$ verifies that $x,y$ are fixed. If $x+iy=0$ with $x,y\in W$, conjugation also gives $x-iy=0$, so addition and subtraction give $x=y=0$. Hence $V=W\oplus_\R iW$, and the natural map $W_\C\to V$ is an isomorphism. Such a $W$ is called a real form. Conversely a real subspace with this unique decomposition defines a conjugation by $J(x+iy)=x-iy$; the scalar formula above proves conjugate linearity and $J^2=I$. It is the only conjugation fixing $W$ pointwise, since those values determine its values on $x+iy$. When $W$ is finite dimensional, the real-form condition is equivalent to any real basis of $W$ being a complex basis of $V$: expansion proves one direction, and separating complex coefficients proves the other. Consequently $\dim_\R W=\dim_\C V$. A complex-linear operator preserves $W$ if and only if it commutes with $J$: apply the commutation equation to fixed vectors for one direction; for the other write $v=x+iy$ and use preservation on $x,y$.

If $b$ is a real positive definite inner product on $W$, there is a unique complex inner product extending it, in our first-slot-linear convention:
\[h(x+iy,u+iv)=b(x,u)+b(y,v)+i\bigl(b(y,u)-b(x,v)\bigr).
\]
Expansion verifies additivity, linearity in the first variable, conjugate linearity in the second, and conjugate symmetry. On the diagonal it equals $b(x,x)+b(y,y)$, which is positive unless $x=y=0$. Any extension must have the displayed expansion by its two scalar-linearity rules, proving uniqueness. It also gives $h(Jz,Jw)=\overline{h(z,w)}=h(w,z)$, so $J$ is antiunitary. This construction uses only the inner-product axioms, not a spectral theorem.

\paragraph{Realification of a complex operator.}
If $V$ has complex basis $e_1,\ldots,e_n$, its realification $V_\R$ has real basis $e_1,\ldots,e_n,ie_1,\ldots,ie_n$, by separating complex coordinates; hence its real dimension is $2n$. For the complex matrix $A=B+iC$ of $T$, the matrix of the same map on $V_\R$ is
\[A_\R=\begin{pmatrix}B&-C\\ C&B\end{pmatrix}.
\]
Over $\C$, the invertible coordinate transformation $(x,y)\mapsto(x+iy,x-iy)$ intertwines $A_\R$ with $\diag(A,\overline A)$. Its inverse is $(u,v)\mapsto((u+v)/2,(u-v)/(2i))$, so no unproved change-of-basis assertion is hidden here. Therefore
\[\chi_{T_\R}(t)=\chi_T(t)\overline{\chi_T}(t),\qquad
\det_\R T_\R=\abs{\det_\C T}^2,\qquad
\tr_\R T_\R=2\operatorname{Re}\tr_\C T.
\]
The bar on a polynomial conjugates its coefficients and fixes its formal variable. The product is fixed by conjugation and is a real polynomial. Both determinant and characteristic identities include singular operators and the zero space.

Declare the interleaved real basis $(e_1,ie_1,\ldots,e_n,ie_n)$ positive. A different complex basis is related by some invertible complex matrix $P$; the real determinant of its coordinate change is $\abs{\det_\C P}^2>0$. Passing from grouped to interleaved coordinates conjugates every transition matrix by the same permutation, leaving its determinant unchanged. Thus this orientation is independent of the chosen complex basis: it is our canonical complex orientation. The grouping permutation itself has sign $(-1)^{n(n-1)/2}$, so grouped and interleaved volume generators must not be identified without this sign. Greub Section 7.9 declares the grouped basis positive; this is another basis-independent convention, related to ours by precisely that sign. The determinant and characteristic-polynomial identities above do not depend on this convention. The empty ordered basis gives the zero-dimensional orientation.

% END modules/M05.09.tex


% BEGIN modules/M05.10.tex
\chapter{Invariant decompositions and canonical forms}
Let $V$ be an $n$-dimensional vector space over a field $\K$, and let $T\in\End_\K(V)$. Splitting is always over the field explicitly named. The proofs of rational and Jordan form below include the polynomial-matrix reduction that produces the decompositions; classification is not an external assumption.

\Kentry{10.triangularization}{The triangularization criterion}
\begin{theorem}
The following are equivalent: $T$ has an upper triangular matrix in a $\K$-basis; $\chi_T$ splits over $\K$; $m_T$ splits over $\K$.
\end{theorem}
\begin{proof}
An upper triangular matrix has characteristic polynomial equal to the product of $x$ minus its diagonal entries, so triangularity implies splitting of $\chi_T$. Since $m_T\mid\chi_T$, splitting of $\chi_T$ implies splitting of $m_T$.

For the remaining implication induct on $n$. At $n=0$ the empty basis works. If $n>0$, $m_T$ has positive degree, so its splitting supplies a root $\lambda\in\K$. By \kref{08.eigen_roots} there is a nonzero $v_1$ with $Tv_1=\lambda v_1$. The line $L=\K v_1$ is invariant. The quotient operator on $V/L$ is annihilated by $m_T$, so its minimal polynomial divides $m_T$ and splits. By induction choose a triangularizing basis of $V/L$, and lift its vectors to $v_2,\ldots,v_n$. The vectors $v_1,\ldots,v_n$ are a basis. For each $j$, the quotient statement says that $Tv_j$ differs from a combination of $v_2,\ldots,v_j$ by a vector in $L$; hence $Tv_j\in\Span(v_1,\ldots,v_j)$. This is exactly upper triangularity.
\end{proof}

\Kentry{10.simultaneous_triangular}{Simultaneous triangularization}
\begin{theorem}
A possibly infinite pairwise commuting family of triangularizable operators on $V$ has a common upper triangular matrix basis.
\end{theorem}
\begin{proof}
First treat a finite family. On a nonzero space it has a common eigenvector. To see this, start with the whole space and consider the operators successively. At each step the space already chosen is a nonzero common eigenspace for the preceding operators and is invariant under every remaining operator, by commutation. The next restriction is annihilated by the original split minimal polynomial, so its own minimal polynomial splits; it therefore has an eigenvalue and a nonzero eigenspace inside the chosen space. Replace the chosen space by that eigenspace. After finitely many steps choose a nonzero vector in the final space. For an empty family choose any nonzero vector.

Now induct on $\dim V$. The zero case is immediate. A common eigenvector spans a common invariant line $L$. The quotient operators commute and each is annihilated by its original split minimal polynomial, so they remain triangularizable by the preceding criterion. The induction hypothesis gives a common triangular basis in $V/L$. Lift it and prepend the common eigenvector. For every operator, the same quotient calculation used in \kref{10.triangularization} shows that it preserves every initial span of the lifted basis. Thus all its matrices are upper triangular in that basis.

For an arbitrary family, choose a finite subfamily spanning its linear span in $\End_\K(V)$, as in \kref{09.simultaneous_diagonal}. The finite result triangularizes that subfamily. Every other member has an upper triangular matrix because it is a linear combination of those matrices. No field closure beyond the stated triangularizability hypotheses is required.
\end{proof}

\Kentry{10.kernel_stabilization}{Stabilizing kernels and images}
\begin{theorem}
For $k\ge0$, $\ker T^k\subseteq\ker T^{k+1}$ and $\im T^{k+1}\subseteq\im T^k$. If two consecutive kernels are equal, all later kernels are equal. Both sequences stabilize by exponent $n$.
\end{theorem}
\begin{proof}
The inclusions follow by applying one extra copy of $T$. Suppose $\ker T^k=\ker T^{k+1}$ and $T^{k+2}v=0$. Then $Tv\in\ker T^{k+1}=\ker T^k$, so $v\in\ker T^{k+1}$. Thus $\ker T^{k+2}=\ker T^{k+1}$; repetition proves permanent equality. Until stabilization, each inclusion is strict and increases dimension by at least one. Starting from $\ker T^0=0$, there can be at most $n$ increases, so $\ker T^n=\ker T^{n+j}$ for all $j\ge0$. Rank-nullity makes the dimensions of the corresponding images equal, and their inclusions make the images themselves equal. When $n=0$, both statements already hold at exponent zero.
\end{proof}

\Kentry{10.nilpotent}{Nilpotent operators and their index}
An operator $N$ is nilpotent if $N^r=0$ for some positive integer $r$. On a nonzero finite-dimensional space its index is the least such positive integer. If $\dim V=n>0$ and $N$ is nilpotent, the stabilized kernel $\ker N^n$ must be $V$: it equals every later kernel, including one known to equal $V$. Thus $N^n=0$ and the index is at most $n$. On the zero space every power, including $N^0=I$, is zero, but the least \emph{positive} index is $1$; it must not be asserted to be at most zero. The only eigenvalue of a nilpotent operator on a nonzero space is $0$: $Nv=\lambda v$ implies $0=N^rv=\lambda^rv$, and a field has no nonzero nilpotent scalar.

\Kentry{10.fitting}{The Fitting decomposition}
\begin{theorem}
For every $T$,
\[V=\ker T^n\oplus\im T^n.\]
Both summands are invariant. On the first, $T$ is nilpotent; on the second, it is invertible.
\end{theorem}
\begin{proof}
If $v=T^nu\in\ker T^n$, then $T^{2n}u=0$. Stabilization gives $u\in\ker T^n$, so $v=0$. Thus the two subspaces intersect trivially. Rank-nullity says their dimensions sum to $n$, proving that they sum to $V$. Powers of $T$ commute with $T$, which gives invariance. On the kernel, the $n$th power vanishes if $n>0$; the zero-space case is nilpotent as just explained. On the image, $T(\im T^n)=\im T^{n+1}=\im T^n$, so the restriction is surjective and therefore invertible in finite dimension. For $n=0$, both summands are zero and the statements hold for the unique map.
\end{proof}

\Kentry{10.generalized}{Generalized eigenspaces}
For $\lambda\in\K$, put $N_\lambda=T-\lambda I$ and define
\[G_\lambda(T)=\ker N_\lambda^n.\]
Stabilization applied to $N_\lambda$ proves that this kernel equals the union of the kernels of all its nonnegative powers and is unchanged by increasing the exponent beyond $n$. A generalized eigenvector is a \emph{nonzero} vector killed by some positive power of $N_\lambda$. The space $G_\lambda$ is invariant because $T$ commutes with $N_\lambda$. On $G_\lambda$, $T=\lambda I+N_\lambda$ with a nilpotent restriction.

\Kentry{10.generalized_decomposition}{Generalized eigenspace decomposition}
\begin{theorem}
If $\chi_T$ splits over $\K$, then
\[V=\bigoplus_{\lambda\in\operatorname{spec}_\K(T)}G_\lambda(T).\]
On each summand the operator is $\lambda I$ plus a nilpotent operator.
\end{theorem}
\begin{proof}
The minimal polynomial splits and can be written $m_T=\prod_i(x-\lambda_i)^{a_i}$ with distinct $\lambda_i$ and $a_i\ge1$. Primary decomposition gives summands $W_i=\ker(T-\lambda_iI)^{a_i}$. They are contained in $G_{\lambda_i}$ because $a_i\le\deg m_T\le n$. On $W_j$ with $j\ne i$, $T-\lambda_iI$ is invertible: a Bezout identity between $x-\lambda_i$ and $(x-\lambda_j)^{a_j}$ evaluates to an inverse. Therefore a vector in $\bigoplus_jW_j$ killed by a power of $T-\lambda_iI$ has zero components in all $W_j$, $j\ne i$. This proves $G_{\lambda_i}=W_i$. A scalar not among the $\lambda_i$ has $G_\lambda=0$ by the same argument. The nilpotent assertion follows from the defining kernel equation. The zero space has the empty decomposition.
\end{proof}

\Kentry{10.cyclic}{Cyclic subspaces and vector annihilators}
For $v\in V$, let $C_T(v)=\Span\{T^jv:j\ge0\}$. It is invariant. The map
\[\K[x]\longrightarrow C_T(v),\qquad p\longmapsto p(T)v\]
is surjective and linear. Its kernel is an ideal, nonzero since $m_T$ belongs to it. By the same minimal-degree division proof as \kref{08.minpoly}, it has a unique monic generator $d_v$. Hence it induces a $\K[x]$-module isomorphism $\K[x]/(d_v)\simeq C_T(v)$, where $x$ acts as $T$. If $r=\deg d_v$, division gives representatives of degree less than $r$, and such a representative belongs to $(d_v)$ only if zero. Thus $v,Tv,\ldots,T^{r-1}v$ are a basis of the cyclic space and its dimension is $r$. For $v=0$, $d_v=1$ and the list is empty.

\Kentry{10.companion}{Companion matrices}
Let $d=x^r+c_{r-1}x^{r-1}+\cdots+c_0$ be monic of positive degree. Multiplication by $x$ on $\K[x]/(d)$ has, in the basis $1,x,\ldots,x^{r-1}$, the companion matrix $C(d)$ specified by
\[
C(d)_{i+1,i}=1\quad(1\le i<r),\qquad
C(d)_{i,r}=-c_{i-1}\quad(1\le i\le r),
\]
with all other entries zero. Thus $C(x+c_0)=[-c_0]$, and the same specification handles every positive size without a special display convention. The vector $1$ is cyclic with annihilator $d$. A polynomial annihilates the entire multiplication operator exactly when it annihilates $1$: one direction is immediate, and the other follows from $p(x)q(x)=q(x)p(x)$ in the quotient. Thus its minimal polynomial is $d$. Cayley--Hamilton says $d$ divides its monic characteristic polynomial of degree $r$; equal degree and monicity give equality.

\Kentry{10.invariant_factor_computation}{Smith reduction over the polynomial ring}
\begin{theorem}\label{aux:smith-polynomial}
A matrix over $R=\K[x]$ can be transformed by invertible elementary row and column operations into a diagonal matrix with nonzero monic entries $d_1,\ldots,d_r$ satisfying $d_1\mid\cdots\mid d_r$, followed by zeros. For each $k\le r$, the monic gcd $\Delta_k$ of its $k\times k$ minors equals $d_1\cdots d_k$; put $\Delta_0=1$. Thus the nonzero diagonal entries are uniquely determined by $d_k=\Delta_k/\Delta_{k-1}$.
\end{theorem}
\begin{proof}
If the matrix is zero, stop. Otherwise move a nonzero entry to position $(1,1)$ by row and column swaps, and call it $a$. Divide every other entry in the first column by $a$ using Euclidean division. Subtract its quotient times row $1$ from that row. If a nonzero remainder occurs, swap it to position $(1,1)$ and restart. The new pivot has strictly smaller degree. When the first column can be cleared, perform the same procedure on the first row by column operations. A nonzero remainder again gives a smaller-degree pivot. Because nonnegative degrees cannot decrease indefinitely, eventually all entries in both the first column and first row outside the pivot are zero.

At that stage write the matrix as $\diag(a,B)$. If an entry $b$ of $B$ is not divisible by $a$, add its row to row $1$. The first entry in row $1$ remains $a$, while another entry in that row becomes $b$. Division by $a$ and a column operation produce a nonzero remainder of degree less than $\deg a$; move it to the pivot and repeat the clearing procedure. This further strict decrease can occur only finitely often. Thus eventually the pivot divides every entry of $B$. Multiply row $1$ by the inverse of its leading coefficient to make the pivot monic.

Apply the same algorithm inductively to $B$. All its entries are multiples of $a$, and row and column operations confined to $B$ preserve this property. Hence its first nonzero final diagonal entry is divisible by $a$. The induction hypothesis orders the remaining entries by divisibility. This proves existence and termination. Every operation used is invertible over $R$: swaps, adding a polynomial multiple of a different row or column, and multiplication by a nonzero scalar. No division by an arbitrary nonunit of $R$ is used as a matrix operation.

For uniqueness, consider the ideal generated by the $k\times k$ minors. A row addition changes any such minor into the old minor plus a polynomial multiple of a minor using the replacement row, or leaves it unchanged; if the replacement row is already among the selected rows the extra determinant vanishes. Row swaps and unit rescalings also preserve the generated ideal. The same reasoning applies to columns, and inverse operations give equality of ideals in both directions. For a diagonal matrix with the divisibility chain, every nonzero $k$-minor is a product $d_{i_1}\cdots d_{i_k}$ with $i_j\ge j$, so it is divisible by $d_1\cdots d_k$. The leading $k$-minor equals this product. Consequently the ideal is generated by that product, whose monic generator is $\Delta_k$. This proves the formula and uniqueness.
\end{proof}
For $xI-A$, the determinant is the nonzero monic polynomial $\chi_A$, so Smith reduction has $n$ nonzero diagonal entries. Discard the unit entries $1$; the remaining entries will be exactly the invariant factors of the operator, as the next proof establishes. The gcd-of-minors formulas give an effective computation, not just an existence assertion.

\Kentry{10.rational}{Cyclic decomposition and rational canonical form}
\begin{theorem}
Every finite-dimensional operator is a direct sum of cyclic operators with monic nonconstant annihilators $d_1\mid\cdots\mid d_r$. It therefore has a matrix $\diag(C(d_1),\ldots,C(d_r))$ over its original field.
\end{theorem}
\begin{proof}
In a chosen basis let $A=[T]$ and make $V$ an $R=\K[x]$-module by letting $xv=Tv$. Define
\[\Phi:R^n\longrightarrow V,\qquad \sum_{j\ge0}x^jv_j\longmapsto\sum_{j\ge0}A^jv_j,
\]
where each $v_j\in\K^n$ is identified with a vector of $V$. This map is onto because constant vectors are allowed. Its kernel is exactly $(xI-A)R^n$. One inclusion follows from evaluation. For the other, for every $j\ge1$ and constant vector $v$,
\[x^jv-A^jv=(xI-A)\sum_{\ell=0}^{j-1}x^{j-1-\ell}A^\ell v.\]
Adding these identities shows that any vector polynomial differs from the constant vector $\Phi(f)$ by an element of $(xI-A)R^n$. Thus a vector polynomial in the kernel belongs to that image. The induced map is an $R$-module isomorphism
\[R^n/(xI-A)R^n\simeq V.\]
Smith reduction gives invertible $U,W$ with $U(xI-A)W=\diag(e_1,\ldots,e_n)$. Right multiplication by $W$ does not change the image submodule, and left multiplication by $U$ induces an isomorphism on quotients. Hence
\[V\simeq\bigoplus_{i=1}^n R/(e_i).
\]
The unit entries give zero summands; discard them. For each remaining $d_i=e_j$, the quotient has basis $1,x,\ldots,x^{\deg d_i-1}$ and multiplication by $x$ is $C(d_i)$. Transport these bases to $V$ and concatenate them. This yields the asserted matrix and the divisibility chain. If $V=0$, the matrix and list are empty.
\end{proof}

\paragraph{Relative cyclic decomposition and admissible subspaces.}
An invariant subspace $W\subseteq V$ is \emph{$T$-admissible} if, for every $f\in\K[x]$ and $v\in V$ with $f(T)v\in W$, there exists $w\in W$ satisfying $f(T)v=f(T)w$. This condition means that a polynomial equation whose right side belongs to $W$, if solvable in $V$, is already solvable in $W$.
\begin{theorem}
An invariant subspace $W$ is $T$-admissible if and only if it has a $T$-invariant complement. More precisely, if $W$ is admissible, let $d_1\mid\cdots\mid d_r$ be the nonconstant invariant factors of the induced operator on $V/W$. There are vectors $u_1,\ldots,u_r\in V$ with vector annihilators $d_1,\ldots,d_r$ such that
\[V=W\oplus C_T(u_1)\oplus\cdots\oplus C_T(u_r).
\]
\end{theorem}
\begin{proof}
Suppose first that $V=W\oplus L$ with both summands invariant. Write $v=w+\ell$. If $f(T)v\in W$, then $f(T)\ell=f(T)v-f(T)w$ belongs to $W$ and also to $L$, by invariance. Hence it is zero, proving admissibility.

Conversely, apply the already proved rational decomposition to the quotient $V/W$. Choose cyclic generators $\overline v_i$ with annihilators $d_i$ and choose lifts $v_i\in V$. Then $d_i(T)v_i\in W$, so admissibility supplies $w_i\in W$ with $d_i(T)v_i=d_i(T)w_i$. Put $u_i=v_i-w_i$. Its quotient class remains $\overline v_i$ and $d_i(T)u_i=0$. Define an $\K[x]$-linear map
\[F:\bigoplus_i\K[x]/(d_i)\longrightarrow V,\qquad
([q_i])_i\longmapsto\sum_i q_i(T)u_i.
\]
It is well defined by the annihilating equations. Composing with the quotient map $\pi:V\to V/W$ gives the rational-decomposition isomorphism sending the $i$th class of $1$ to $\overline v_i$. Thus $F$ is injective, and its image $L$ intersects $W$ trivially: if $F(z)\in W$, then $(\pi F)(z)=0$ and $z=0$. Surjectivity of $\pi F$ says that for every $v$ there is $\ell\in L$ with $v-\ell\in W$. Hence $V=W\oplus L$. The map $F$ also proves that $L$ is the direct sum of its cyclic images and that each $u_i$ has exactly annihilator $d_i$, not a proper divisor. All summands are invariant. If $W=V$, the quotient and the list of factors are empty, giving complement zero. This proof lifts the existing quotient rational factors; it does not assume a complement in order to construct them.
\end{proof}
The uniqueness of the ordered factor list follows from the quotient operator's invariant-factor uniqueness proved next. Hoffman--Kunze orders the relative factors in the reverse divisibility order; reversing the list gives our $d_1\mid\cdots\mid d_r$. Neither the complement nor the individual cyclic generators is asserted to be unique.

\Kentry{10.invariant_factors}{Uniqueness and the similarity criterion}
\begin{theorem}
The nonconstant polynomials $d_1\mid\cdots\mid d_r$ in rational form are uniquely determined by $T$. Their product is $\chi_T$ and, if $V\ne0$, their largest member is $m_T$. Two matrices over $\K$ are similar if and only if their invariant-factor lists agree.
\end{theorem}
\begin{proof}
If $A'=P^{-1}AP$, then $xI-A'=P^{-1}(xI-A)P$. Invertible constant matrices act by invertible row and column operations over $\K[x]$: Gaussian elimination over $\K$ expresses each as a product of elementary matrices. Thus the determinantal ideals, and hence the Smith entries, are unchanged. This proves independence of basis and uniqueness of the list produced above. Conversely, suppose an operator is already a direct sum of cyclic quotients with a divisibility chain. A presentation for the sum is the diagonal matrix of their annihilators, and its $k$th determinantal ideal is generated by the product of its first $k$ entries. Equivalently, one may use the basis independence just proved and the companion-block representative of that same sum. The presentation reduction in the preceding proof shows that the nonunit Smith entries coincide with this chain. For completeness, the presentation of a single companion block reduces to $\diag(1,\ldots,1,d)$: eliminate successively the relations $xe_i-e_{i+1}$, leaving only $d e_1=0$. For several blocks these eliminations give their diagonal annihilators, whose chain is already in Smith form.

Characteristic polynomials multiply under invariant direct sums, and each companion block has characteristic polynomial $d_i$. Thus $\chi_T=\prod_i d_i$. Minimal polynomials take the monic least common multiple, which is $d_r$ because of divisibility. If two matrices have the same list, both are similar to the same companion-block matrix and hence to each other. The converse follows from basis independence. In dimension zero the list is empty, and both polynomials are $1$.
\end{proof}

\Kentry{10.cyclic_test}{When a cyclic vector exists}
If $v$ is cyclic for $V$, its vector annihilator equals the operator's minimal polynomial, because every vector is $q(T)v$. Its degree is $n$, so the monic divisibility $m_T\mid\chi_T$ forces equality. Conversely, if $m_T=\chi_T$ and $n>0$, rational form gives
\[d_r=m_T=\chi_T=d_1\cdots d_r.\]
Cancelling $d_r$ in $\K[x]$ shows that the product of preceding factors is $1$. Since they were nonconstant and monic, there are none. Thus rational form has a single cyclic summand and a cyclic vector. On $V=0$, the vector $0$ generates the zero space, so the cyclic-vector convention includes it and again $m_T=\chi_T=1$.

\paragraph{The commutant criterion for cyclicity.}
An operator $T$ has a cyclic vector if and only if every operator commuting with $T$ is a polynomial in $T$.
If $v$ is cyclic and $ST=TS$, write $Sv=p(T)v$. Every vector is $q(T)v$, and
\[S(q(T)v)=q(T)Sv=q(T)p(T)v=p(T)(q(T)v),
\]
so $S=p(T)$. Conversely, on a nonzero space without a cyclic vector, rational form has at least two nonzero summands $R/(d_1),\ldots,R/(d_r)$, where $R=\K[x]$ and $d_1\mid d_r$. Reduction modulo $d_1$ gives a well-defined, nonzero $R$-linear map $R/(d_r)\to R/(d_1)$: a difference divisible by $d_r$ is divisible by $d_1$, and $[1]$ maps to the nonzero class $[1]$. Define $S$ on the full direct sum to be this map from the last summand to the first, and zero on every other summand. It is $R$-linear, hence commutes with multiplication by $x$, namely $T$. Every polynomial in $T$ preserves each summand separately, whereas $S$ sends the last summand nontrivially into the disjoint first one. Thus $S$ is not a polynomial in $T$. Transport through the rational-decomposition isomorphism gives the required counterexample in $\End(V)$. The zero space has one endomorphism, which is polynomial in its unique $T$, and is cyclic by the stated convention.

\Kentry{10.jordan_chains}{Constructing independent Jordan chains}
\begin{theorem}
A nilpotent operator has a basis that is a disjoint union of chains $v_1,\ldots,v_a$ satisfying $Nv_1=0$ and $Nv_{j+1}=v_j$.
\end{theorem}
\begin{proof}
For a nilpotent $N$, its minimal polynomial divides $x^r$ for some $r$, so every invariant factor in rational form is a positive power $x^a$. On the corresponding cyclic summand $\K[x]/(x^a)$, take the ordered vectors
\[x^{a-1},x^{a-2},\ldots,x,1.\]
They are a basis by uniqueness of remainders of degree less than $a$, and multiplication by $x$ kills the first and takes each later vector to its predecessor. Transport these bases through the direct-sum module isomorphism from rational form. Independence is guaranteed by the quotient bases and the direct sum; it is not assumed from arbitrary choices of preimages. The zero space has the empty union of chains.
\end{proof}

\Kentry{10.jordan_existence}{Jordan form over a splitting field}
\begin{theorem}
If $\chi_T$ splits over $\K$, then some basis gives a block diagonal matrix with blocks
\[J_a(\lambda)=\begin{pmatrix}\lambda&1&&0\\&\lambda&\ddots&\\&&\ddots&1\\0&&&\lambda\end{pmatrix}.
\]
The ones are on the superdiagonal.
\end{theorem}
\begin{proof}
Generalized eigenspace decomposition gives $V=\bigoplus_\lambda G_\lambda$, and on $G_\lambda$ the operator $N=T-\lambda I$ is nilpotent. Apply the preceding chain construction to this nilpotent restriction. In a chain ordered from the killed vector to its successive preimages, $Tv_1=\lambda v_1$ and $Tv_j=\lambda v_j+v_{j-1}$ for $j>1$. These are precisely the columns of $J_a(\lambda)$. Concatenating all chains over all summands gives the asserted basis. Splitting is used only to have the generalized eigenspace decomposition over the stated field; no scalar extension is hidden in the assertion.
\end{proof}

\Kentry{10.jordan_uniqueness}{Uniqueness of block sizes}
For one nilpotent block of length $a$, $\dim\ker N^j=\min(j,a)$, because in its chain the first $\min(j,a)$ vectors are killed and the remaining images are distinct chain vectors. Dimensions add over blocks. Hence
\[b_j:=\dim\ker N^j-\dim\ker N^{j-1}\]
is the number of blocks of length at least $j$. The number of blocks of length exactly $j$ is $b_j-b_{j+1}$. Since sufficiently large kernels equal the whole space, these formulas recover every block size. For a general split operator apply them to $(T-\lambda I)|_{G_\lambda}$ for each eigenvalue. The kernels and their dimensions are intrinsic, so the block data are unique up to block permutation. Conversely, equal block data give identical block matrices after reordering and thus similar operators.

\Kentry{10.semisimple}{Semisimplicity over the stated field}
Call $T$ semisimple if every invariant subspace has an invariant complement. In module language, invariant subspaces are exactly $\K[x]$-submodules, so this is the assertion that every submodule is a direct summand.
\begin{theorem}
$T$ is semisimple if and only if $m_T$ has no repeated irreducible factor over $\K$.
\end{theorem}
\begin{proof}
Suppose $m_T=\prod_i p_i$ with distinct monic irreducibles. Primary decomposition gives $V=\bigoplus_iV_i$, where $p_i(T)=0$ on $V_i$. The field $F_i=\K[x]/(p_i)$ acts on $V_i$ by polynomial evaluation. If $U\subset V$ is invariant, every polynomial projection onto $V_i$ preserves $U$, so $U=\bigoplus_i(U\cap V_i)$. Each intersection is an $F_i$-subspace. Extend an $F_i$-basis of it to an $F_i$-basis of $V_i$; this is a finite-dimensional basis extension, since $F_i$-independence implies $\K$-independence. Let $W_i$ be the span of the added basis vectors. It is an $F_i$-subspace and hence $T$-invariant. Then $\bigoplus_iW_i$ is the required complement of $U$.

Conversely, suppose $p^a$ with $a\ge2$ occurs in $m_T$, and let $W=\ker p(T)^a$ be its primary summand. On $W$ the operator $Q=p(T)|_W$ is nilpotent and nonzero. Indeed, if $p(T)|_W=0$, combining annihilators on all primary summands would reduce the exponent of $p$ in $m_T$, contrary to its minimality. Set $U=\ker Q$, a $T$-invariant subspace. Semisimplicity would supply a $T$-invariant complement $L$ to $U$ in $V$. Since $U\subset W$, decomposing any $w\in W$ as $u+\ell$ gives $\ell=w-u\in W$, so $W=U\oplus(L\cap W)$. On $L\cap W$, $Q$ is both nilpotent and injective, since its kernel is the intersection with $U$. An injective nilpotent map can act only on the zero space: apply injectivity repeatedly to $Q^r v=0$. Thus $L\cap W=0$, forcing $W=U$ and $Q=0$, a contradiction.
\end{proof}
Diagonalizability further requires the distinct irreducibles to be linear. Scalar extension can alter this distinction. The minimal polynomial of a matrix is unchanged by extension of the field: by the rank--minor criterion proved in \kref{05.adjugate}, linear independence of a finite list of its powers is characterized by a nonzero maximal minor of their flattened coordinate columns, and such a minor remains nonzero over an extension. In particular $I,A,\ldots,A^{d-1}$ remain independent when $d=\deg m_A$, while the original monic degree-$d$ polynomial still annihilates the extended matrix. Minimality and uniqueness therefore give the same minimal polynomial, not merely the same degree. If an irreducible polynomial acquires repeated roots in an extension, squarefreeness in the sense of distinct irreducible factors over the original field does not imply a product of distinct linear factors there. In characteristic zero, each irreducible is separable: its derivative is nonzero of lower degree, hence coprime to it. Bezout's identity remains true after scalar extension, so no repeated root appears.

\Kentry{10.additive_decomposition}{Diagonalizable and nilpotent parts}
\begin{theorem}
If $m_T$ splits over $\K$, there is a unique decomposition $T=D+N$ with $D$ diagonalizable, $N$ nilpotent, and $DN=ND$. Both parts are polynomials in $T$.
\end{theorem}
\begin{proof}
Let $V=\bigoplus_iG_{\lambda_i}(T)$ and let $P_i$ be the polynomial projections supplied by primary decomposition. Set $D=\sum_i\lambda_iP_i$ and $N=T-D$. On $G_{\lambda_i}$, $D=\lambda_iI$ and $N=T-\lambda_iI$, a nilpotent restriction. A common exponent at least all the nilpotency indices kills $N$ on every summand, so $N$ is nilpotent on $V$. The decomposition makes $D$ diagonalizable by choosing bases of the summands, and the polynomial formulas imply commutation.

For uniqueness, suppose $T=D'+N'$ is another such pair. Commutation implies that every eigenspace $E_\lambda(D')$ is invariant under $N'$, and on it $T-\lambda I=N'$ is nilpotent. Thus $E_\lambda(D')\subseteq G_\lambda(T)$. Since the eigenspaces of $D'$ sum to $V$ and distinct generalized eigenspaces of $T$ form a direct sum, these inclusions must be equalities for the nonzero summands: decompose a vector of $G_\lambda(T)$ into the eigenspaces of $D'$ and use directness of the generalized eigenspaces to remove all other components. Therefore $D'$ acts by $\lambda$ on every $G_\lambda(T)$, which characterizes $D$; then $N'=T-D'=N$. The zero-space case is included.
\end{proof}

\Kentry{10.semisimple_decomposition}{Semisimple and nilpotent parts over subfields of the complex numbers}
\begin{theorem}
If $\K\subset\C$, every finite-dimensional $T$ has a unique commuting decomposition $T=S+N$, where $S$ is semisimple over $\K$ and $N$ is nilpotent. Both are polynomials in $T$ with coefficients in $\K$.
\end{theorem}
\begin{proof}
The zero space is immediate. Write $m=m_T=\prod_i p_i^{a_i}$, set $f=\prod_i p_i$ and $r=\max_i a_i$. Work in the commutative algebra $R=\K[x]/(m)$, and let $J$ be the ideal generated by the image of $f$. Then $J^r=0$. Since $\K$ has characteristic zero, $f$ and $f'$ are coprime. We construct $y\in R$ with $y\equiv x\pmod J$ and $f(y)=0$.

Start with $y_0=x$. If $y\equiv x\pmod J$, then $f'(y)\equiv f'(x)\pmod J$. Bezout makes $f'(x)$ invertible modulo $J$. Every element invertible modulo a nilpotent ideal is invertible: if $uv=1-z$ with $z\in J$ and $z^r=0$, then $v(1+z+\cdots+z^{r-1})$ is an inverse of $u$. Thus define
\[y_{k+1}=y_k-f(y_k)f'(y_k)^{-1}.\]
For a polynomial over any commutative algebra, expansion of each monomial gives $f(y+h)=f(y)+f'(y)h+h^2q$ for some algebra element $q$. Taking $h=-f(y)f'(y)^{-1}$ cancels the first two terms. Consequently, if $f(y_k)\in J^s$, then $f(y_{k+1})\in J^{2s}$, and $y_{k+1}\equiv x\pmod J$. Initially $f(y_0)\in J$; after finitely many steps with $2^k\ge r$, obtain $f(y_k)=0$. Every operation occurred inside $R$, so $y_k$ has a representative polynomial $g\in\K[x]$.

Set $S=g(T)$ and $N=T-S$. The equation $f(y_k)=0$ gives $f(S)=0$, so $m_S$ divides the squarefree polynomial $f$, making $S$ semisimple. Also $x-y_k\in J$, hence $(T-S)^r=0$. Both are polynomials in $T$ and commute.

For uniqueness, extend the matrices to $\C$. A semisimple operator over $\K\subset\C$ has squarefree minimal polynomial over $\K$; separability in characteristic zero and invariance of the minimal polynomial under scalar extension, proved in \kref{10.semisimple}, make its complex extension diagonalizable. Thus any commuting semisimple--nilpotent decomposition over $\K$ becomes a diagonalizable--nilpotent decomposition over $\C$. Uniqueness in the preceding entry makes the two extended matrices equal, and injectivity of $\K\hookrightarrow\C$ makes the original matrices equal. This proof does not assert the theorem over an imperfect field of positive characteristic.
\end{proof}

\Kentry{10.complex_square_roots}{Square roots of invertible complex operators}
\begin{lemma}
If $Q^r=0$ over a field of characteristic zero, then $I+Q$ has a polynomial square root.
\end{lemma}
\begin{proof}
For $r=1$, $Q=0$ and use $I$. For $r>1$, define $a_0=1$, $a_1=1/2$, and, for $2\le k<r$,
\[a_k=-\frac12\sum_{i=1}^{k-1}a_i a_{k-i}.\]
In $q(t)=\sum_{k=0}^{r-1}a_kt^k$, the coefficient of $t$ in $q(t)^2$ is $1$, and for $2\le k<r$ the coefficient is $2a_k+\sum_{i=1}^{k-1}a_i a_{k-i}=0$. Thus $q(t)^2\equiv1+t\pmod{t^r}$, and substitution of $Q$ gives $q(Q)^2=I+Q$. No convergence or infinite analytic series is involved.
\end{proof}
\begin{theorem}
Every invertible operator on a finite-dimensional complex space has a square root.
\end{theorem}
\begin{proof}
On $G_\lambda(T)$ write $T=\lambda I+Q$, with $Q$ nilpotent. Invertibility implies $\lambda\ne0$. Choose $c\in\C$ with $c^2=\lambda$ and apply the lemma to $Q/\lambda$. Then $R_\lambda=cq(Q/\lambda)$ has square $T|_{G_\lambda}$. On the direct sum of the generalized eigenspaces define $R$ as the direct sum of these maps. Squaring preserves the blocks and gives $R^2=T$. On the zero space use its unique map.
\end{proof}
Neither uniqueness nor positivity is asserted: $I$ and $-I$ are distinct roots of $I$ on a nonzero complex space. Singular operators need not have roots. For $J=\begin{pmatrix}0&1\\0&0\end{pmatrix}$, any root $R$ would commute with $J=R^2$. Writing out $RJ=JR$ forces $R=\begin{pmatrix}a&b\\0&a\end{pmatrix}$. But $R^2=J$ would imply $a^2=0$ and $2ab=1$, an impossibility.

\Kentry{10.transpose_similarity}{Similarity to the transpose over the complex field}
For a Jordan block $J_a(\lambda)$ let $F_a$ be the permutation matrix reversing the order of its basis. Since $F_a^{-1}=F_a$, direct inspection of the diagonal and adjacent entries gives $F_a^{-1}J_a(\lambda)F_a=J_a(\lambda)^t$. For a block diagonal Jordan matrix $J$, use the block diagonal matrix $F$ of these reversals. If $A=PJP^{-1}$, then
\[A^t=(P^{-1})^tJ^tP^t=(P^{-1})^tF^{-1}JFP^t.
\]
Thus $A^t$ is similar to $J$, as is $A$, and they are similar to one another. The proof uses existence of Jordan form over $\C$. A theorem over arbitrary fields can be proved using rational forms, but it is not being inferred from an unavailable Jordan form over those fields.

% END modules/M05.10.tex


% BEGIN modules/M05.11.tex
\chapter{Bilinear, quadratic, and Hermitian forms}
Unless a different domain is displayed, a form is on a finite-dimensional vector space $V$. Bilinear forms are over a field $\K$. Hermitian forms are over $\C$ and are linear in the first argument, conjugate-linear in the second. Their diagonal values need not be positive.

\Kentry{11.bilinear}{Bilinear forms and pairings}
A bilinear pairing $B:V\times W\to\K$ is linear in each argument with the other held fixed. A bilinear form is the case $W=V$. Addition and scalar multiplication of pairings are defined pointwise and preserve bilinearity. In bases $(e_i)$ and $(f_j)$, expansion in each variable gives
\[B\Bigl(\sum_i x_ie_i,\sum_j y_jf_j\Bigr)=\sum_{i,j}x_iy_j B(e_i,f_j).
\]
Thus a pairing is uniquely specified by these finitely many values, and every array of such values defines a pairing by the displayed formula. This verifies existence and uniqueness of the coordinate construction rather than assuming it from a dimension count.

\Kentry{11.form_matrix}{Matrices of forms and congruence}
For a form and a basis $(e_i)$, set $A_{ij}=B(e_i,e_j)$. With column coordinate vectors, $B(x,y)=x^tAy$. If the columns of $P$ are a new basis expressed in the old one, then $x=Px'$ and $y=Py'$, so
\[B(x,y)=(x')^t P^tAPy'.\]
The new matrix is $P^tAP$. This is congruence, not similarity $P^{-1}AP$. The distinction is substantive: replacing coordinates in a form transforms both input slots, while replacing coordinates for an operator transforms its input and its output in opposite directions. For example, the one-dimensional matrices $[1]$ and $[4]$ over $\R$ are congruent by $P=[2]$, but are not similar.

\Kentry{11.radicals}{Left and right radicals}
The left radical is $\operatorname{rad}_L B=\{v:B(v,w)=0\text{ for all }w\}$; the right radical is defined with the slots interchanged. Each is a subspace, either by the defining linear equations or as the kernel of the corresponding map to $V^*$. In coordinates,
\[\operatorname{rad}_L B=\ker A^t,\qquad \operatorname{rad}_R B=\ker A.
\]
For $A=\begin{pmatrix}0&1\\0&0\end{pmatrix}$ the left radical is the second coordinate line and the right radical is the first, so they need not coincide. For symmetric or alternating forms they do coincide, by symmetry or skew-symmetry.

\Kentry{11.nondegenerate}{Nondegeneracy and rank}
Define $\rank B$ as the rank of the map $v\mapsto B(v,\cdot)$ into $V^*$; in coordinates this is $\rank A^t=\rank A$, so it is independent of basis. Rank-nullity gives
\[\rank B=n-\dim\operatorname{rad}_L B=n-\dim\operatorname{rad}_R B.
\]
For a square finite matrix, rank $n$, zero kernel, and invertibility are equivalent. Therefore the left radical is zero if and only if the right radical is zero, if and only if $A$ is invertible. A form satisfying these equivalent conditions is nondegenerate. For a pairing on two different spaces, the two nondegeneracy conditions must be kept separate until dimensions are specified.

\paragraph{A nondegenerate restriction and one-sided orthogonal complements.}
For this paragraph the ambient space $V$ may have arbitrary dimension, but $W\subseteq V$ is finite dimensional. Let $B$ be bilinear, or over $\R$ or $\C$ let it be linear in the first slot and conjugate-linear in the second (conjugation is the identity over $\R$). Put
\[W^{\perp_R}=\{v\in V:B(w,v)=0\text{ for every }w\in W\}.
\]
It is a subspace even in the sesquilinear case: the defining value on $av+bv'$ is $\overline aB(w,v)+\overline bB(w,v')$. If $B|_{W\times W}$ is nondegenerate, then
\[V=W\oplus W^{\perp_R}.
\]
Here is an explicit projection proof, also establishing the converse criterion $W\cap W^{\perp_R}=0$ if and only if the restriction is nondegenerate. Choose a basis $e_1,\ldots,e_r$ of $W$ and let $H_{ij}=B(e_i,e_j)$. In the sesquilinear case, for $v\in V$ solve
\[H\overline a=b(v),\qquad b(v)_i=B(e_i,v),
\]
and put $Ev=\sum_j a_je_j$. Invertibility of $H$ makes the solution unique, and $B(e_i,Ev)=B(e_i,v)$ makes $v-Ev\in W^{\perp_R}$. For $v\in W$ uniqueness gives $Ev=v$. For a vector in $W\cap W^{\perp_R}$, its coordinates satisfy $H\overline a=0$, hence are zero. The coordinate formula $a=\overline{H^{-1}b(v)}$ is linear in $v$, so $E$ is a linear projection. In the bilinear case omit the conjugates and the same proof applies. Conversely a singular $H$ has a nonzero null vector, giving a nonzero vector of $W\cap W^{\perp_R}$ (conjugate its coordinates in the sesquilinear convention). Thus the intersection criterion is equivalent to nonsingularity. For $W=0$, the empty equations give $E=0$ and $W^{\perp_R}=V$.

There is a corresponding left complement $W^{\perp_L}=\{v:B(v,w)=0\text{ for all }w\in W\}$. Apply the same argument to $B^{\mathrm{op}}(u,v)=B(v,u)$ in the bilinear case, or $B^{\mathrm{op}}(u,v)=\overline{B(v,u)}$ in the sesquilinear case, to obtain $V=W\oplus W^{\perp_L}$. No symmetry is assumed, so the left and right complements need not coincide. These are algebraic direct sums and should not be silently called orthogonal decompositions for an unrelated positive inner product.

\Kentry{11.symmetric_quadratic}{Symmetric and quadratic forms}
A bilinear form is symmetric if $B(u,v)=B(v,u)$. Its diagonal function $q(v)=B(v,v)$ satisfies $q(av)=a^2q(v)$ and has bilinear polar function
\[q(u+v)-q(u)-q(v)=B(u,v)+B(v,u).
\]
More generally, a quadratic form is a function $q:V\to\K$ with the stated degree-two homogeneity and a bilinear polar function. In characteristic two that polar function does not determine $q$: on $\K$, the functions $q_0(x)=0$ and $q_1(x)=x^2$ both have zero polar function, but $q_1(1)=1$. Thus the terms ``quadratic form'' and ``symmetric bilinear form'' cannot be identified without a characteristic hypothesis and an explicit conversion.

\Kentry{11.polarization}{Polarization when two is invertible}
If $\operatorname{char}\K\ne2$ and $q(v)=B(v,v)$ for symmetric $B$, the preceding expansion gives
\[B(u,v)=\frac{q(u+v)-q(u)-q(v)}2.
\]
Conversely, if $q$ is a quadratic form, set $B(u,v)$ equal to its polar function divided by $2$. It is symmetric and bilinear by definition. Its diagonal satisfies
\[B(v,v)=\frac{q(2v)-2q(v)}2=\frac{4q(v)-2q(v)}2=q(v).
\]
These constructions are inverse. The division by $2$ explains precisely why this proof does not extend to characteristic two.
In fact the converse proof needs only $q(2v)=4q(v)$ and bilinearity of the polar function: its diagonal calculation already recovers $q(v)$ from the symmetric bilinear $B$. That representation then implies full homogeneity $q(av)=a^2q(v)$. This proves the weaker-hypothesis characterization as well.

\Kentry{11.diagonal_congruence}{Orthogonal bases for symmetric forms}
\begin{theorem}
Over a field of characteristic different from two, every symmetric bilinear form has an orthogonal basis; the vectors with zero diagonal coefficient form a basis of its radical.
\end{theorem}
\begin{proof}
Induct on $n=\dim V$. The empty basis works when $n=0$. If $B=0$, every basis works. Otherwise some $e$ satisfies $B(e,e)\ne0$: if all diagonal values were zero, polarization would force $B=0$. Put $d=B(e,e)$ and
\[W=\{v:B(v,e)=0\}.
\]
Every $v$ decomposes as $d^{-1}B(v,e)e+w$, with $w\in W$. The intersection $\K e\cap W$ is zero because $d\ne0$, so $V=\K e\oplus W$. Symmetry makes the summands mutually orthogonal. Apply the induction hypothesis to $B|_W$ and prepend $e$. In a diagonal basis a vector is in the radical exactly when each coordinate times the corresponding diagonal entry is zero; hence precisely the zero-diagonal basis vectors span the radical. This proof uses only invertibility of $2$, so it establishes the characteristic-not-two extension of the narrower characteristic-zero source statement.
\end{proof}

\Kentry{11.inertia}{Sylvester's law over the real field}
\begin{theorem}
A real symmetric form has matrix $\diag(I_p,-I_q,0_r)$ in some basis, and the three integers $p,q,r$ are unique.
\end{theorem}
\begin{proof}
Diagonalize by the preceding theorem. Divide each vector with nonzero diagonal value $d$ by $\sqrt{\abs d}$ to change that value to $1$ or $-1$. Reorder the basis. The radical has dimension $r$, so $r$ is intrinsic. In this normal form let $V_+$ be the first $p$ coordinates and $V_{-0}$ the last $q+r$. Every nonzero vector in $V_+$ has positive diagonal value. If a subspace $U$ has dimension greater than $p$, then
\[\dim U+\dim V_{-0}>n,
\]
so $U\cap V_{-0}\ne0$ by the dimension formula. A vector in this intersection has nonpositive diagonal value. Thus $p$ is the largest possible dimension of a subspace on which the form is positive definite; this characterization is independent of basis. Repeating the argument with the negative of the form makes $q$ intrinsic. These invariants also prove that two normal forms with different triples cannot be congruent. Signature $p-q$ alone omits the radical and does not classify general degenerate forms.
\end{proof}

\Kentry{11.complex_classification}{Complex symmetric bilinear forms}
Over $\C$, diagonalization of a symmetric bilinear form gives entries $d_i$. For every nonzero $d_i$ choose $c_i$ with $c_i^2=d_i$, and divide its basis vector by $c_i$; the new diagonal entry is $1$. This gives $\diag(I_s,0)$, where $s=\rank B$. Rank is invariant under congruence because multiplication by the invertible matrices $P^t$ and $P$ preserves matrix rank. Hence two such forms are congruent exactly when their ranks agree. The rescaling uses squares, not squared absolute values, because the form is bilinear; it is not the classification of Hermitian forms.

\Kentry{11.alternating}{Alternating and skew-symmetric forms}
A bilinear form is alternating if $B(v,v)=0$ for every $v$. Expanding $0=B(u+v,u+v)$ gives $B(u,v)=-B(v,u)$ over every field. Conversely skew-symmetry gives $2B(v,v)=0$, which implies alternation when $2\ne0$ in the field. In characteristic two the form $B(x,y)=xy$ on a one-dimensional space is skew-symmetric because $-1=1$, yet $B(1,1)=1$. Thus the alternating hypothesis cannot be replaced by matrix skew-symmetry in that characteristic.

When $\operatorname{char}\K\ne2$, every bilinear form has the unique decomposition $B=S+A$ into a symmetric and an alternating form:
\[S(u,v)=\frac{B(u,v)+B(v,u)}2,\qquad
A(u,v)=\frac{B(u,v)-B(v,u)}2.
\]
Bilinearity follows by adding the two slotwise bilinear functions; swapping $u,v$ fixes $S$ and negates $A$, and $A(v,v)=0$. The formulas add to $B$. For uniqueness the difference of two symmetric parts is both symmetric and alternating, so its value equals its own negative, and is zero because $2$ is invertible. Consequently $B(v,v)=S(v,v)$: the alternating part contributes nothing to the quadratic form. This gives the explicit bridge from arbitrary bilinear representatives of a quadratic form to its unique symmetric representative.

\Kentry{11.symplectic_blocks}{Normal form for alternating forms}
\begin{theorem}
Over any field, an alternating form has a basis with matrix a direct sum of blocks $J=\begin{pmatrix}0&1\\-1&0\end{pmatrix}$ and a zero block. Its rank is even; nondegeneracy implies even dimension.
\end{theorem}
\begin{proof}
If the form is zero, choose any basis. Otherwise there are $u,v$ with $B(u,v)\ne0$. Rescale $v$ to obtain $B(u,v)=1$. Alternation implies $B(u,u)=B(v,v)=0$ and $B(v,u)=-1$, so the restriction to $H=\Span(u,v)$ has matrix $J$ and is nondegenerate. For any $x$, set
\[w=x-B(x,v)u+B(x,u)v.
\]
Direct calculation gives $B(w,u)=B(w,v)=0$. Therefore $V=H+H^\perp$, where $H^\perp$ consists of vectors orthogonal to both $u,v$. The intersection is zero: for $au+bv$ in it, pairing with $u$ gives $-b=0$, and pairing with $v$ gives $a=0$. Skew-symmetry makes orthogonality mutual. Thus $V=H\oplus H^\perp$, and the restriction on the complement remains alternating. Each step removes two dimensions, so induction ends with blocks $J$ and a subspace on which the form vanishes. The final zero block is the radical, while each $J$ has rank two. No division by two occurred, so the proof includes characteristic two, where the displayed $-1$ equals $1$.
\end{proof}

\Kentry{11.hermitian}{Sesquilinear and Hermitian forms}
A complex sesquilinear form in our convention is linear in its first argument and conjugate-linear in its second. In a basis with $H_{ij}=B(e_i,e_j)$ it is
\[B(x,y)=x^tH\overline y.
\]
It is Hermitian if $B(y,x)=\overline{B(x,y)}$, equivalently $H=(\overline H)^t$. In that case $B(v,v)$ is real because it equals its conjugate. With a basis change $P$, the matrix becomes $P^tH\overline P$. This convention is equivalent to the more usual quadratic-matrix expression $x^*Ax$: put $A=H^t$, then $x^tH\overline x=x^*Ax$ after transposing the scalar. This conversion will be used explicitly in the principal-minor test.

\paragraph{Triangular Gram reduction for a general sesquilinear form.}
Let $B$ be any sesquilinear form on a finite-dimensional complex space (or a bilinear form over $\R$). It need not be Hermitian or positive. Fix an ordered basis $e_1,\ldots,e_n$, write $H_{ij}=B(e_i,e_j)$, and suppose every leading principal determinant $\det H[1\!:\!k,1\!:\!k]$ is nonzero. Then there is a unique unit upper triangular matrix $P$ such that
\[P^tH\overline P\ \text{is lower triangular}.
\]
In the Hoffman--Kunze convention $A_{ij}=B(e_j,e_i)=H_{ji}$, this is equivalently the assertion that
\[P^*AP\ \text{is upper triangular},\qquad P^*=\overline P^{\,t}.
\]
Indeed, the two displayed matrices are transposes of one another; the leading determinants of $A$ and $H$ agree. This is a one-sided triangular statement, not a diagonalization theorem for an arbitrary form.

To prove it, set $W_k=\Span(e_1,\ldots,e_k)$. Each restriction to $W_k$ is nondegenerate by its leading determinant. The preceding one-sided decomposition supplies a projection $E_k:V\to W_k$ whose kernel is $W_k^{\perp_R}$; take $E_0=0$. Define
\[v_k=e_k-E_{k-1}e_k.
\]
Then $v_k-e_k\in W_{k-1}$ and $v_k\in W_{k-1}^{\perp_R}$. Consequently the matrix $P$ with columns $v_k$ in the original basis is unit upper triangular, hence invertible. For $i<j$, $v_i\in W_{j-1}$ and therefore $B(v_i,v_j)=0$. These are exactly the entries above the diagonal of $P^tH\overline P$, so that matrix is lower triangular.

For uniqueness, let another unit upper triangular matrix $Q$ have columns $w_k$ and satisfy the same lower triangular condition. Its first $k-1$ columns form a basis of $W_{k-1}$, and the zero entries above the diagonal give $B(w_i,w_k)=0$ for $i<k$. Hence $w_k\in W_{k-1}^{\perp_R}$. Also $w_k-e_k\in W_{k-1}$. In the direct sum $V=W_{k-1}\oplus W_{k-1}^{\perp_R}$ there is a unique such representative of $e_k+W_{k-1}$, namely $v_k$. Thus every column agrees and $Q=P$. The empty matrix is the unique solution when $n=0$. If the form is Hermitian, its transformed matrix is both Hermitian and lower triangular and is therefore diagonal; the general proof did not require that extra hypothesis.

\Kentry{11.hermitian_diagonal}{Recovering a complex form from its diagonal}
\begin{theorem}
For any complex sesquilinear form $B$ linear in the first slot, let $q(v)=B(v,v)$. Then
\[
 B(u,v)=\tfrac14\bigl(q(u+v)-q(u-v)+i q(u+iv)-i q(u-iv)\bigr).
\]
Moreover, $B$ is Hermitian if and only if $q(v)$ is real for every $v$.
\end{theorem}
\begin{proof}
Expanding the first difference gives $2(B(u,v)+B(v,u))$, while the second difference gives $2i(B(v,u)-B(u,v))$. Multiplying the latter by $i$ and adding leaves $4B(u,v)$, proving the formula without a Hermitian assumption. If $q$ is real-valued, put $a=B(u,v)$ and $b=B(v,u)$. The same expansions say that $a+b$ and $i(b-a)$ are real. Hence $a+b=\overline a+\overline b$ and $b-a=-\overline b+\overline a$. Adding these equations gives $2b=2\overline a$, so $b=\overline a$. Conversely the Hermitian identity with $u=v$ makes $q$ real. Over $\R$, a nonzero skew-symmetric bilinear form, such as the $2\times2$ block $J$, vanishes on the diagonal, so this complex statement has no unrestricted real analogue.
\end{proof}

\Kentry{11.hermitian_inertia}{Hermitian inertia without a spectral-theorem assumption}
\begin{theorem}
Every finite-dimensional Hermitian form is congruent, in the first-slot-linear convention, to $\diag(I_p,-I_q,0_r)$; the three integers are unique.
\end{theorem}
\begin{proof}
If the form is nonzero, diagonal recovery shows that some $e$ has $d=B(e,e)\ne0$, and $d$ is real. The vector $v-d^{-1}B(v,e)e$ is orthogonal to $e$; the Hermitian identity makes the orthogonality mutual. Exactly as in the symmetric diagonalization proof, $V=\C e\oplus\{v:B(v,e)=0\}$, an orthogonal direct sum. Induction produces a diagonal matrix with real entries. Dividing each nonzero-diagonal basis vector by the positive real number $\sqrt{\abs d}$ reduces it to $1$ or $-1$.

For uniqueness use complex subspace dimensions. In a displayed normal form, the positive coordinate space has dimension $p$. Any complex subspace of larger dimension intersects the negative-plus-zero coordinate space nontrivially and hence cannot be positive definite. Thus $p$ is the largest complex dimension of a positive-definite subspace. Applying this to $-B$ makes $q$ intrinsic, and the radical dimension is $r$. This proves uniqueness without assuming the spectral theorem, avoiding a circular route through operator diagonalization.
\end{proof}

\Kentry{11.definiteness}{Definiteness, semidefiniteness, and isotropy}
For a real symmetric or complex Hermitian form, positive definite means $B(v,v)>0$ for every $v\ne0$; positive semidefinite means $B(v,v)\ge0$ for every $v$. Negative variants apply these definitions to $-B$. Indefinite means that some vector has positive value and another has negative value. In the inertia normal form these cases respectively mean $(q,r)=(0,0)$, $q=0$, the corresponding reversed conditions, and $p,q>0$. These equivalences follow immediately by summing the signed squares or squared absolute values of coordinates. An isotropic vector is a nonzero vector with $B(v,v)=0$. For $\diag(1,-1)$, $(1,1)^t$ is isotropic but the form is nondegenerate, so isotropy need not mean radical membership. For a positive semidefinite form, the normal form has only positive and zero coordinates; hence zero diagonal value does imply radical membership. The zero-dimensional form is positive definite under the universal definition, since there are no nonzero vectors.

\Kentry{11.principal_minors}{Leading and all-principal-minor tests}
\begin{theorem}
Let $A$ be a real symmetric or complex Hermitian matrix, with quadratic expression $q_A(x)=x^*Ax$ (over $\R$, $*$ is transpose). Then $q_A$ is positive definite if and only if every leading principal minor is positive. It is positive semidefinite if and only if every principal minor is nonnegative.
\end{theorem}
\begin{proof}
For the strict assertion induct on the size. The empty matrix satisfies both sides vacuously. For positive size write
\[A=\begin{pmatrix}a&r^*\\r&B\end{pmatrix}.
\]
If $a>0$, let $C=B-a^{-1}rr^*$. Expansion gives
\[q_A(z,y)=a\abs{z+a^{-1}r^*y}^2+y^*Cy.
\]
Thus $A$ is positive definite exactly when $a>0$ and $C$ is positive definite. Elementary block row elimination gives $\det A_k=a\det C_{k-1}$ for each leading block $A_k$ of size $k\ge2$, where $C_{k-1}$ is the leading block of $C$. If all leading minors of $A$ are positive, the first is $a>0$ and all leading minors of $C$ are positive; induction and the completed-square formula give positive definiteness. Conversely positive definiteness gives $a>0$; minimizing the square by $z=-a^{-1}r^*y$ makes $C$ positive definite, and induction plus the determinant formula makes all leading minors positive.

For the semidefinite assertion, first suppose $A$ is positive semidefinite. Every principal restriction remains so. Its inertia normal form has only nonnegative entries, and determinant changes under congruence by the positive factor $\abs{\det P}^2$. Therefore its determinant is nonnegative.

Conversely suppose all principal minors of $A$ are nonnegative. For any leading block $A_k$ and any $\varepsilon>0$, determinant multilinearity, selecting either a column of $A_k$ or the corresponding column of $\varepsilon I$, yields
\[\det(A_k+\varepsilon I)=\sum_{J\subseteq\{1,\ldots,k\}}\varepsilon^{k-|J|}\det A_k[J,J].
\]
Indeed, once the identity columns outside $J$ are chosen, their fixed row positions leave precisely the principal submatrix indexed by $J$, with positive cofactor sign; the empty minor is $1$. Every summand is nonnegative and the empty-set summand is $\varepsilon^k>0$. All leading minors of $A+\varepsilon I$ are therefore positive, so the strict criterion gives $x^*Ax+\varepsilon\norm{x}_2^2>0$ for $x\ne0$. If $x^*Ax$ were negative, choosing $0<\varepsilon<-x^*Ax/\norm{x}_2^2$ would contradict this inequality. Hence $x^*Ax\ge0$ for every $x$.
\end{proof}
Leading minors alone do not suffice for semidefiniteness: $A=\diag(0,-1)$ has leading minors $0,0$ but takes the negative value $-1$ on the second basis vector. In the first-slot-linear form convention the matrix is $H=A^t$; corresponding principal determinants agree under transpose, so the test has the same content.

\Kentry{11.isometry_groups}{Groups preserving a nondegenerate form}
For a nondegenerate symmetric, Hermitian, or alternating form $B$, define its isometry group as the invertible linear maps $S$ with $B(Su,Sv)=B(u,v)$ for all $u,v$. The identity preserves the form, composition preserves it by two applications of the defining equality, and, for an invertible preserving $S$, substituting $u=S^{-1}a$, $v=S^{-1}b$ gives $B(S^{-1}a,S^{-1}b)=B(a,b)$. Thus these maps form a group. In finite dimension an arbitrary linear form-preserving map is already invertible when the form is nondegenerate: $Sv=0$ implies $B(v,w)=B(Sv,Sw)=0$ for every $w$, so $v=0$, and an injective endomorphism is bijective. In coordinates the bilinear condition is $P^tAP=A$; the Hermitian first-slot-linear condition is $P^tH\overline P=H$. Positive definite symmetric and Hermitian forms give orthogonal and unitary groups; a nondegenerate alternating form gives a symplectic group and requires even dimension. These are distinct form-preserving conditions.

\Kentry{11.form_operator}{Representing forms by operators}
Let $h$ be a positive definite real or complex inner product, linear in the first slot, and let $B$ be a bilinear or sesquilinear form with the same slot convention. Choose an $h$-orthonormal basis $(e_j)$, whose existence is proved in \kref{12.gram_schmidt}, and define
\[Au=\sum_j B(u,e_j)e_j.
\]
Linearity of $B$ in $u$ makes $A$ linear. If $v=\sum_jv_je_j$, then
\[h(Au,v)=\sum_j B(u,e_j)\overline{v_j}=B(u,v).
\]
If $A'$ has the same property, $h((A-A')u,v)=0$ for every $v$, and choosing $v=(A-A')u$ gives $(A-A')u=0$. Thus $A$ is unique, making the construction independent of the chosen basis. The form is Hermitian (symmetric over $\R$) exactly when
\[h(Au,v)=\overline{h(Av,u)}=h(u,Av),
\]
which is the defining identity for a self-adjoint operator, using the existence and uniqueness of adjoints in \kref{13.adjoint}. Finally $B(u,u)=h(Au,u)$ makes the positive definite and semidefinite conditions correspond term for term. This entry depends on the elementary inner-product construction, not on spectral classification; the preceding algebraic form classifications did not use this entry.

\Kentry{11.pfaffian}{The matrix Pfaffian and its normalization}
Let $A$ be a real skew-symmetric $2m\times2m$ matrix, and let $e^1,\ldots,e^{2m}$ be the ordered dual coordinate basis. Set
\[\omega_A=\sum_{i<j}A_{ij}e^i\wedge e^j.
\]
Define $\Pf(A)$ by
\[\frac1{m!}\underbrace{\omega_A\wedge\cdots\wedge\omega_A}_{m\text{ factors}}
 =\Pf(A)e^1\wedge\cdots\wedge e^{2m}.
\]
For $m=0$ both sides are the scalar $1$, so $\Pf$ of the empty matrix is $1$. Expansion of the wedge product, discarding terms with repeated indices and sorting the remaining indices, gives the equivalent signed-pairing formula
\[\Pf(A)=\frac1{2^m m!}\sum_{\sigma\in S_{2m}}\operatorname{sgn}(\sigma)
 \prod_{j=1}^m A_{\sigma(2j-1),\sigma(2j)}.
\]
The factor $2^m$ accounts for swapping the two entries within each pair, which changes both its permutation sign and its matrix entry sign. The factor $m!$ accounts for the order of the pairs, whose degree-two wedge factors commute. The exterior algebra and its top-degree determinant action are constructed independently in \kref{15.exterior} and \kref{15.determinant_line}.

\paragraph{Direct sums.}
For skew-symmetric matrices $A$ of size $2r$ and $B$ of size $2s$, order the coordinate basis with all coordinates of $A$ before all coordinates of $B$. Then
\[\Pf(\diag(A,B))=\Pf(A)\Pf(B).
\]
Indeed the associated two-form is $\omega_A+\omega_B$. Its degree-two summands commute under wedge product. In $(\omega_A+\omega_B)^{r+s}$, every term except the one with $r$ factors from $A$ and $s$ from $B$ vanishes: more than $r$ $A$-factors has degree exceeding $2r$ in the first coordinate space, and fewer forces more than $s$ $B$-factors. The surviving term has coefficient $\binom{r+s}{r}$. Dividing by $(r+s)!$ leaves
\[\frac{\omega_A^r}{r!}\wedge\frac{\omega_B^s}{s!},
\]
whose ordered volume coefficient is the product of the two Pfaffians. This includes $r=0$ or $s=0$ under the empty Pfaffian convention. A different interleaving of the two coordinate blocks must include the determinant sign of that permutation; the displayed formula fixes the order explicitly.

\begin{theorem}
For every real $2m\times2m$ matrix $P$, including singular $P$,
\[\Pf(P^tAP)=\det(P)\Pf(A),\qquad \Pf(A)^2=\det A.
\]
\end{theorem}
\begin{proof}
Let $P$ act as a linear map on the coordinate space. The pullback of $\omega_A$ has value $\omega_A(Pu,Pv)=u^tP^tAPv$, hence equals $\omega_{P^tAP}$. Pullback respects wedge multiplication, and on the top exterior power it multiplies by $\det P$. Applying it to the defining equation for $\Pf(A)$ proves the transformation identity. This argument does not use an inverse of $P$.

For the square identity, alternating-form normal form supplies an invertible $P$ with $P^tAP$ a direct sum of $r$ blocks $J=\begin{pmatrix}0&1\\-1&0\end{pmatrix}$ and a zero block. If $r<m$, its top wedge is zero, so its Pfaffian and determinant are both zero. If $r=m$, its two-form is $\sum_{j=1}^m e^{2j-1}\wedge e^{2j}$; the only nonzero terms in its $m$th wedge use every summand once, giving $m!$ times the ordered volume form. Thus its Pfaffian is $1$, and its determinant is also $1$, the product of the block determinants. In either case its Pfaffian squared equals its determinant. The transformation identity and determinant multiplicativity now give
\[(\det P)^2\Pf(A)^2=(\det P)^2\det A.
\]
Cancel the nonzero scalar $(\det P)^2$. The empty case was already fixed by convention.
\end{proof}

\paragraph{Bridge to Greub's skew-operator convention.}
This bridge is a real Euclidean specialization; Greub also treats nondegenerate bilinear metrics more generally. For an orthonormal vector basis $e_1,\ldots,e_{2m}$, Greub's map is
\[\Phi(a\wedge b)x=h(a,x)b-h(b,x)a.
\]
Consequently $\Phi(e_i\wedge e_j)$ has entry $-1$ at $(i,j)$ and $1$ at $(j,i)$. For an operator with skew matrix $A$, its inverse image is therefore $z=-\sum_{i<j}A_{ij}e_i\wedge e_j$. In Greub's divided-power notation, $z^m$ means the ordinary $m$-fold wedge divided by $m!$. If $a=c\,e_1\wedge\cdots\wedge e_{2m}$ with $c\ne0$, his scalar is
\[\operatorname{Pf}_a(A)=h(z^m,a)=c(-1)^m\Pf(A).
\]
Indeed the coefficient of $z^m$ in the orthonormal volume basis is $(-1)^m\Pf(A)$, and the squared norm of that volume basis is $1$. Thus
\[\operatorname{Pf}_a(A)^2=c^2\det A=h(a,a)\det A,
\]
which explains both the sign and the volume-norm factor. They cannot be discarded when changing conventions. The formulas used for this bridge occur in Greub Sections 8.2 and 8.4 (physical PDF pages 201, 207--209), and were checked against the rendered formulas rather than OCR alone.

% END modules/M05.11.tex


% BEGIN modules/M05.12.tex
\chapter{Inner products and orthogonal geometry}
Throughout this chapter the scalar field is $\mathbb F=\R$ or $\C$. Inner products are linear in the first argument. Unless a finite-dimensional hypothesis is stated, the ambient space can be infinite-dimensional; all sums displayed here are finite.

\Kentry{12.inner_product}{Real and complex inner products}
\begin{definition}
An inner product on an $\mathbb F$-vector space $V$ is a map $h:V\times V\to\mathbb F$ such that
\[
h(au+bv,w)=a h(u,w)+b h(v,w),\qquad
h(v,u)=\overline{h(u,v)},\qquad h(v,v)>0\quad(v\ne0).
\]
The last condition includes that $h(v,v)$ is real. We write $\ip{u}{v}=h(u,v)$. Over $\R$, conjugation is the identity, so this is a positive-definite symmetric bilinear form.
\end{definition}
Hermitian symmetry and first-slot linearity imply
$\ip{u}{av+bw}=\bar a\ip{u}{v}+\bar b\ip{u}{w}$, and $\ip{0}{v}=\ip{v}{0}=0$.
Indeed, conjugate the corresponding first-slot identities. Thus the definition agrees with the sesquilinear convention of \kref{11.hermitian}. On $\mathbb F^n$ the formula $\ip{x}{y}=\sum_j x_j\bar y_j$ satisfies the axioms because $\sum_j|x_j|^2$ is positive unless all coordinates vanish. On the zero space the zero-valued form satisfies positive definiteness vacuously.

\Kentry{12.convention_bridge}{Changing the linearity convention}
\begin{proposition}
If $h$ is linear in its first argument, then
$g(u,v)=h(v,u)=\overline{h(u,v)}$ is conjugate-linear in its first argument and linear in its second. This operation is its own inverse and preserves norms and orthogonality.
\end{proposition}
\begin{proof}
For example $g(au+bv,w)=h(w,au+bv)=\bar a g(u,w)+\bar b g(v,w)$, while $g(u,av+bw)=a g(u,v)+b g(u,w)$. Hermitian symmetry and positivity are unchanged, because $g(v,v)=h(v,v)$; a scalar is zero exactly when its conjugate is zero. Applying the argument swap twice returns $h$. In particular, the coefficient $h(v,e)$ in an orthonormal expansion becomes $g(e,v)$, not $g(v,e)$. This is the required translation of formulas as well as symbols.
\end{proof}

\Kentry{12.norm}{The induced norm and distance}
\begin{definition}
Define $\norm{v}=\sqrt{\ip{v}{v}}$, using the nonnegative real square root, and $d(u,v)=\norm{u-v}$.
\end{definition}
These are well defined because the diagonal values are nonnegative real numbers. Positivity gives $\norm{v}=0$ exactly when $v=0$, and sesquilinearity gives
$\norm{av}^2=|a|^2\norm{v}^2$, hence $\norm{av}=|a|\norm{v}$. The triangle inequality, and therefore the remaining metric axiom for $d$, are proved below.

\Kentry{12.cauchy_schwarz}{Cauchy--Schwarz, including equality}
\begin{theorem}
For $u,v\in V$, $|\ip{u}{v}|\le\norm{u}\norm{v}$. Equality holds exactly when the two-vector list $(u,v)$ is linearly dependent.
\end{theorem}
\begin{proof}
If $v=0$, both sides vanish and the list is dependent. If $v\ne0$, put
$a=\ip{u}{v}/\norm{v}^2$ and $r=u-av$. Direct calculation gives $\ip{r}{v}=0$ and
\[
0\le\norm{r}^2=\norm{u}^2-\frac{|\ip{u}{v}|^2}{\norm{v}^2}.
\]
Multiplication by the positive number $\norm{v}^2$ and nonnegative square roots give the inequality. Equality holds exactly when $\norm r=0$, equivalently $u=av$. With $v\ne0$, this is precisely linear dependence: in a nontrivial relation $bu+cv=0$, the coefficient $b$ cannot vanish, so $u=-c b^{-1}v$. Conversely, scalar multiples give equality by absolute homogeneity.
\end{proof}

\Kentry{12.norm_laws}{Norm laws and Pythagoras}
\begin{theorem}
The induced norm is a norm. If $u\perp v$, meaning $\ip{u}{v}=0$, then
$\norm{u+v}^2=\norm u^2+\norm v^2$.
\end{theorem}
\begin{proof}
The positivity and homogeneity axioms were proved in \kref{12.norm}. Expanding the square and using Cauchy--Schwarz gives
\[
\norm{u+v}^2=\norm u^2+2\operatorname{Re}\ip{u}{v}+\norm v^2
\le (\norm u+\norm v)^2.
\]
Both sides to be compared are nonnegative, so taking square roots proves the triangle inequality. If $u\perp v$, the cross term is zero, which proves Pythagoras. Consequently $d(u,v)$ is symmetric, is zero precisely on the diagonal, and satisfies $d(u,w)\le d(u,v)+d(v,w)$ by applying the triangle inequality to $u-w=(u-v)+(v-w)$.
\end{proof}

\begin{corollary}[Equality in the triangle inequality]
For real or complex inner-product vectors $u,v$,
$\norm{u+v}=\norm u+\norm v$ if and only if one vector is zero or $v=t u$ for a real $t>0$. Equivalently, for $u\ne0$, equality is exactly $v=t u$ with real $t\ge0$.
\end{corollary}
\begin{proof}
Zero vectors give equality directly. For two nonzero vectors, squaring shows that equality is equivalent to
$\operatorname{Re}\ip uv=\norm u\norm v$.
The inequalities $\operatorname{Re}\ip uv\le|\ip uv|\le\norm u\norm v$ then both become equalities. Cauchy--Schwarz equality gives $v=c u$ for some nonzero scalar $c$. With the first-slot-linear convention, $\ip uv=\bar c\norm u^2$; equality of its real part with its modulus forces $c$ to be positive real. Conversely a positive real multiple gives $\norm{u+tu}=(1+t)\norm u$. Thus complex linear dependence alone is insufficient: relative complex phases must align. For example $v=iu\ne0$ is dependent on $u$ but does not give triangle equality.
\end{proof}

\Kentry{12.parallelogram}{Parallelogram and polarization identities}
\begin{theorem}
For an inner-product norm,
\[
\norm{u+v}^2+\norm{u-v}^2=2\norm u^2+2\norm v^2.
\]
Over $\R$, $\ip{u}{v}=\tfrac14(\norm{u+v}^2-\norm{u-v}^2)$. Over $\C$ with our convention,
\[
\ip{u}{v}=\tfrac14\bigl(\norm{u+v}^2-\norm{u-v}^2
 +i\norm{u+iv}^2-i\norm{u-iv}^2\bigr).
\]
\end{theorem}
\begin{proof}
Expansion of the first two squares gives respectively
$\norm u^2+\norm v^2\pm2\operatorname{Re}\ip{u}{v}$. Their sum and difference give the first two assertions. If $z=\ip{u}{v}$, then $\ip{u}{iv}=-iz$, so the difference of the two squares involving $iv$ is $4\operatorname{Im}z$. Substitution proves the complex formula. These formulas show that an inner product is determined by its norm. They do not assert that every norm comes from an inner product.
\end{proof}

\Kentry{12.orthonormal}{Orthogonal and orthonormal families}
\begin{definition}
A family $(e_j)_{j\in J}$ is orthonormal if $\ip{e_j}{e_k}=\delta_{jk}$. A family is orthogonal if distinct members are orthogonal; normalization additionally requires that every vector be nonzero.
\end{definition}
\begin{proposition}
Every orthonormal family is linearly independent, and a finite orthogonal family of nonzero vectors becomes orthonormal upon replacing $v_j$ by $v_j/\norm{v_j}$.
\end{proposition}
\begin{proof}
If a finite sum $\sum_{j\in F}a_j e_j=0$, taking the inner product with $e_k$, for $k\in F$, gives $a_k=0$. This is the definition of independence even for an infinite family. For the normalization assertion, the denominators are positive real numbers, and direct sesquilinear expansion gives zero off the diagonal and one on the diagonal. A zero vector cannot be normalized and would already make a list dependent.
\end{proof}

\Kentry{12.gram_schmidt}{Gram--Schmidt with preservation of initial spans}
\begin{theorem}
For an independent finite list $v_1,\ldots,v_m$, define recursively
\[
w_j=v_j-\sum_{i<j}\ip{v_j}{e_i}e_i,\qquad e_j=w_j/\norm{w_j}.
\]
All denominators are positive. The list $(e_j)$ is orthonormal and
$\Span(e_1,\ldots,e_j)=\Span(v_1,\ldots,v_j)$ for every $j$.
\end{theorem}
\begin{proof}
At $j=1$, independence gives $v_1\ne0$, so the assertions hold. Suppose they hold through $j-1$. Then
$\ip{w_j}{e_k}=\ip{v_j}{e_k}-\sum_{i<j}\ip{v_j}{e_i}\delta_{ik}=0$ for $k<j$. If $w_j=0$, then $v_j$ lies in the span of the earlier $e_i$, hence in the span of the earlier $v_i$, contradicting independence. Thus normalization is valid. The formula for $w_j$ places $e_j$ in the old initial span extended by $v_j$, while the rearranged formula places $v_j$ in the span of $e_1,\ldots,e_j$. Together with the induction hypothesis, these give equality of the two initial spans. This completes the induction. For a dependent input list a zero residual is instead omitted if only a basis of its span is desired; it must never be divided by.
\end{proof}

\Kentry{12.on_basis}{Existence and extension of orthonormal bases}
\begin{theorem}
Every finite-dimensional inner-product space has an orthonormal basis, and every orthonormal list extends to such a basis.
\end{theorem}
\begin{proof}
Apply \kref{02.basis_extension} to obtain a basis. Gram--Schmidt produces an orthonormal list with the same full span, so it is a basis. For extension, first extend the given orthonormal list to an ordinary basis. Gram--Schmidt leaves its initial orthonormal vectors unchanged: all preceding inner products are zero and each norm is one. Hence the resulting orthonormal basis extends the original list. In dimension zero the empty list is the required basis, and no normalization is performed.
\end{proof}

\Kentry{12.fourier_parseval}{Fourier coefficients, Parseval, and Bessel}
\begin{theorem}
If $e_1,\ldots,e_n$ is an orthonormal basis, then
\[
v=\sum_{j=1}^n\ip{v}{e_j}e_j,\qquad
\norm v^2=\sum_{j=1}^n|\ip{v}{e_j}|^2.
\]
For an orthonormal list, not necessarily spanning, the sum of these squared coefficients is at most $\norm v^2$.
\end{theorem}
\begin{proof}
Expand $v=\sum a_j e_j$ using the basis property. Taking the inner product with $e_k$ gives $a_k=\ip{v}{e_k}$. Expanding the squared norm then leaves only diagonal terms. For a list of length $m$, let $p=\sum_{j=1}^m\ip{v}{e_j}e_j$ and $r=v-p$. The calculation $\ip r{e_k}=0$ shows $r\perp p$. Pythagoras gives $\norm v^2=\norm p^2+\norm r^2\ge\norm p^2=\sum_j|\ip v{e_j}|^2$. This proof also covers the empty list.
\end{proof}

\Kentry{12.orthogonal_complement}{Orthogonal complements}
\begin{definition}
For a subspace $U\subseteq V$, put
$U^\perp=\{v\in V:\ip vu=0\text{ for every }u\in U\}$.
\end{definition}
\begin{proposition}
$U^\perp$ is a subspace and $U\cap U^\perp=\{0\}$. Moreover $U\subseteq(U^\perp)^\perp$.
\end{proposition}
\begin{proof}
First-slot linearity preserves all equations $\ip vu=0$ under addition and scalar multiplication and places zero in $U^\perp$. A vector in the intersection satisfies $\ip vv=0$, hence is zero. Hermitian symmetry shows that each $u\in U$ is orthogonal to every vector in $U^\perp$, proving the last inclusion. The intersection assertion uses positivity and would fail for a general degenerate or indefinite bilinear form.
\end{proof}

\Kentry{12.orthogonal_decomposition}{Decomposition along a finite-dimensional subspace}
\begin{theorem}
If $U$ is finite-dimensional, then $V=U\oplus U^\perp$, even if $V$ is not finite-dimensional. If $V$ is finite-dimensional, then
$\dim U^\perp=\dim V-\dim U$ and $(U^\perp)^\perp=U$.
\end{theorem}
\begin{proof}
Choose an orthonormal basis $e_1,\ldots,e_m$ of $U$. For $v\in V$ set $u=\sum_j\ip v{e_j}e_j$. Then $u\in U$ and $\ip{v-u}{e_k}=0$ for each basis vector. Conjugate-linearity in the second slot extends this to every vector of $U$, so $v-u\in U^\perp$. If $v=u+w=u'+w'$ are two such decompositions, $u-u'=w'-w$ lies in the intersection and is zero. This proves both existence and uniqueness. In finite dimension, concatenating bases of the summands proves the dimension formula. The inclusion of $U$ in its double perpendicular was already proved, and applying the dimension formula twice shows their dimensions agree; they are therefore equal. The first decomposition also shows that the double-perpendicular equality holds whenever $U$ is finite-dimensional, regardless of the ambient dimension: write $x=u+w$ and use $x\perp w$ to obtain $\norm w^2=0$.
\end{proof}

\Kentry{12.projection}{Orthogonal projection and basis independence}
\begin{definition}
For finite-dimensional $U$, define $P_Uv$ to be the unique $U$-component in the preceding decomposition. Thus
$P_Uv=\sum_j\ip v{e_j}e_j$ in any orthonormal basis of $U$.
\end{definition}
\begin{proof}[Verification]
Uniqueness of the component makes the formula independent of the chosen basis. Adding decompositions and multiplying them by a scalar preserves membership in both summands, so uniqueness also proves linearity. For $u\in U$, its decomposition is $u+0$, hence $P_Uu=u$. Consequently $P_U^2=P_U$ and $\im P_U=U$. Its kernel is exactly $U^\perp$. Finally, decomposing $x=u+w$ and $y=u'+w'$ gives
$\ip{P_Ux}{y}=\ip u{u'}=\ip x{P_Uy}$; this is the self-adjointness identity whenever the adjoint is defined.
\end{proof}

\Kentry{12.best_approximation}{The unique nearest point}
\begin{theorem}
For finite-dimensional $U$, $P_Uv$ is the unique minimizer of $\norm{v-u}$ over $u\in U$. Precisely,
\[
\norm{v-u}^2=\norm{v-P_Uv}^2+\norm{P_Uv-u}^2.
\]
\end{theorem}
\begin{proof}
The two terms in $v-u=(v-P_Uv)+(P_Uv-u)$ lie in $U^\perp$ and $U$, so they are orthogonal. Pythagoras proves the identity. The second squared norm is nonnegative and vanishes exactly when $u=P_Uv$, proving existence and uniqueness of the minimum.
\end{proof}

\begin{theorem}[Conditional nearest-point criterion]
Let $U$ be any linear subspace of an inner-product space $V$, with no dimension or completeness assumption. For fixed $v\in V$ and $u_0\in U$, the vector $u_0$ minimizes $\norm{v-u}$ over $u\in U$ if and only if $v-u_0\in U^\perp$. If a minimizer exists, it is unique. For finite-dimensional $U$, the preceding projection theorem guarantees existence.
\end{theorem}
\begin{proof}
If $r=v-u_0\perp U$, then for every $u\in U$ the two vectors $r$ and $u_0-u$ are orthogonal. Hence
$\norm{v-u}^2=\norm r^2+\norm{u_0-u}^2$, proving minimality and uniqueness. Conversely suppose $u_0$ minimizes the distance. For any nonzero $z\in U$, put $a=\ip rz/\norm z^2$. Then $r-az\perp z$ and
$\norm{r-az}^2=\norm r^2-|\ip rz|^2/\norm z^2$.
But $u_0+az\in U$, so minimality forbids a strictly smaller value, forcing $\ip rz=0$. The same vanishing is automatic for $z=0$, so $r\in U^\perp$. No existence assertion for an arbitrary infinite-dimensional subspace was used.
\end{proof}

\begin{corollary}[Contractive projections]
For an idempotent linear map $P:V\to V$ on any inner-product space, $\norm{Px}\le\norm x$ for every $x$ if and only if $\im P\perp\ker P$. In that case $P$ is the orthogonal projection onto its image.
\end{corollary}
\begin{proof}
Idempotence gives $V=\im P\oplus\ker P$, by $x=Px+(I-P)x$. If the two summands are orthogonal, Pythagoras gives the asserted inequality. Conversely, for $u\in\im P$ and $z\in\ker P$, the inequality applied to $u-az$ gives $\norm u\le\norm{u-az}$ for every scalar $a$. The zero vector is therefore a nearest point to $u$ in the line $\K z$ (here $\K=\R$ or $\C$); the conditional criterion just proved gives $u\perp z$. The direct-sum decomposition is consequently the orthogonal one and identifies $P$ with its projection.
\end{proof}

\Kentry{12.riesz}{Finite-dimensional Riesz representation}
\begin{theorem}
For finite-dimensional $V$, each linear functional $\varphi:V\to\mathbb F$ has a unique representing vector $w$ such that $\varphi(v)=\ip vw$ for all $v$. The correspondence $w\mapsto\ip{-}{w}$ is conjugate-linear over $\C$ and linear over $\R$.
\end{theorem}
\begin{proof}
Choose an orthonormal basis and define $w=\sum_j\overline{\varphi(e_j)}e_j$. If $v=\sum_j a_j e_j$, then
$\ip vw=\sum_j a_j\varphi(e_j)=\varphi(v)$, proving existence. If $w'$ also represents $\varphi$, then $\ip v{w-w'}=0$ for all $v$; take $v=w-w'$ to obtain $w=w'$. The scalar behavior follows from conjugate-linearity of the second slot. In the zero space the unique functional and representing vector are both zero. No corresponding representation theorem for arbitrary algebraic functionals on an infinite-dimensional space is asserted here.
\end{proof}

% END modules/M05.12.tex


% BEGIN modules/M05.13.tex
\chapter{Adjoints and finite-dimensional spectral theory}
Here all inner-product spaces are finite-dimensional over $\mathbb F=\R$ or $\C$, unless an explicitly algebraic assertion is made over another field. The inner product is linear in its first argument. We write $T^\dagger$ for the adjoint and reserve $T^+$ for the pseudoinverse.

\Kentry{13.adjoint}{Construction of the adjoint}
\begin{theorem}
For a linear map $T:V\to W$ there is a unique linear map $T^\dagger:W\to V$ satisfying
\[
\ip{Tv}{w}_W=\ip{v}{T^\dagger w}_V\qquad(v\in V,\ w\in W).
\]
\end{theorem}
\begin{proof}
For fixed $w$, the left side is a linear functional in $v$. By \kref{12.riesz} it has a unique representing vector, which we call $T^\dagger w$. For $w=aw_1+bw_2$, the functional is $\bar a\ip{Tv}{w_1}+\bar b\ip{Tv}{w_2}$, represented by $aT^\dagger w_1+bT^\dagger w_2$. Uniqueness proves linearity. Any map satisfying the displayed identity has the same representing vector at every $w$, proving uniqueness. This argument also covers zero domain or codomain.
\end{proof}

\Kentry{13.adjoint_laws}{Algebra and matrices of adjoints}
\begin{theorem}
For maps of compatible types,
\[
(S+T)^\dagger=S^\dagger+T^\dagger,\quad
(aT)^\dagger=\bar a T^\dagger,\quad
(ST)^\dagger=T^\dagger S^\dagger,\quad
(T^\dagger)^\dagger=T.
\]
In orthonormal bases, the matrix of $T^\dagger$ is the conjugate transpose of the matrix of $T$.
\end{theorem}
\begin{proof}
Add the two defining adjoint identities to prove the sum law by uniqueness. For the scalar law, $\ip{aTv}{w}=a\ip v{T^\dagger w}=\ip v{\bar a T^\dagger w}$. For composition, $\ip{STv}{u}=\ip{Tv}{S^\dagger u}=\ip v{T^\dagger S^\dagger u}$. Taking conjugates of $\ip{Tv}{w}=\ip v{T^\dagger w}$ gives $\ip{T^\dagger w}{v}=\ip w{Tv}$, proving the last law. If $e_j$ and $f_i$ are the chosen bases, then $A_{ij}=\ip{Te_j}{f_i}=\ip{e_j}{T^\dagger f_i}=\overline{\ip{T^\dagger f_i}{e_j}}$, so the adjoint matrix entry $(j,i)$ is $\overline{A_{ij}}$. In particular $I^\dagger=I$, and if $T$ is invertible, applying the composition law to both inverse identities gives $(T^{-1})^\dagger=(T^\dagger)^{-1}$.
\end{proof}

\Kentry{13.adjoint_kernel}{Kernels, images, and orthogonal complements}
\begin{theorem}
For $T:V\to W$,
\[
\ker T^\dagger=(\im T)^\perp,\quad
\im T^\dagger=(\ker T)^\perp,\quad
\rank T^\dagger=\rank T,\quad
\ker(T^\dagger T)=\ker T.
\]
\end{theorem}
\begin{proof}
A vector $w$ is in $\ker T^\dagger$ exactly when $\ip v{T^\dagger w}=0$ for every $v$: the reverse implication follows by taking $v=T^\dagger w$. By the adjoint identity this is equivalent to $w\perp\im T$. Apply the same identity to $T^\dagger$ and use $(T^\dagger)^\dagger=T$ to get $\ker T=(\im T^\dagger)^\perp$. Double orthogonal complements in finite dimension yield the image formula. It implies $\rank T^\dagger=\dim V-\dim\ker T=\rank T$ by rank--nullity. Finally,
$\ip{T^\dagger Tv}{v}=\ip{Tv}{Tv}=\norm{Tv}^2$.
Thus $T^\dagger Tv=0$ implies $Tv=0$, and the other inclusion is immediate from composition.
\end{proof}

\Kentry{13.self_adjoint}{Self-adjointness and its limits}
\begin{definition}
An endomorphism $T$ is self-adjoint if $T=T^\dagger$. It is normal if $TT^\dagger=T^\dagger T$.
\end{definition}
Self-adjointness implies normality because both products then equal $T^2$. The converse fails: on $\R^2$, the matrix
$J=\begin{pmatrix}0&-1\\1&0\end{pmatrix}$ satisfies $J^\dagger=-J$ and $JJ^\dagger=J^\dagger J=I$, but $J\ne J^\dagger$ and $\det(tI-J)=t^2+1$ has no real root. Consequently a real normal operator need not have a real eigenbasis.

\begin{lemma}\label{aux:selfadjoint-diagonal-test}
If $D$ is self-adjoint and $\ip{Dv}{v}=0$ for every $v$, then $D=0$.
\end{lemma}
\begin{proof}
Set $b(u,v)=\ip{Du}{v}$. The adjoint identity shows $b(v,u)=\overline{b(u,v)}$. Expanding $b(u+v,u+v)-b(u-v,u-v)$ gives $4\operatorname{Re}b(u,v)$, and over $\C$ the analogous difference with $iv$ gives $4\operatorname{Im}b(u,v)$. These differences vanish by hypothesis; over $\R$ the real part is already the whole scalar. Hence $b(u,v)=0$ for all $u,v$. Taking $v=Du$ proves $\norm{Du}^2=0$, so $Du=0$ for every $u$.
\end{proof}

Over $\C$, one has stronger diagonal tests without assuming $T$ self-adjoint in advance:
\[
\bigl(\ip{Tv}{v}=0\text{ for all }v\bigr)\Longleftrightarrow T=0,
\qquad
\bigl(\ip{Tv}{v}\in\R\text{ for all }v\bigr)\Longleftrightarrow T=T^\dagger.
\]
Apply the full sesquilinear diagonal criterion in \kref{11.hermitian_diagonal} to $b(u,v)=\ip{Tu}{v}$ and the operator--form correspondence in \kref{11.form_operator}. Those entries prove that zero diagonal implies zero form and that real diagonal is equivalent to Hermitian symmetry; the latter is precisely the adjoint identity $T=T^\dagger$. Thus the earlier self-adjoint lemma over both fields is not mistaken for the stronger arbitrary-operator complex theorem. Over $\R$, the nonzero rotation $J$ above has $\ip{Jv}{v}=0$ for every $v$, so that stronger zero-diagonal assertion is false there.

\Kentry{13.unitary}{Isometries and unitary maps}
\begin{definition}
A linear map $U:V\to W$ is an isometry if $\ip{Ux}{Uy}=\ip xy$ for every $x,y$. A surjective isometry is called unitary over $\C$, or orthogonal over $\R$.
\end{definition}
\begin{proposition}
An isometry is injective and is characterized by $U^\dagger U=I_V$. For an endomorphism of a finite-dimensional space, this condition also implies $UU^\dagger=I_V$, and $U^{-1}=U^\dagger$.
\end{proposition}
\begin{proof}
The defining equality is equivalent to $\ip{x}{(U^\dagger U-I)y}=0$ for all $x,y$, hence to $U^\dagger U=I$. It implies $\norm{Ux}=\norm x$ and therefore injectivity. For an endomorphism, rank--nullity makes an injective map surjective, and multiplication of $U^\dagger U=I$ by its inverse gives $U^\dagger=U^{-1}$, hence the other product identity. More generally, a surjective isometry between distinct spaces is bijective by the same norm argument and again has inverse $U^\dagger$. Preservation of norms alone also characterizes a linear isometry, by the polarization identities in \kref{12.parallelogram}.
\end{proof}

\Kentry{13.normal}{Normal operators as a separate class}
An endomorphism is normal precisely when $TT^\dagger=T^\dagger T$. Self-adjoint maps and unitary endomorphisms are normal: in the first case both products are $T^2$, in the second both are $I$. The rotation $J$ above is a normal map that is not self-adjoint. Normality is therefore retained as its own condition rather than identified with self-adjointness. If $U$ is unitary, the adjoint laws show $(UTU^{-1})^\dagger=UT^\dagger U^{-1}$; multiplying verifies that normality is invariant under unitary equivalence.

\Kentry{13.normal_basic}{Elementary identities for normal maps}
\begin{theorem}
Normality is equivalent to $\norm{Tv}=\norm{T^\dagger v}$ for every $v$. For normal $T$,
\[
\ker T=\ker T^\dagger,\qquad \im T=\im T^\dagger,
\qquad V=\ker T\mathbin{\oplus^{\perp}}\im T.
\]
Every $T-\lambda I$ is normal, and
$Tv=\lambda v$ if and only if $T^\dagger v=\bar\lambda v$.
Distinct eigenspaces are orthogonal. Over $\C$, the operators
\[
A=(T+T^\dagger)/2,\qquad B=(T-T^\dagger)/(2i)
\]
are self-adjoint, $T=A+iB$, and $T$ is normal exactly when $AB=BA$.
\end{theorem}
\begin{proof}
The difference $D=T^\dagger T-TT^\dagger$ is self-adjoint, and
$\ip{Dv}{v}=\norm{Tv}^2-\norm{T^\dagger v}^2$. Normality makes this zero. Conversely, equality of the norms makes this zero for every $v$, and Lemma~\ref{aux:selfadjoint-diagonal-test} gives $D=0$. Equality of norms immediately identifies the two kernels. By \kref{13.adjoint_kernel}, their orthogonal complements identify the two images and give $\im T=(\ker T)^\perp$; the displayed orthogonal decomposition follows.

The adjoint of $T-\lambda I$ is $T^\dagger-\bar\lambda I$. Expanding the two products shows their difference is still $TT^\dagger-T^\dagger T$, so the shifted map is normal. Its two kernels coincide by the just-proved assertion, giving the eigenvector equivalence. If $Tu=\alpha u$ and $Tv=\beta v$, then
\[
\alpha\ip uv=\ip{Tu}{v}=\ip{u}{T^\dagger v}
=\ip u{\bar\beta v}=\beta\ip uv.
\]
When $\alpha\ne\beta$, division by their nonzero difference gives $\ip uv=0$. Finally the adjoint laws verify $A^\dagger=A$, $B^\dagger=B$ and $T=A+iB$. Direct multiplication gives
$AB-BA=(T^\dagger T-TT^\dagger)/(2i)$, proving the equivalence with commutation. Conversely, any expression $T=A+iB$ with self-adjoint $A,B$ forces the displayed formulas by adding and subtracting the adjoint identity, so the assertion does not depend on an unspecified choice of parts.
\end{proof}

\Kentry{13.self_adjoint_eigenvalues}{Eigenvalues of self-adjoint maps}
\begin{theorem}
Every eigenvalue of a complex self-adjoint operator is real, and distinct eigenspaces of a self-adjoint operator over $\R$ or $\C$ are orthogonal. Orthogonality also holds for complex normal operators, whose eigenvalues need not be real.
\end{theorem}
\begin{proof}
If $Tv=\lambda v$ with $v\ne0$ and $T=T^\dagger$, then
$\lambda\norm v^2=\ip{Tv}{v}=\ip v{Tv}=\bar\lambda\norm v^2$.
The norm square is positive, so $\lambda=\bar\lambda$. Orthogonality follows from \kref{13.normal_basic}, since self-adjoint maps are normal. Multiplication by $i$ on a nonzero complex inner-product space is unitary and normal, and has eigenvalue $i$; this proves that real eigenvalues cannot be claimed for all normal operators.
\end{proof}

\Kentry{13.normal_reducing}{Invariant subspaces reduce normal maps}
\begin{theorem}
For a normal endomorphism $T$ over $\R$ or $\C$, every $T$-invariant subspace $M$ is also $T^\dagger$-invariant, and $M^\perp$ is invariant under both maps. Every endomorphism $S$ commuting with $T$ also commutes with $T^\dagger$.
\end{theorem}
\begin{proof}
Choose orthonormal bases of $M$ and $M^\perp$. In their concatenation, invariance gives
\[
T=\begin{pmatrix}A&B\\0&C\end{pmatrix},\qquad
T^\dagger=\begin{pmatrix}A^*&0\\B^*&C^*\end{pmatrix},
\]
where a star denotes conjugate transpose of a matrix. The upper left blocks of $TT^\dagger=T^\dagger T$ imply
$AA^*+BB^*=A^*A$. The trace of $AA^*$ and of $A^*A$ is in both cases $\sum_{i,j}|A_{ij}|^2$. Taking traces therefore gives $\sum_{i,j}|B_{ij}|^2=0$. Every summand is nonnegative, so every entry of $B$ is zero. Both matrices are block diagonal, proving all the invariance assertions. Their restrictions to each summand are normal because the two block products still agree.

For the commutant assertion, equip $V\oplus V$ with the direct-sum inner product. The operator $\widetilde T(x,y)=(Tx,Ty)$ is normal. If $ST=TS$, the graph $G=\{(x,Sx):x\in V\}$ is $\widetilde T$-invariant, since $\widetilde T(x,Sx)=(Tx,STx)$. The first part makes $G$ invariant under $\widetilde T^\dagger$. Thus $(T^\dagger x,T^\dagger Sx)$ lies in $G$, which means $T^\dagger Sx=ST^\dagger x$ for every $x$. Finite dimension is essential to the trace argument; no infinite-dimensional extension is being assumed.
\end{proof}

\Kentry{13.complex_spectral}{The complex spectral theorem}
\begin{theorem}
A complex endomorphism has an orthonormal eigenbasis if and only if it is normal.
\end{theorem}
\begin{proof}
A diagonal matrix in an orthonormal basis commutes with its conjugate transpose, proving one implication. For the converse, induct on $n=\dim V$. At $n=0$ the empty basis works. If $n>0$, \kref{09.existence} gives an eigenvalue and a nonzero eigenvector; normalize it to $e$. By \kref{13.normal_basic}, $T^\dagger e=\bar\lambda e$. Consequently the orthogonal complement of $\mathbb C e$ is invariant under $T$ and its adjoint: for $x\perp e$, the two adjoint identities give $\ip{Tx}{e}=0$ and $\ip{T^\dagger x}{e}=0$. The adjoint of the restriction is the restriction of the adjoint, so the restriction is normal. Its dimension is $n-1$, and the induction hypothesis provides an orthonormal eigenbasis there. Append $e$ to obtain the required basis of $V$.
\end{proof}

\Kentry{13.real_spectral}{The real spectral theorem}
\begin{theorem}
A real endomorphism has an orthonormal eigenbasis if and only if it is self-adjoint.
\end{theorem}
\begin{proof}
A real diagonal matrix is its own transpose, proving necessity. For sufficiency, represent a self-adjoint $T$ in a real orthonormal basis by a real symmetric matrix $A$. Regard the same matrix as an operator on $\mathbb C^n$ with its standard inner product. Its conjugate transpose is still $A$, so it is complex self-adjoint. If $n>0$, it has a complex eigenpair $Az=\lambda z$ by \kref{09.existence}, and \kref{13.self_adjoint_eigenvalues} makes $\lambda$ real. Write $z=x+iy$ with real vectors $x,y$. Equality of real and imaginary parts yields $Ax=\lambda x$ and $Ay=\lambda y$; at least one of $x,y$ is nonzero. Hence $T$ has a real eigenvector, which can be normalized to $e$.

The real subspace $e^\perp$ is invariant, because $\ip{Tx}{e}=\ip x{Te}=\lambda\ip xe=0$. The restriction is self-adjoint. Induction on dimension, with the empty basis at dimension zero, now supplies an orthonormal eigenbasis of $e^\perp$, and adjoining $e$ completes the proof. The rotation in \kref{13.self_adjoint} shows why real normality alone is insufficient.
\end{proof}

\Kentry{13.schur}{Orthonormal triangularization}
\begin{theorem}
Every endomorphism of a finite-dimensional complex inner-product space has an upper-triangular matrix in an orthonormal basis. Over either $\R$ or $\C$, an endomorphism already known to be triangularizable can be triangularized orthonormally.
\end{theorem}
\begin{proof}
Suppose $v_1,\ldots,v_n$ is a triangularizing basis. Its initial spans $V_j=\Span(v_1,\ldots,v_j)$ satisfy $T(V_j)\subseteq V_j$, because the first $j$ matrix columns have zero entries below row $j$. Gram--Schmidt gives an orthonormal basis $e_1,\ldots,e_n$ with the same initial spans. Thus $Te_j\in\Span(e_1,\ldots,e_j)$, exactly the upper-triangular condition. Over $\C$, \kref{10.triangularization} provides the starting basis; its proof uses existence of an eigenvector and induction on the quotient dimension, not a spectral theorem. The empty basis handles dimension zero. Normality is not required anywhere in this argument.
\end{proof}

\Kentry{13.real_normal_blocks}{Orthogonal blocks for real normal maps}
\begin{theorem}
A real normal endomorphism has, in a real orthonormal basis, one-dimensional real blocks and two-dimensional blocks
$\begin{pmatrix}a&-b\\b&a\end{pmatrix}$ with $b>0$.
For an orthogonal endomorphism the one-dimensional blocks are $1$ or $-1$, and every two-dimensional block has $a^2+b^2=1$.
\end{theorem}
\begin{proof}
Use a real orthonormal basis to view $T$ as a real matrix $A$ and extend its coefficients to $\mathbb C^n$. The equality $AA^t=A^tA$ makes the complex extension normal. For $n>0$ it has an eigenvalue by \kref{09.existence}. If that eigenvalue is real, taking a nonzero real or imaginary part of an eigenvector gives a real eigenline, and \kref{13.normal_reducing} splits off its orthogonal complement.

Otherwise, conjugating the eigenpair lets us choose $\lambda=a+ib$ with $b>0$ and eigenvector $z=x+iy$. Since $A$ is real, $\bar z$ has eigenvalue $\bar\lambda\ne\lambda$. Normal-eigenspace orthogonality gives
\[
0=\ip z{\bar z}=\sum_j z_j^2
=\norm x^2-\norm y^2+2i\ip xy.
\]
Hence $x\perp y$ and $\norm x=\norm y=c>0$; positivity of $c$ follows from $z\ne0$. Comparing real and imaginary parts gives $Ax=ax-by$ and $Ay=bx+ay$. The orthonormal pair $e=x/c$, $f=-y/c$ therefore satisfies $Te=ae+bf$ and $Tf=-be+af$. Its span is invariant and has the asserted block. Its orthogonal complement is invariant and the restriction is normal by \kref{13.normal_reducing}. Induction on the remaining dimension, with dimension zero as base, completes the decomposition. In a one-dimensional orthogonal block, preservation of norm forces the scalar to have square one. For the displayed two-dimensional block, its transpose times itself is $(a^2+b^2)I$, so orthogonality is exactly $a^2+b^2=1$.
\end{proof}

\Kentry{13.algebraic_projections}{Algebraic eigenspace projections}
\begin{theorem}
Let $T$ be diagonalizable over a field $\K$, with distinct eigenvalues $\lambda_1,\ldots,\lambda_r$. Then
\[
p_i(t)=\prod_{j\ne i}\frac{t-\lambda_j}{\lambda_i-\lambda_j},\qquad P_i=p_i(T)
\]
is the projection onto $E_{\lambda_i}$ along the sum of the other eigenspaces. The projections sum to $I$, satisfy $P_i^2=P_i$, and satisfy $P_iP_j=0$ for $i\ne j$.
\end{theorem}
\begin{proof}
All denominators are nonzero. For $v\in E_{\lambda_j}$, polynomial evaluation gives $p_i(T)v=p_i(\lambda_j)v=\delta_{ij}v$. The eigenspaces span $V$ by diagonalizability, so this calculation determines the maps on all vectors and proves every stated identity component by component. In the zero space the spectrum is empty, the sum of projections is the zero map, and that map is also the identity of the zero space. No inner product enters this proof, so these projections need not be orthogonal.
\end{proof}

\Kentry{13.spectral_projections}{Orthogonal spectral projections}
\begin{theorem}
For real self-adjoint or complex normal $T$, its distinct eigenspaces give orthogonal projections $P_\lambda$ such that
$I=\sum_\lambda P_\lambda$ and $T=\sum_\lambda\lambda P_\lambda$. These are precisely the interpolation polynomials from \kref{13.algebraic_projections}.
\end{theorem}
\begin{proof}
The appropriate spectral theorem provides a full orthonormal eigenbasis. Grouping its vectors by eigenvalue gives an orthogonal direct sum of eigenspaces. The formula in \kref{12.projection} selects exactly the coordinates in one group. Thus $P_\lambda$ is the identity on that eigenspace and zero on the other groups. Summing these actions gives the two displayed decompositions. The interpolation polynomial has exactly the same actions on a spanning eigenbasis, so it is the same operator. This also proves pairwise orthogonality of the projections as maps: $P_\lambda P_\mu=0$ when $\lambda\ne\mu$, and each is self-adjoint.
\end{proof}

\Kentry{13.functional_calculus}{Functions on a finite spectrum}
\begin{definition}
For diagonalizable $T$ over $\K$ and any function $f:\{\lambda_1,\ldots,\lambda_r\}\to\K$, set $f(T)=\sum_i f(\lambda_i)P_i$, where the $P_i$ are the algebraic projections.
\end{definition}
\begin{proposition}
This construction is linear and multiplicative in $f$, sends the constant function $1$ to $I$, agrees with polynomial evaluation, and depends only on the values on the spectrum. In the orthogonal spectral setting over $\R$ or $\C$, $f(T)^\dagger=\bar f(T)$.
\end{proposition}
\begin{proof}
On $E_{\lambda_i}$, $f(T)$ is multiplication by $f(\lambda_i)$. Addition, scalar multiplication and composition therefore give the corresponding operations on the functions, and the identity assertion follows from $\sum_i P_i=I$. A polynomial $q(T)$ has the same action $q(\lambda_i)$ on every eigenspace, so the two definitions agree. Alternatively, interpolation constructs a polynomial representing any specified finite set of values. Since the eigenspaces span, equal values give equal operators. In the orthogonal case each $P_i^\dagger=P_i$, so the adjoint laws give $f(T)^\dagger=\sum_i\overline{f(\lambda_i)}P_i$. The empty-spectrum, zero-dimensional case satisfies all these identities because its identity map is zero.
\end{proof}

\Kentry{13.simultaneous_spectral}{Common orthonormal eigenbases}
\begin{theorem}
A possibly infinite pairwise commuting family of real self-adjoint operators has a common orthonormal eigenbasis. A possibly infinite pairwise commuting family of complex normal operators also has one.
\end{theorem}
\begin{proof}
Induct on $\dim V$, with the empty basis at zero. If every member of the family is scalar, any orthonormal basis works. Otherwise choose a nonscalar member $T$. Its spectral theorem shows that it has at least two distinct eigenspaces: a diagonalizable map with only one eigenvalue is scalar. Each eigenspace has smaller dimension. If $S$ commutes with $T$ and $Tv=\lambda v$, then $T(Sv)=S(Tv)=\lambda Sv$, so every member preserves each eigenspace. In the real case, restricting a self-adjoint map to such a subspace remains self-adjoint, because its perpendicular is invariant by the adjoint identity. In the complex case \kref{13.normal_reducing} shows that the subspace also reduces each normal $S$, so each restricted map remains normal. The restrictions still commute. Apply the dimension induction in each eigenspace and concatenate the resulting orthonormal bases. Distinct eigenspaces of $T$ are orthogonal, so the concatenation is orthonormal and is an eigenbasis for every member. This also covers an empty family, which falls in the all-scalar case vacuously.
\end{proof}

The induction above is on the dimension, not on the size of the family, so it applies verbatim to arbitrary families. In the real case it also covers commuting normal operators individually diagonalizable over $\R$: their real eigenspaces are mutually orthogonal by \kref{13.normal_basic}, and orthonormal bases inside those eigenspaces give real diagonal matrices, hence self-adjointness.

\begin{theorem}[Joint projections and a single generator]
For either family $\mathcal F$ above, let $W_1,\ldots,W_q$ be its distinct nonzero joint eigenspaces, grouping vectors that have the same eigenvalue for every member of $\mathcal F$. They form an orthogonal direct sum of $V$. Let $P_j$ be the orthogonal projection onto $W_j$. The unital algebra generated by $\mathcal F$ is exactly
\[
\mathcal A=\left\{\sum_{j=1}^q c_jP_j:c_j\in\mathbb F\right\}.
\]
It is closed under adjoints. For any distinct real numbers $t_1,\ldots,t_q$, the self-adjoint operator $D=\sum_jt_jP_j$ generates this same algebra: $\mathcal A=\mathbb F[D]$.
\end{theorem}
\begin{proof}
Take a common orthonormal eigenbasis and partition its finite set of vectors by their eigenvalue functions on $\mathcal F$. Their spans $W_j$ are intrinsic: if $v$ is a joint eigenvector with one such function, expansion in the basis forces every nonzero coefficient to have the same function. Hence the resulting spaces are exactly the joint eigenspaces. Every member of $\mathcal F$ is scalar on each $W_j$, and distinct $W_j$ are orthogonal. Every algebraic sum or product of family members and $I$ consequently lies in $\Span(P_1,\ldots,P_q)$.

For $j\ne k$, the two eigenvalue functions differ, so choose $T_{jk}\in\mathcal F$ whose values on $W_j,W_k$ differ. By \kref{13.algebraic_projections}, a polynomial $Q_{jk}$ in $T_{jk}$ is the projection onto its eigenspace with the value found on $W_j$. It acts as identity on $W_j$ and as zero on $W_k$, and acts by zero or one on each other joint space. The product $\prod_{k\ne j}Q_{jk}$ is therefore identity on $W_j$ and zero on every other joint space, hence equals $P_j$. All factors are in the generated algebra, so each $P_j$ belongs to it. This proves the algebra identity without assuming that $\mathcal F$ is finite; only finitely many pairs of nonzero joint spaces needed separation. Since $P_j^\dagger=P_j$, adjoints conjugate the coefficients and preserve this algebra.

The map $D$ is self-adjoint because all $t_j$ are real. Interpolation gives
\[
P_j=\prod_{k\ne j}\frac{D-t_kI}{t_j-t_k},
\]
by checking its action on each $W_k$. Thus every element of $\mathcal A$ is a polynomial in $D$, and the reverse inclusion follows from $D\in\mathcal A$. If $V\ne0$ and the family is empty, there is a single joint space $V$, its projection is $I$, and the algebra is $\mathbb F I$ (take $D=0$). For $V=0$, the joint-space list is empty, $I=0$, and both algebras consist of the unique zero map; take $D=0$ again.
\end{proof}

\Kentry{13.unitary_spectrum}{Spectral tests within the normal class}
\begin{theorem}
For a complex normal $T$, self-adjointness is equivalent to all eigenvalues being real, and unitarity is equivalent to every eigenvalue having absolute value one.
\end{theorem}
\begin{proof}
In an orthonormal eigenbasis, $T$ has diagonal entries $\lambda_i$, its adjoint has entries $\bar\lambda_i$, and $T^\dagger T$ has entries $|\lambda_i|^2$. Equality with $T$ or with $I$ is therefore equivalent to the respective scalar conditions. The arguments are valid for the empty diagonal as well. Without normality these tests fail: $\begin{pmatrix}1&1\\0&1\end{pmatrix}$ has only the real eigenvalue $1$, of modulus one, but its transpose is different from itself and its columns are not orthonormal.
\end{proof}

\begin{theorem}[Cayley transform and inverse]
On a finite-dimensional complex inner-product space, a self-adjoint $T$ determines a unitary operator
\[
U=(T-iI)(T+iI)^{-1},
\]
and $I-U$ is invertible. Conversely, for every unitary $U$ with $I-U$ invertible,
\[
T=i(I+U)(I-U)^{-1}
\]
is self-adjoint. These constructions are inverse bijections. The alternate convention $(I-iT)(I+iT)^{-1}$ is also unitary; its excluded eigenvalue is $-1$, instead of $1$.
\end{theorem}
\begin{proof}
Choose an orthonormal eigenbasis of the self-adjoint $T$. Its eigenvalues $\lambda$ are real, so $\lambda+i\ne0$. This proves $T+iI$ invertible. The displayed $U$ is diagonal in the same basis, with eigenvalues $z=(\lambda-i)/(\lambda+i)$. They have modulus one because numerator and denominator have equal squared modulus $\lambda^2+1$. Also $z\ne1$, since equality would give $-i=i$. Hence $U$ is unitary and $I-U$ is invertible by their diagonal matrices.

Conversely, a unitary $U$ is normal and has an orthonormal eigenbasis with eigenvalues $|z|=1$. The assumption makes $z\ne1$. In that basis the proposed $T$ has eigenvalues $t=i(1+z)/(1-z)$. They are real, since $\bar z=z^{-1}$ and
\[
\bar t=-i\frac{1+z^{-1}}{1-z^{-1}}
=i\frac{1+z}{1-z}=t.
\]
Thus its diagonal matrix is self-adjoint. Solving $z=(t-i)/(t+i)$ for $t$ gives exactly $t=i(1+z)/(1-z)$; all denominators already proved nonzero, so the two constructions are inverse on every basis vector and hence on operators. For the alternate convention, the corresponding eigenvalue is $(1-i\lambda)/(1+i\lambda)$, of modulus one and never $-1$, by the same scalar calculation. Zero-dimensional spaces satisfy all invertibility and empty-spectrum statements as well.
\end{proof}

\Kentry{13.normal_classification}{Classification up to unitary equivalence}
\begin{theorem}
Two normal operators on finite-dimensional inner-product spaces over the same field $\R$ or $\C$ are unitarily equivalent if and only if their characteristic polynomials are equal.
\end{theorem}
\begin{proof}
A unitary intertwining map is invertible, so its matrix conjugates the two operator matrices. Multiplicativity of determinant yields equality of characteristic polynomials. Conversely, over $\C$ the spectral theorem produces orthonormal eigenbases. Equality of characteristic polynomials gives the same eigenvalues with the same multiplicities, so order the two bases to have matching eigenvalue lists. The linear map sending one basis to the other is unitary and intertwines the two operators.

Over $\R$, use \kref{13.real_normal_blocks}. The one-dimensional block $a$ contributes $t-a$ and a block with $b>0$ contributes the irreducible polynomial $(t-a)^2+b^2$. Unique factorization in $\R[t]$ makes equality of characteristic polynomials equivalent to equality of these block multisets. The condition $b>0$ fixes the representative of each conjugate pair. Reorder the orthonormal block bases to obtain identical block matrices; the basis-to-basis map is then an orthogonal intertwiner. Equality of polynomial degrees also guarantees equal dimensions, including the zero-dimensional case.
\end{proof}

\begin{theorem}[Orthogonal cyclic decomposition of a normal operator]
A normal operator $T$ over $\R$ or $\C$ has an orthogonal decomposition
\[
V=C_T(w_1)\mathbin{\oplus^\perp}\cdots\mathbin{\oplus^\perp}C_T(w_r),
\qquad C_T(w)=\Span\{w,Tw,T^2w,\ldots\},
\]
whose nonconstant monic vector-annihilator polynomials $d_1,\ldots,d_r$ may be ordered so that $d_{j+1}\mid d_j$. The reversed list is the ascending invariant-factor list from \kref{10.invariant_factors}. Each $d_j$ is a product of distinct irreducible factors over the base field. The orthogonal projections onto the primary spaces of $T$ are polynomials in $T$.
\end{theorem}
\begin{proof}
Over $\C$, the spectral theorem expresses $V$ as an orthogonal sum of one-dimensional blocks, on which $T$ acts by a scalar $a$ and the minimal polynomial is $p(t)=t-a$. Over $\R$, \kref{13.real_normal_blocks} adds two-dimensional blocks with matrix
$\begin{pmatrix}a&-b\\b&a\end{pmatrix}$, $b>0$.
On such a block $H$, $(T-aI)^2=-b^2I$, so the irreducible polynomial $p(t)=(t-a)^2+b^2$ annihilates it. Every nonzero $v\in H$ is cyclic: if $v,Tv$ were dependent, $Tv=cv$ for a real $c$, contradicting $((c-a)^2+b^2)v=0$. Thus $v,Tv$ span $H$ and its vector-annihilator is $p$. The same cyclic assertion is immediate on each one-dimensional block.

Group blocks by their irreducible polynomial $p$, and denote them $H_{p,1},\ldots,H_{p,m_p}$. All blocks are mutually orthogonal and invariant. Put $r=\max_p m_p$ if $V\ne0$. For each $j\le r$ choose a nonzero $v_{p,j}\in H_{p,j}$ whenever $m_p\ge j$, and set
\[
w_j=\sum_{m_p\ge j}v_{p,j},\qquad
d_j(t)=\prod_{m_p\ge j}p(t).
\]
The factors are distinct irreducibles and hence pairwise coprime. The polynomial Chinese remainder construction in \kref{08.bezout_decomposition} supplies a polynomial equal to $1$ modulo one of these factors and $0$ modulo all the others. Applying it to $T$ and $w_j$ extracts the corresponding $v_{p,j}$. Since each such vector is cyclic in $H_{p,j}$, this proves
$C_T(w_j)=\bigoplus_{m_p\ge j}H_{p,j}$.
Conversely this direct sum is invariant and contains $w_j$, so no further components enter the cyclic span. A polynomial kills $w_j$ exactly when it kills every component, exactly when it is divisible by each participating $p$, hence by their product $d_j$. Thus $d_j$ is its monic vector-annihilator.

Every block occurs in exactly one of these layers, so the cyclic summands are mutually orthogonal and exhaust $V$. The set of factors in layer $j+1$ is a subset of that in layer $j$, proving $d_{j+1}\mid d_j$. Uniqueness of invariant factors in \kref{10.invariant_factors} identifies the reversed list. Finally the full primary space for $p$ is $\bigoplus_jH_{p,j}$. A Chinese-remainder polynomial which is $1$ modulo $p$ and $0$ modulo every other factor acts as identity on that primary space and as zero on its orthogonal complement. It is therefore its orthogonal projection and is a polynomial in $T$. For $V=0$, take the empty decomposition and factor list; there are no nonzero primary spaces to project onto.
\end{proof}
The summands in this theorem are orthogonal, but the power basis of a cyclic summand need not be orthonormal. Consequently this theorem does not assert unitary similarity to a matrix of ordinary companion blocks.

\Kentry{13.simultaneous_schur}{Simultaneous Schur form for an arbitrary commuting family}
\begin{theorem}
Every possibly infinite pairwise commuting family $\mathcal F\subseteq\End_\C(V)$ has a common orthonormal triangularizing basis.
\end{theorem}
\begin{proof}
By \kref{03.hom_dimension}, the endomorphism space has dimension $(\dim V)^2$. Its subspace $E=\Span(\mathcal F)$ is therefore finite-dimensional. Select a finite subfamily spanning $E$: start with the empty list; if the current span is not $E$, some member of $\mathcal F$ is outside it, since otherwise $E$ would be contained in it. Adjoin that member. Each step increases dimension by one, so the process stops after at most $\dim E$ steps. The selected maps still commute. By \kref{10.simultaneous_triangular}, they have a common triangularizing basis, hence a common invariant complete flag. Gram--Schmidt preserves that flag as proved in \kref{13.schur}, giving an orthonormal common triangularizing basis. Every member of $\mathcal F$ is a linear combination of the selected maps, and a linear combination of upper-triangular matrices is upper triangular. If $E=0$, every family member is zero and any orthonormal basis works. Thus neither the family size nor the zero-space case creates an exception.
\end{proof}

% END modules/M05.13.tex


% BEGIN modules/M05.14.tex
\chapter{Positive operators, singular values, and least squares}
All spaces in this chapter are finite-dimensional inner-product spaces over $\mathbb F=\R$ or $\C$, with first-slot-linear inner product. Dimensions, including zero, are allowed unless explicitly restricted. An adjoint is denoted by a dagger, a pseudoinverse by a plus sign.

\Kentry{14.positive}{Positive semidefinite and positive definite operators}
\begin{definition}
An endomorphism $T$ is positive semidefinite, abbreviated positive, if $T=T^\dagger$ and $\ip{Tv}{v}\ge0$ for every $v$. It is positive definite if it is self-adjoint and the inequality is strict for every nonzero $v$.
\end{definition}
Self-adjointness makes every diagonal value real, since $\ip{Tv}{v}=\ip v{Tv}=\overline{\ip{Tv}{v}}$. Over $\R$ the quadratic inequality alone does not imply self-adjointness. For example, if $J=\begin{pmatrix}0&-1\\1&0\end{pmatrix}$, then $\ip{(I+J)v}{v}=\norm v^2$ for real $v$, but $(I+J)^\dagger=I-J\ne I+J$. Thus self-adjointness is part of the definition, not an implicit consequence of the inequality.

\Kentry{14.positive_spectrum}{Spectral characterization and factorization of positivity}
\begin{theorem}
A self-adjoint operator is positive exactly when its eigenvalues are nonnegative, and positive definite exactly when they are all positive. Positive operators are exactly the operators of the form $R^\dagger R$; for the converse construction one can take $R$ to be an endomorphism of the same space. A positive operator is positive definite exactly when it is invertible.
\end{theorem}
\begin{proof}
The appropriate real or complex spectral theorem gives an orthonormal eigenbasis $e_i$, with real eigenvalues $\lambda_i$. For $v=\sum_i a_i e_i$,
\[
\ip{Tv}{v}=\sum_i\lambda_i|a_i|^2.
\]
If the quadratic form is nonnegative, putting $v=e_i$ gives $\lambda_i\ge0$. Conversely, nonnegative eigenvalues make every term nonnegative. Strict positivity for nonzero $v$ is equivalent to every $\lambda_i>0$: necessity again uses $e_i$, and sufficiency uses that some coefficient of a nonzero vector is nonzero.

For any $R:V\to W$, the adjoint laws make $R^\dagger R$ self-adjoint, and
$\ip{R^\dagger Rv}{v}=\norm{Rv}^2\ge0$.
Conversely, for positive $T$ define $Re_i=\sqrt{\lambda_i}e_i$ and extend linearly. Its diagonal is real, so $R^\dagger=R$, and $R^\dagger Re_i=\lambda_i e_i=Te_i$; equality on a basis proves the factorization. Finally, a diagonal matrix is invertible exactly when all its diagonal entries are nonzero, which for nonnegative $\lambda_i$ is equivalent to strict positivity. On the zero space the empty diagonal is invertible, and positive definiteness is vacuous; the statements remain valid.
\end{proof}

\begin{corollary}[Zero quadratic values of a positive operator]
For positive $T$, $\ip{Tv}{v}=0$ if and only if $Tv=0$.
\end{corollary}
\begin{proof}
In the orthonormal eigenbasis from the preceding proof,
$\ip{Tv}{v}=\sum_i\lambda_i|a_i|^2$ with $\lambda_i\ge0$.
A finite sum of nonnegative real numbers is zero exactly when every term is zero. This is equivalent to $a_i=0$ at every positive $\lambda_i$, and therefore to $Tv=\sum_i\lambda_i a_i e_i=0$. The reverse implication is also immediate from the inner-product definition. Zero eigenvalues are never divided by, and the empty sum gives the zero-space case.
\end{proof}

\Kentry{14.positive_root}{The unique positive square root}
\begin{theorem}
Each positive $T$ has a unique positive $S$ with $S^2=T$. On its orthogonal eigenspace $E_\lambda$, this root acts by $\sqrt\lambda$.
\end{theorem}
\begin{proof}
The construction $Se_i=\sqrt{\lambda_i}e_i$ in the preceding proof gives a positive map whose square is $T$. Equivalently, \kref{13.functional_calculus} defines it without choosing vectors inside eigenspaces, by the function $\lambda\mapsto\sqrt\lambda$.

For uniqueness, let $R$ be another positive square root. Since $R^2=T$, multiplication gives $RT=TR=R^3$. Thus $R$ preserves each eigenspace of $T$. Its restriction to such an eigenspace is self-adjoint and positive, because the inner product and its quadratic values restrict, and the adjoint identity holds on the invariant subspace. Apply its spectral theorem there. Every restricted eigenvalue $\mu$ is nonnegative and satisfies $\mu^2=\lambda$, because $R^2$ is $\lambda I$ on $E_\lambda$. Hence every restricted eigenvalue is $\sqrt\lambda$, so the restriction equals $\sqrt\lambda I$ on a spanning eigenbasis. This holds for every $E_\lambda$, whose direct sum is $V$, and therefore $R=S$. In particular the zero eigenspace is handled by $\mu^2=0$, not by division by an eigenvalue.
\end{proof}

\Kentry{14.singular_values}{Singular values and their multiplicities}
\begin{definition}
For $T:V\to W$, put $n=\dim V$. The singular values including zeros are the $n$ nonnegative square roots of the eigenvalues of $T^\dagger T$, counted with multiplicity and arranged in nonincreasing order. We write $s_1\ge\cdots\ge s_r>0$ for the positive values only; the remaining $n-r$ values are zero.
\end{definition}
The definition is valid because $T^\dagger T$ is positive. The spectral theorem and the characteristic polynomial show that its eigenvalue multiset is independent of the chosen orthonormal eigenbasis. Also $\ker(T^\dagger T)=\ker T$ by \kref{13.adjoint_kernel}; thus the zero eigenvalue has multiplicity $\dim\ker T$, and $r=n-\dim\ker T=\rank T$. In particular $r\le\min(\dim V,\dim W)$. If $n=0$, the entire list is empty. If $T=0$ and $n>0$, the full list consists of zeros while the positive list is empty.

\Kentry{14.svd}{Singular value decomposition}
\begin{theorem}
There are orthonormal bases $e_1,\ldots,e_n$ of $V$ and $f_1,\ldots,f_m$ of $W$ such that
\[
Te_i=s_i f_i\ (1\le i\le r),\qquad Te_i=0\ (r<i\le n).
\]
Consequently
$Tv=\sum_{i=1}^r s_i\ip v{e_i}f_i$,
and the matrix of $T$ in these bases is rectangular diagonal with nonnegative diagonal entries.
\end{theorem}
\begin{proof}
Choose an orthonormal eigenbasis of $T^\dagger T$, ordered so that its positive eigenvalues are $s_1^2,\ldots,s_r^2$. For $i\le r$ define $f_i=Te_i/s_i$. The adjoint identity gives
\[
\ip{Te_i}{Te_j}=\ip{e_i}{T^\dagger Te_j}=s_j^2\delta_{ij}.
\]
Dividing by the positive real numbers $s_i s_j$ shows that these $f_i$ are orthonormal. Extend them to an orthonormal basis of $W$ using \kref{12.on_basis}. For $i>r$, $\norm{Te_i}^2=\ip{T^\dagger Te_i}{e_i}=0$, so $Te_i=0$. Expand an arbitrary $v$ in the $e_i$ basis and apply linearity to get the displayed formula. The images of its first $r$ vectors are independent, and all images lie in their span, confirming that $r$ is the rank. If $r=0$, the formula is the empty sum and $T=0$; the two bases can be chosen independently. This includes zero domain or codomain.
\end{proof}

\begin{corollary}[Determinant and all singular values]
For any endomorphism $T$ on an $n$-dimensional real or complex inner-product space, including singular $T$, list all its singular values $s_1,\ldots,s_n$ including zeros. Then
\[
|\det T|=\sqrt{\det(T^\dagger T)}=\prod_{i=1}^n s_i.
\]
\end{corollary}
\begin{proof}
The Leibniz determinant sum and transposition identity give
$\det(T^\dagger)=\overline{\det T}$ in orthonormal coordinates. Multiplicativity yields
$\det(T^\dagger T)=\overline{\det T}\det T=|\det T|^2$.
The nonnegative real square root is therefore $|\det T|$. In an orthonormal eigenbasis of $T^\dagger T$, its diagonal entries are all $s_i^2$, so its determinant is $\prod_i s_i^2$. As every $s_i$ is nonnegative, the square root of this product is $\prod_i s_i$. For singular $T$ in positive dimension at least one $s_i$ is zero and both sides vanish. For $n=0$ the determinant and empty product are one. In particular the product of only the positive singular values is not the right formula for singular $T$.
\end{proof}

\Kentry{14.pseudoinverse}{Projection followed by inverse restriction}
\begin{definition}
For $T:V\to W$, let $M=(\ker T)^\perp$ and $N=\im T$. The restriction $R=T|_M:M\to N$ is invertible. Define
\[
T^+w=R^{-1}(P_Nw),
\]
viewing the result in the subspace $M\subseteq V$.
\end{definition}
\begin{proof}[Verification]
The kernel of $R$ is $M\cap\ker T=\{0\}$. To prove surjectivity, write any preimage $v$ of a vector $Tv\in N$ as $v=k+m$ with $k\in\ker T$ and $m\in M$; then $Tv=Tm=Rm$. Hence $R$ is a bijective linear map, whose inverse is linear. Since orthogonal projection is linear, the displayed construction is a well-defined linear map $W\to V$. It is characterized by the two conditions
\[
T^+w\in(\ker T)^\perp,\qquad T(T^+w)=P_{\im T}w.
\]
Indeed, any vector satisfying them is the unique inverse image under $R$. This characterization also proves basis independence. If $T=0$, the domain of $R$ and its codomain are both zero spaces, and $T^+=0$.
\end{proof}

\Kentry{14.penrose}{The four Penrose identities and uniqueness}
\begin{theorem}
The products $TT^+$ and $T^+T$ are respectively $P_{\im T}$ and $P_{(\ker T)^\perp}$. In particular,
\[
TT^+T=T,\quad T^+TT^+=T^+,\quad
(TT^+)^\dagger=TT^+,\quad (T^+T)^\dagger=T^+T.
\]
Conversely these four equations characterize $T^+$ among all linear maps $X:W\to V$.
\end{theorem}
\begin{proof}
The first projection identity is part of the defining characterization. For $v=k+m$ in $\ker T\oplus(\ker T)^\perp$, $Tv=Tm\in\im T$, so $T^+Tv=R^{-1}(Tm)=m$. This is the second projection identity. A projection fixes its image; moreover $\im T^+\subseteq(\ker T)^\perp$. These facts give the first two Penrose equations. Orthogonal projections are self-adjoint by \kref{12.projection}, giving the other two.

Suppose $X$ satisfies the four equations with $X$ in place of $T^+$. Then $(TX)^2=TXTX=(TXT)X=TX$, and $TX$ is self-adjoint. Its image lies in $\im T$, and $TXT=T$ shows it acts as the identity on $\im T$. Hence its image is exactly $\im T$. A self-adjoint idempotent $P$ is the orthogonal projection onto its image: it splits each $w$ as $Pw+(I-P)w$, and these components are orthogonal because $P(I-P)=0$ and $P=P^\dagger$. Thus $TX=P_{\im T}$.

Similarly $XT$ is a self-adjoint idempotent. Its kernel equals $\ker T$: one inclusion is immediate; if $XTv=0$, then $Tv=TXTv=0$ by $TXT=T$. Hence $XT=P_{(\ker T)^\perp}$, since an orthogonal projection has image the perpendicular of its kernel. Finally $XTX=X$ says that $Xw=(XT)Xw$ lies in $(\ker T)^\perp$. Together with $TXw=P_{\im T}w$, this is precisely the two-condition characterization of $T^+w$. Therefore $X=T^+$.
\end{proof}

\Kentry{14.least_squares}{All least-squares minimizers}
\begin{theorem}
For every $w\in W$, the minimizers of $\norm{Tv-w}$ over $v\in V$ are exactly
$T^+w+\ker T$. Among them, $T^+w$ uniquely has smallest norm. The least-squares minimizer is unique exactly when $T$ is injective.
\end{theorem}
\begin{proof}
As $v$ ranges over $V$, $Tv$ ranges over $N=\im T$. By \kref{12.best_approximation}, the unique nearest vector of $N$ to $w$ is $P_Nw$. Thus $v$ minimizes the residual exactly when $Tv=P_Nw=TT^+w$, which is equivalent to $v-T^+w\in\ker T$. Since $T^+w\perp\ker T$, for every such $k$,
\[
\norm{T^+w+k}^2=\norm{T^+w}^2+\norm k^2.
\]
The second term vanishes exactly at $k=0$, proving minimum-norm uniqueness. The entire affine set has one element exactly when $\ker T=\{0\}$, equivalently when $T$ is injective. The assertion includes $T=0$: then every vector minimizes the same residual $\norm w$, and zero alone has minimum norm.
\end{proof}

\Kentry{14.normal_equations}{Normal equations and the injective formula}
\begin{theorem}
A vector $v$ minimizes $\norm{Tv-w}$ exactly when
$T^\dagger(Tv-w)=0$. If $T$ is injective, there is one minimizer and it equals
$ (T^\dagger T)^{-1}T^\dagger w$.
\end{theorem}
\begin{proof}
The nearest-point condition is $Tv=P_{\im T}w$, equivalently $Tv-w\in(\im T)^\perp$, because $Tv$ already lies in $\im T$. By \kref{13.adjoint_kernel}, this perpendicular space is $\ker T^\dagger$, proving the equivalence. If $T$ is injective, $\ker(T^\dagger T)=\ker T=0$, so the square endomorphism $T^\dagger T$ is invertible. Multiplying the equation $T^\dagger Tv=T^\dagger w$ by this inverse gives the formula; substitution shows that it does solve the equation. If $T$ is not injective, the same kernel identity makes $T^\dagger T$ singular, and this inverse formula is not used.
\end{proof}

\Kentry{14.svd_pseudoinverse}{SVD formulas for adjoints and pseudoinverses}
\begin{theorem}
Using the positive singular-value expansion of \kref{14.svd},
\[
T^\dagger w=\sum_{i=1}^r s_i\ip w{f_i}e_i,
\qquad
T^+w=\sum_{i=1}^r s_i^{-1}\ip w{f_i}e_i.
\]
Only positive singular values are inverted.
\end{theorem}
\begin{proof}
Let $A(w)$ denote the first displayed sum. Direct expansion gives
\[
\ip v{A(w)}=\sum_i s_i\overline{\ip w{f_i}}\ip v{e_i}
=\sum_i s_i\ip v{e_i}\ip{f_i}w=\ip{Tv}w.
\]
The uniqueness of the adjoint therefore makes $A=T^\dagger$. The vectors $e_1,\ldots,e_r$ span $(\ker T)^\perp$ and $f_1,\ldots,f_r$ span $\im T$, as follows from the SVD formula and its independent images. The second sum belongs to the former span, and applying $T$ to it gives $\sum_i\ip w{f_i}f_i=P_{\im T}w$. The defining characterization of the pseudoinverse proves the second identity. When $r=0$, both sums vanish, in agreement with the definitions.
\end{proof}

\Kentry{14.operator_norm}{Operator norm and the largest singular value}
\begin{definition}
For $T:V\to W$ define
\[
\norm T=\sup\{\norm{Tv}:v\in V,\ \norm v\le1\}.
\]
The unit ball, unlike the unit sphere, always contains zero. In particular this definition gives zero when $V=\{0\}$.
\end{definition}
\begin{theorem}
The supremum exists and equals $s_1$ if $T\ne0$, and equals zero if $T=0$. Furthermore $\norm{T^\dagger}=\norm T$ and $\norm{ST}\le\norm S\norm T$ whenever composition is defined.
\end{theorem}
\begin{proof}
For the SVD coefficients, Parseval and orthonormality give
\[
\norm{Tv}^2=\sum_{i=1}^r s_i^2|\ip v{e_i}|^2
\le s_1^2\norm v^2\quad(r>0).
\]
Thus the set defining the supremum is nonempty and bounded, and $v=e_1$ attains $s_1$. If $r=0$, every output is zero and the supremum is zero. This proves existence and the formula without an empty-spectrum maximum. It also proves $\norm{Tv}\le\norm T\norm v$ for every $v$, including zero. The adjoint formula in \kref{14.svd_pseudoinverse} has the same positive singular values with the two orthonormal lists interchanged, so its largest value, or zero if the list is empty, is the same.

For $\norm v\le1$, the two pointwise bounds give $\norm{STv}\le\norm S\norm{Tv}\le\norm S\norm T$. Taking the supremum proves submultiplicativity. Positivity, absolute homogeneity, and the triangle inequality for this operator norm follow respectively from the pointwise norm axioms: if all outputs have norm zero then $T=0$; scaling multiplies every bound by $|a|$; and $\norm{(S+T)v}\le\norm{Sv}+\norm{Tv}\le\norm S+\norm T$ on the unit ball. Thus this is indeed a norm on the vector space of linear maps.
\end{proof}

\Kentry{14.low_rank}{Best approximation of bounded rank}
\begin{theorem}
Let $s_1\ge\cdots\ge s_r>0$ be the positive singular values of $T$. For any integer $k\ge0$, set
\[
T_kv=\sum_{i=1}^{\min(k,r)}s_i\ip v{e_i}f_i.
\]
Then $\rank T_k\le k$, and it minimizes $\norm{T-S}$ among maps with $\rank S\le k$. The minimum is $s_{k+1}$ if $k<r$ and zero if $k\ge r$.
\end{theorem}
\begin{proof}
The image of $T_k$ lies in the span of its at most $k$ displayed output vectors, so its rank is at most $k$. If $k\ge r$, then $T_k=T$ and zero error is the smallest possible value. Suppose $k<r$. The remaining SVD sum gives $\norm{T-T_k}=s_{k+1}$ by the operator-norm calculation, or directly by the same squared-coefficient bound with equality at $e_{k+1}$.

Let $S$ have rank at most $k$. On the $(k+1)$-dimensional subspace $E=\Span(e_1,\ldots,e_{k+1})$, its restriction has rank at most $k$. Rank--nullity gives a nonzero $x\in E\cap\ker S$; divide by its positive norm so that $\norm x=1$. Writing $x=\sum_{i=1}^{k+1}a_i e_i$ gives $\sum_i|a_i|^2=1$ and
\[
\norm{(T-S)x}^2=\norm{Tx}^2
=\sum_{i=1}^{k+1}s_i^2|a_i|^2\ge s_{k+1}^2.
\]
Therefore $\norm{T-S}\ge s_{k+1}$, proving optimality. This argument includes $k=0$, where every rank-zero map is zero. No uniqueness of the minimizer is asserted.
\end{proof}

\Kentry{14.frobenius}{Frobenius norm, trace, and comparison inequalities}
\begin{theorem}
In orthonormal bases, define $\norm T_F^2=\sum_{i,j}|A_{ij}|^2$ for the matrix $A$ of $T$. This is independent of the orthonormal bases and satisfies
\[
\norm T_F^2=\tr(T^\dagger T)=\sum_{i=1}^r s_i^2,
\qquad
\norm T\le\norm T_F\le\sqrt{\rank T}\,\norm T.
\]
\end{theorem}
\begin{proof}
For the domain basis $e_j$, Parseval in the codomain gives $\sum_i|A_{ij}|^2=\norm{Te_j}^2$. Also $\ip{T^\dagger Te_j}{e_j}=\norm{Te_j}^2$. Summing over $j$ identifies the squared matrix expression with the sum of the diagonal entries of $T^\dagger T$, namely its trace. Trace is basis-independent by \kref{04.trace}, while the adjoint is intrinsically defined, proving independence of both orthonormal bases. In an eigenbasis of $T^\dagger T$ the trace is $\sum_{i=1}^r s_i^2$. If $r>0$, then
$s_1^2\le\sum_i s_i^2\le r s_1^2$; take nonnegative square roots and use \kref{14.operator_norm}. If $r=0$, all three quantities in the inequality are zero, so the inequalities hold there too. The trace identity in particular proves invariance under unitary changes of coordinates, as needed below.
\end{proof}

\Kentry{14.polar}{Polar decomposition of an endomorphism}
\begin{theorem}
Every endomorphism $T:V\to V$ can be written
$T=U\sqrt{T^\dagger T}$ with $U$ unitary over $\C$ or orthogonal over $\R$. If $T$ is invertible, this $U$ is unique. A singular map can have more than one such $U$.
\end{theorem}
\begin{proof}
Choose full orthonormal domain and codomain bases as in the SVD. Since both spaces are $V$, the two bases have the same length $n$. Define $Ue_i=f_i$ for all $1\le i\le n$. A map carrying one orthonormal basis to another preserves every inner product by expanding coordinates and is surjective, so it is unitary. The positive root $P=\sqrt{T^\dagger T}$ satisfies $Pe_i=s_i e_i$ for $i\le r$ and $Pe_i=0$ otherwise. Thus $UPe_i=Te_i$ for each basis vector, proving $UP=T$. If $T$ is invertible, $P$ is invertible, so necessarily $U=TP^{-1}$. If $T$ is singular on a nonzero space, at least one $e_i$ has $Pe_i=0$; replacing the corresponding $f_i$ by $-f_i$ defines a different unitary map with the same product $UP$. The zero-dimensional endomorphism is invertible and has just its unique unitary map. The nonunique completion in the singular case is distinct from a canonical partial-isometry formulation.
\end{proof}

\Kentry{14.qr_cholesky}{Reduced QR and Cholesky factorizations}
\begin{theorem}
For a full-column-rank $m$-by-$n$ real or complex matrix $A$, there is a unique reduced factorization $A=QR$, where $Q$ has $n$ orthonormal columns and $R$ is $n$-by-$n$ upper triangular with positive real diagonal. For a positive-definite Hermitian $n$-by-$n$ matrix $H$, there is a unique factorization $H=R^\dagger R$ with $R$ upper triangular and positive real diagonal.
\end{theorem}
\begin{proof}
Apply Gram--Schmidt to the independent columns $a_1,\ldots,a_n$ of $A$, producing $q_j$. The initial-span property gives
$a_j=\sum_{i\le j}r_{ij}q_i$, where $r_{ij}=\ip{a_j}{q_i}$ for $i<j$, and $r_{jj}$ is the positive norm of the residual at step $j$. Let $Q$ have columns $q_i$ and set $r_{ij}=0$ for $i>j$. These formulas give $A=QR$ and the required properties. Independent $n$ columns imply $n\le m$; the case $n=0$ uses the unique empty factors.

For uniqueness, suppose $QR=Q'R'$. Both upper-triangular factors have nonzero diagonal and are invertible. Their inverses are upper triangular: solving a column of $R'x=e_j$ by backward substitution gives $x_i=0$ for $i>j$, then determines the remaining entries. The two column spaces of $Q,Q'$ equal $\im A$. Therefore $Q'=QC$ where $C=Q^\dagger Q'$, and $C$ is unitary: $C^\dagger C=Q'^\dagger QQ^\dagger Q'=Q'^\dagger Q'=I$, since $QQ^\dagger$ is the projection onto their common column space. The equation becomes $R=CR'$, so $C=R(R')^{-1}$ is upper triangular with positive real diagonal.

An upper-triangular unitary matrix is diagonal with diagonal entries of modulus one. To see this, its first column has only its first entry possibly nonzero, and norm one forces that entry to have modulus one. Orthogonality of every other column with the first forces every other entry in the first row to vanish. The remaining lower-right block is again upper triangular and unitary. Repeating on the size, down to the empty block, proves the claim. Positive real diagonal therefore forces $C=I$, giving $Q=Q'$ and $R=R'$.

For Cholesky existence, take the positive root $S=\sqrt H$. It is invertible, so its square QR factorization gives $S=QR$. Since $S=S^\dagger$, we have $H=S^\dagger S=R^\dagger Q^\dagger QR=R^\dagger R$. If $R_1^\dagger R_1=R_2^\dagger R_2$, then $C=R_1R_2^{-1}$ is upper triangular with positive real diagonal, and
$C^\dagger C=(R_2^{-1})^\dagger R_1^\dagger R_1R_2^{-1}=I$.
The preceding upper-triangular unitary argument gives $C=I$, proving uniqueness. This theorem does not assume that rank-deficient or merely semidefinite inputs have the same positive-diagonal factorization.
\end{proof}

\Kentry{14.ellipsoids}{Ellipsoids and volume change}
\begin{theorem}
Let $T$ be an invertible endomorphism of an $n$-dimensional inner-product space over $\mathbb F=\R$ or $\C$. Let $B_{\le} =\{v:\norm v\le1\}$ and $B_{<}=\{v:\norm v<1\}$. In an orthonormal output basis from the SVD,
\[
T(B_{\le})=\left\{\sum_{i=1}^n y_if_i:
                 \sum_{i=1}^n\frac{|y_i|^2}{s_i^2}\le1\right\},
\qquad
T(B_{<})=\left\{\sum_{i=1}^n y_if_i:
                 \sum_{i=1}^n\frac{|y_i|^2}{s_i^2}<1\right\}.
\]
Here every $s_i>0$. These are respectively the closed and open ellipsoid with principal directions $f_i$ and lengths $s_i$. More generally, an invertible operator maps every such open (respectively closed) ellipsoid to another open (respectively closed) ellipsoid. Over $\R$, the volume scale is $\prod_i s_i=|\det T|$.
\end{theorem}
\begin{proof}
Invertibility makes the SVD rank $r=n$. For $v=\sum_i x_i e_i$, Parseval gives $\norm v^2=\sum_i|x_i|^2$ and $Tv=\sum_i s_i x_i f_i$. Set $y_i=s_i x_i$; division is permitted because every $s_i>0$. The condition $\sum_i|x_i|^2\le1$ is exactly the first ellipsoid inequality, and the condition with $<$ is exactly the second. Conversely each output satisfying the chosen inequality has the preimage with $x_i=y_i/s_i$. This proves both equalities, not merely inclusions. Over $\R$ the endpoints of a principal semiaxis in the closed ellipsoid are $\pm s_i f_i$; over $\C$ the one-complex-dimensional direction has coordinate modulus bounded by $s_i$. The strict-inequality ellipsoid excludes its boundary, so those endpoints are not asserted to lie in it.

For a given ellipsoid in an orthonormal basis $g_i$ with positive lengths $a_i$, define the invertible positive diagonal map $Sg_i=a_i g_i$. The same coordinate calculation makes the ellipsoid $S(B_{<})$ or $S(B_{\le})$. Its image under $T$ is $(TS)(B_{<})$ or $(TS)(B_{\le})$, which has the asserted ellipsoid form by applying the preceding SVD argument to the invertible $TS$. A translated ellipsoid has its center translated by $T$ because $T(c+E)=Tc+T(E)$.

The determinant identity from \kref{14.svd} gives $|\det T|=\prod_i s_i$. In the real case Theorem~\ref{aux:volume-change} applies to either measurable ball and its image and gives the volume scale. This measure bridge is independent of the SVD argument. No real $n$-dimensional volume interpretation is being asserted for a complex $n$-space. For $n=0$, both balls and both ellipsoids consist of the unique zero vector, the products are empty products equal to one, and real zero-dimensional volume is one. No singular ellipsoid formula divides by a zero singular value.
\end{proof}

\Kentry{14.eigenvalue_frobenius_bound}{Eigenvalue squares and the Frobenius norm}
\begin{theorem}
For a complex square matrix $A$ with eigenvalues $\lambda_1,\ldots,\lambda_n$ counted with algebraic multiplicity,
\[
\sum_i|\lambda_i|^2\le\sum_{i,j}|A_{ij}|^2.
\]
Equality holds exactly when $A$ is normal, relative to the standard inner product.
\end{theorem}
\begin{proof}
Schur triangularization gives a unitary change of basis in which the matrix is upper triangular, say $B$. Its characteristic polynomial is the product of $t-B_{ii}$ by the triangular determinant formula, so its diagonal entries list the eigenvalues with algebraic multiplicity. By unitary invariance of the Frobenius norm,
\[
\sum_i|\lambda_i|^2=\sum_i|B_{ii}|^2
\le\sum_{i,j}|B_{ij}|^2=\norm A_F^2.
\]
Equality forces every off-diagonal entry of $B$ to vanish, because the omitted summands are nonnegative. Then $B$ is diagonal and hence normal, and unitary conjugation preserves normality, so $A$ is normal. Conversely, if $A$ is normal, the complex spectral theorem gives an orthonormal eigenbasis, in which the Frobenius norm is exactly the sum of the eigenvalue squares. Basis independence then gives equality for the original matrix. The assertion holds for the empty matrix, with both sums zero.
\end{proof}

\Kentry{14.positive_det_bound}{A determinant bound from positive diagonal values}
\begin{theorem}
For a positive $T$ and an orthonormal basis $e_1,\ldots,e_n$,
\[
\det T\le\prod_{i=1}^n\ip{Te_i}{e_i}.
\]
If $T$ is positive definite, equality holds exactly when this basis is an eigenbasis.
\end{theorem}
\begin{proof}
Let $S=\sqrt T$. Then $T=S^\dagger S$, and the column vectors $Se_i$ in any orthonormal codomain basis satisfy
$\norm{Se_i}^2=\ip{Te_i}{e_i}$.
Hadamard's inequality from \kref{05.hadamard} gives
$|\det S|\le\prod_i\norm{Se_i}$.
Square both sides. Determinant multiplicativity and the conjugate-transpose determinant formula give $\det T=\overline{\det S}\det S=|\det S|^2$, yielding the desired inequality.

If $T$ is positive definite, then $S$ is invertible, so none of its columns is zero. The equality clause of Hadamard is therefore exactly pairwise orthogonality of the $Se_i$. But
$\ip{Se_i}{Se_j}=\ip{Te_i}{e_j}$,
so this says that the matrix of $T$ in the chosen orthonormal basis has zero off-diagonal entries, equivalently that every $e_i$ is an eigenvector. Conversely, such a diagonal matrix gives equality by multiplying its diagonal entries. No equality characterization of this form is asserted for singular $T$, where zero columns can force equality for other reasons. The empty basis again satisfies the assertion with products equal to one.
\end{proof}

% END modules/M05.14.tex


% BEGIN modules/M05.15.tex
\chapter{Multilinear, tensor, exterior, and symmetric constructions}
Throughout this chapter $\K$ is a field. Dimension formulas have finite-dimensional inputs. Constructions by universal properties do not depend on a chosen basis. Characteristic-zero restrictions used to compare the classical source conventions are recorded separately from constructions valid in every characteristic.

\Kentry{15.multilinear}{Multilinear maps}
A function $b:V_1\times\cdots\times V_p\to W$ is multilinear if it is linear in each argument with the other arguments fixed. Pointwise addition and scalar multiplication make the collection a vector space: apply each distributive and scalar identity in the selected slot. For $p=0$ the domain is the singleton empty tuple and a map is a chosen element of $W$; this convention gives the space $W$. For $p>0$, if any argument is zero, linearity in that slot gives value zero. Expanding one slot at a time gives
\[
 b\left(\sum_{i_1}a_{1i_1}e_{1i_1},\ldots,\sum_{i_p}a_{pi_p}e_{pi_p}\right)
 =\sum_{i_1,\ldots,i_p}\left(\prod_{j=1}^p a_{ji_j}\right)b(e_{1i_1},\ldots,e_{pi_p}).
\]
This identity is obtained by $p$ applications of linearity, and fixes the coefficient convention used below.


\paragraph{Restriction, descent, and currying.}
Restrictions of a multilinear map to subspaces remain multilinear. A bilinear $b:V\times W\to H$ descends to $(V/U)\times(W/S)$ exactly when $b(U,W)=b(V,S)=0$: changing representatives adds only these terms, and necessity follows by setting a quotient argument to zero. More generally a map into $H/L$ descends when these terms lie in $L$. A bilinear map on two subspaces extends to the containing spaces by choosing complements and declaring it zero whenever either argument belongs to a chosen complement. The extension is usually not canonical. Currying gives
\[
 \operatorname{Bil}(V,W;H)\simeq\Hom(V,\Hom(W,H)),\qquad b\longmapsto(v\mapsto b(v,\cdot)).
\]
The displayed map and its inverse, evaluation, are linear; linearity in each slot proves that they take values in the stated spaces.

\paragraph{The correct factorization test.}
For bilinear $b:V\times W\to G$ and $c:V\times W\to H$, write $\widetilde b,\widetilde c$ for their tensor linearizations from \kref{15.tensor_universal}. A linear $f:G\to H$ with $c=f\circ b$ exists exactly when $\ker\widetilde b\subseteq\ker\widetilde c$. Necessity follows by composition. For sufficiency define $f_0(\widetilde b(t))=\widetilde c(t)$ on $\im\widetilde b$; the kernel condition makes this well-defined and linear. Extend it to $G$ by a complement. Given one extension, all others differ from it by a linear map factoring through $G/\im\widetilde b$. Thus the solution set is an affine space modeled on $\Hom(G/\im\widetilde b,H)$, not generally a vector space with the inherited zero. The exact uniqueness condition is that this model space vanish. Separate radical inclusions do not suffice: over $\R^2$, $b(x,y)=x_1y_1+x_2y_2$ and $c(x,y)=x_1y_1+2x_2y_2$ both have zero left and right radicals, but no scalar map $\R\to\R$ takes $b$ to $c$. This corrects Greub's problem 5(b) following Section 1.3.

\Kentry{15.tensor_universal}{The tensor product universal property}
A tensor product of $V,W$ is a space $U$ with a bilinear map $\tau:V\times W\to U$ such that for every vector space $H$ and every bilinear $b:V\times W\to H$ there is a unique linear $\widetilde b:U\to H$ with $\widetilde b\tau=b$. Write $U=V\otimes W$ and $\tau(v,w)=v\otimes w$. Uniqueness entails that elementary tensors span $U$: if their span $S$ were proper, the quotient map $U\to U/S$ and the zero map would be distinct factorizations of the same zero bilinear map. Thus equality of linear maps out of a tensor product may be tested on elementary tensors. For $p$ factors the analogous property represents $p$-linear maps; the nullary product is $\K$, with distinguished element $1$.

\Kentry{15.tensor_existence}{Existence and canonical uniqueness}
Let $F=\K^{(V\times W)}$ be the free vector space of the foundational chapter and let $R$ be spanned by
\[
 [v+v',w]-[v,w]-[v',w],\quad [av,w]-a[v,w],
\]
and the corresponding two expressions for the second slot. Set $U=F/R$ and let $v\otimes w$ be the class of $[v,w]$. The four relations say precisely that this map is bilinear. A bilinear $b$ extends uniquely to $F$ by the free-vector property and annihilates each generator of $R$, hence descends to $U$ by \kref{06.quotient_lift}. Every factorization must agree on the spanning elementary tensors, so is unique. If $U'$ has the same property, the distinguished bilinear maps give linear maps $f:U\to U'$ and $g:U'\to U$ with $gf$ and $fg$ identity on elementary tensors; they are identities everywhere. Thus the isomorphism is unique among those preserving elementary tensors. The identical free-space construction with one relation for each slot constructs any finite tensor product.

\Kentry{15.tensor_basis}{Tensor bases and dimension}
\begin{theorem}
If $(e_i)_{i\in I}$ and $(f_j)_{j\in J}$ are bases of $V,W$, then $(e_i\otimes f_j)_{(i,j)\in I\times J}$ is a basis. For finite dimensions, $\dim(V\otimes W)=\dim V\dim W$.
\end{theorem}
\begin{proof}
Expand each $v,w$ in its finite-support basis expression. Bilinearity expands $v\otimes w$ as a combination of the claimed tensors, which therefore span. Let $e^i,f^j$ be the coordinate functionals. The bilinear map $(v,w)\mapsto e^i(v)f^j(w)$ induces a linear functional on $V\otimes W$ taking $e_k\otimes f_l$ to $\delta_{ik}\delta_{jl}$. Applying it to any finite linear relation extracts its coefficient at $(i,j)$, which must be zero. Counting the product basis proves the finite formula. If either space is zero, the basis is empty and the tensor product is zero.
\end{proof}


\paragraph{Rank and decomposable tensors.}
In finite dimensions write $z=\sum C_{ij}e_i\otimes f_j$. Its associated map $V^\vee\to W$, $\phi\mapsto\sum C_{ij}\phi(e_i)f_j$, has matrix $C^t$. Consequently the least number of elementary tensors whose sum is $z$ equals $\rank C$. Indeed a sum of $r$ such tensors has image in the span of its $r$ second factors, so its rank is at most $r$. Conversely choose a basis $w_1,\ldots,w_r$ of the image; the coordinate functionals of this map on $V^\vee$ are evaluations at vectors $v_1,\ldots,v_r\in V$ by \kref{07.biddual}, yielding $z=\sum v_i\otimes w_i$. The $v_i$ are independent because the coordinate map onto $\K^r$ is surjective, and the $w_i$ were independent. Thus a minimal expression has independent factor lists, although the lists themselves need not be unique. Rank zero gives the empty sum. This argument also shows that a tensor is one elementary tensor, including zero, exactly when this rank is at most one.

If $v\otimes w=v'\otimes w'\ne0$, choose $\phi\in V^\vee$ with $\phi(v)=1$. Applying $\phi\otimes\id$ shows $w=\phi(v')w'$, so $\phi(v')\ne0$. Substitute $w'=\phi(v')^{-1}w$ and evaluate the second factor with a functional nonzero on $w$ to obtain $v'=\phi(v')v$. Conversely these reciprocal rescalings preserve the tensor. Applying this argument successively proves the corresponding criterion for nonzero elementary tensors with any positive number of factors: the corresponding factors differ by nonzero scalars whose product is one. The nonzero assumption cannot be dropped.

\Kentry{15.tensor_maps}{Tensoring linear maps}
For linear $S:V\to V'$ and $T:W\to W'$, the map $(v,w)\mapsto Sv\otimes Tw$ is bilinear, so defines the unique linear $S\otimes T$ with that formula. On each elementary tensor,
\[
 (S'\otimes T')(S\otimes T)(v\otimes w)=S'Sv\otimes T'Tw.
\]
Spanning proves $(S'\otimes T')(S\otimes T)=(S'S)\otimes(T'T)$. The same test gives $\id\otimes\id=\id$ and linearity in either map argument. In particular tensoring two isomorphisms yields an isomorphism whose inverse is the tensor of the inverses.

\Kentry{15.tensor_coordinates}{Coordinates, Kronecker matrices, and basis change}
Write $t=\sum C_{ij}e_i\otimes f_j$. If $e'_a=\sum_i e_iP_{ia}$ and $f'_b=\sum_j f_jQ_{jb}$, expansion gives $C_{\rm old}=P C_{\rm new}Q^t$, hence
\[
 C_{\rm new}=P^{-1}C_{\rm old}(Q^{-1})^t.
\]
If $Se_i=\sum_k A_{ki}u_k$ and $Tf_j=\sum_l B_{lj}v_l$, then
$(S\otimes T)(e_i\otimes f_j)=\sum_{k,l}A_{ki}B_{lj}u_k\otimes v_l$. Therefore in lexicographic product order its matrix entry with row $(k,l)$ and column $(i,j)$ is $A_{ki}B_{lj}$, the Kronecker product convention. These computations prove the transformation laws; an array without its basis and index order is not an intrinsic tensor.


\paragraph{Trace, determinant, and zero-factor guards.}
For endomorphisms $S$ of an $n$-dimensional space and $T$ of an $m$-dimensional space, the product matrix gives
\[
 \tr(S\otimes T)=\sum_{i,j}S_{ii}T_{jj}=\tr(S)\tr(T),\qquad
 \det(S\otimes T)=(\det S)^m(\det T)^n.
\]
To prove the second formula over any field, factor $S\otimes T=(S\otimes I)(I\otimes T)$. Reorder the product basis to group by the second index; $S\otimes I$ becomes block diagonal with $m$ copies of $S$. In the original grouping, $I\otimes T$ is block diagonal with $n$ copies of $T$. The determinant of a block diagonal matrix is the product of block determinants: in the Leibniz expansion all nonzero terms preserve every block, and their signs and products factor block by block. Determinant multiplicativity proves the formula; neither invertibility nor splitting of a characteristic polynomial was used. For a zero-dimensional factor, the empty determinant and every zeroth power, including $0^0$ in this algebraic convention, equal $1$, making the formula valid. The trace is zero in that case.

Under $V^\vee\otimes V\simeq\End(V)$ the map $X\mapsto AXB$ corresponds to $B^\vee\otimes A$, since it sends $\phi\otimes v$ to $(\phi\circ B)\otimes Av$. Thus its trace is $\tr A\tr B$ and its determinant is $(\det A\det B)^n$. Transposition preserves both trace and determinant, justifying removal of the dual from these formulas.

If both factor dimensions are positive, the determinant formula shows that $S\otimes T$ is invertible exactly when both $S,T$ are invertible. This converse fails for a zero-dimensional factor, because the tensor product is then the zero space and its unique endomorphism is invertible. Likewise, for arbitrary linear maps with two nonzero domain spaces, injectivity of their tensor implies injectivity of each: a nonzero kernel vector tensored with a nonzero vector in the other domain is nonzero by the tensor basis theorem but is sent to zero. If both maps are injective, tensor their left inverses to obtain a left inverse of their tensor map; this implication needs no nonzero-domain restriction. These guards correct the unqualified converses in Greub's exercises after Section 1.19.

\Kentry{15.associativity_symmetry}{Associativity, symmetry, and units}
The assignments
\[
 (u\otimes v)\otimes w\mapsto u\otimes(v\otimes w),\qquad
 v\otimes w\mapsto w\otimes v,\qquad a\otimes v\mapsto av
\]
define isomorphisms $(U\otimes V)\otimes W\simeq U\otimes(V\otimes W)$, $V\otimes W\simeq W\otimes V$, and $\K\otimes V\simeq V$. To justify the first assignment, first factor the bilinear map in $(u,v)$ for each $w$; uniqueness makes the resulting map linear in $w$, so the outer universal property factors it again. The reverse bracketing gives its inverse, since both composites fix elementary threefold tensors. The flip has itself as inverse. The inverse to the unit map is $v\mapsto1\otimes v$, because $a\otimes v=1\otimes av$. This also proves compatibility of all rebracketings: any two routes with the same effect on elementary tensors are equal.

\Kentry{15.tensor_direct_sum}{Distributing over direct sums and subspace inclusions}
Let $V=\bigoplus_{i=1}^rV_i$. The map $V\otimes W\to\bigoplus_i(V_i\otimes W)$ sends $v\otimes w$ to $(\pi_i(v)\otimes w)_i$. Its inverse sums the maps $\iota_i\otimes\id_W$. Their composites are identities by $\pi_i\iota_j=\delta_{ij}\id$ and $\sum_i\iota_i\pi_i=\id$. The same construction distributes in the other variable. If $U\subseteq V$, choose a complement using \kref{06.complement}; then $V=U\oplus U'$ and the displayed isomorphism identifies $U\otimes W\to V\otimes W$ with a direct-summand inclusion. Equivalently, tensor a left inverse of $U\hookrightarrow V$ to obtain a left inverse of its tensor map. For unrestricted spaces that complement uses the explicitly declared basis-choice principle; for finite spaces it uses finite basis extension.

\Kentry{15.tensor_quotient}{Tensor products and quotients}
\begin{theorem}
For subspaces $U\subseteq V$, $S\subseteq W$ there is a canonical isomorphism
\[
 (V/U)\otimes(W/S)\simeq(V\otimes W)/(U\otimes W+V\otimes S).
\]
\end{theorem}
\begin{proof}
The tensor of the quotient maps sends $v\otimes w$ to $[v]\otimes[w]$ and vanishes on the indicated sum, giving one direction. Conversely send $([v],[w])$ to the class of $v\otimes w$ in the right-hand quotient. Replacing $v$ by $v+u$ adds $u\otimes w$; replacing $w$ by $w+s$ adds $v\otimes s$ and possibly $u\otimes s$. All vanish in that quotient. The assignment is therefore well-defined and bilinear, hence induces the reverse map. Both composites fix the spanning classes of elementary tensors. This proves bijectivity without equating a quotient with an arbitrarily chosen complement.
\end{proof}

\Kentry{15.tensor_kernel}{Images and kernels of tensor maps}
For finite-dimensional maps $S:V\to V'$ and $T:W\to W'$, identify tensor products of subspaces through their injective natural maps. Then
\[
 \im(S\otimes T)=\im S\otimes\im T,\qquad
 \ker(S\otimes T)=\ker S\otimes W+V\otimes\ker T.
\]
Choose $V=\ker S\oplus V_0$ and $W=\ker T\oplus W_0$. The restrictions $S|_{V_0}:V_0\to\im S$ and $T|_{W_0}:W_0\to\im T$ are isomorphisms: injectivity follows from the complements and surjectivity by subtracting kernel components. Distributing the tensor products gives four summands. On the three containing a kernel factor the tensor map is zero; on $V_0\otimes W_0$ it is an isomorphism onto $\im S\otimes\im T$. This proves both assertions and covers zero maps and zero spaces.

\Kentry{15.tensor_intersections}{Intersections of tensor subspaces}
For subspaces $U,U'\subseteq V$ and $W,W'\subseteq Z$,
\[
 (U\otimes W)\cap(U'\otimes W')=(U\cap U')\otimes(W\cap W')
\]
inside $V\otimes Z$, in the finite-dimensional setting here. Choose a basis of $I=U\cap U'$, extend it separately by lists $A,A'$ to bases of $U,U'$. The combined list for $I,A,A'$ is independent: a relation between its $U$ and $U'$ portions places each portion in $I$, and independence in the two extended bases kills the $A,A'$ coefficients and then the $I$ coefficients. Extend the combined list to a basis of $V$. Do the same with $W,W'$ inside $Z$. In the product basis, $U\otimes W$ is the coordinate span of pairs drawn from the corresponding two subsets, and $U'\otimes W'$ is the span of the other pair of subsets. Uniqueness of coordinates makes the intersection the span of the intersection of these index sets, exactly the displayed tensor product.

\Kentry{15.tensor_dual_hom}{Tensors, duals, and linear maps}
For finite-dimensional $V$ and any $W$, the canonical map
\[
 \Phi:V^\vee\otimes W\longrightarrow\Hom(V,W),\quad
 \Phi(\phi\otimes w)(v)=\phi(v)w
\]
is an isomorphism. It is induced by a bilinear assignment. For a basis $e_i$ with dual $e^i$, define $\Psi(T)=\sum_i e^i\otimes Te_i$. Then $\Phi\Psi(T)(v)=\sum_i e^i(v)Te_i=Tv$. Conversely, $\Psi\Phi(\phi\otimes w)=\sum_i\phi(e_i)e^i\otimes w=\phi\otimes w$. Spanning proves the second inverse identity. The inverse is basis-independent because the canonical map $\Phi$ has a unique inverse. Finite dimension of $V$ is essential for this finite expansion; with infinite $V$ elementary finite sums produce only finite-rank maps.


\paragraph{Composition, the unit tensor, and trace pairing.}
For finite-dimensional $V$, equip $V^\vee\otimes V$ with the product
\[
 (\phi\otimes v)\circ(\psi\otimes w)=\phi(w)\psi\otimes v.
\]
The right side is multilinear, so defines a bilinear product. Under $\Phi$ it becomes composition: applied to $x$, both sides give $\phi(w)\psi(x)v$. Since $\Phi$ is bijective, associativity follows from associativity of composition. The unit is $\sum_i e^i\otimes e_i$, because it maps to $\id_V$; this characterization proves basis independence. For unrestricted $V$, the same map identifies tensors with finite-rank maps: surjectivity onto that space follows by expanding the image in a finite basis, and injectivity follows by writing a tensor with independent vector factors and evaluating its zero operator on every input to force all functional coefficients to vanish. It gives no unit for nonzero infinite-dimensional $V$: the identity has infinite rank.

For finite $V,W$, the pairing
\[
 \Hom(V,W)\times\Hom(W,V)\longrightarrow\K,
 \qquad(S,T)\longmapsto\tr(TS)
\]
 is perfect. With bases $e_i,f_a$, the map $E_{ai}:e_j\mapsto\delta_{ij}f_a$ and the map $F_{jb}:f_c\mapsto\delta_{bc}e_j$ pair by $\delta_{ij}\delta_{ab}$. These are dual bases, proving nondegeneracy, including when one space is zero. For rank-one maps the pairing is $\tr((\psi\otimes v)(\phi\otimes w))=\psi(w)\phi(v)$; this is the corresponding tensor evaluation pairing, not the positive Hermitian Hilbert--Schmidt pairing.

\paragraph{Linear maps on a tensor product and on an endomorphism space.}
For finite $V,W,V',W'$, the map
\[
 \Hom(V,V')\otimes\Hom(W,W')\longrightarrow\Hom(V\otimes W,V'\otimes W'),
 \quad S\otimes T\longmapsto S\otimes T
\]
 is an isomorphism: tensor products of the matrix units send one product-basis vector to one product-basis vector and all others to zero. They are exactly the matrix-unit basis of the target. This verifies both surjectivity and injectivity, rather than using a dimension count without a map. Evaluating elementary tensors and dual elementary tensors likewise gives $(S\otimes T)^\vee=S^\vee\otimes T^\vee$ under the perfect pairing of \kref{15.tensor_pairing}.

Put $A=\End(V)$, $\dim V=n$. The map
\[
 A\otimes A\longrightarrow\End(A),\qquad a\otimes b\longmapsto(X\mapsto aXb)
\]
 is a linear isomorphism. In matrix units,
$E_{ij}E_{ab}E_{kl}=\delta_{ja}\delta_{bk}E_{il}$, so its images are precisely all matrix units on the vector space with basis $E_{ab}$. For $n=0$ both spaces are zero. Composition of left/right multiplications has the rule
$L_aR_b\,L_cR_d=L_{ac}R_{db}$. Therefore the algebra isomorphism is with $A\otimes A^{\mathrm{op}}$, where the second product is reversed, and not with the ordinary tensor-product algebra. This last distinction is essential in using Greub's Section 1.29. The stronger theorem in Section 1.30 classifying algebra endomorphisms as inner belongs to the separately excluded central-simple-algebra extension.

\Kentry{15.tensor_pairing}{The perfect dual tensor pairing}
For finite-dimensional $V,W$ the formula
\[
 (\phi\otimes\psi)(v\otimes w)=\phi(v)\psi(w)
\]
defines a linear map $V^\vee\otimes W^\vee\to(V\otimes W)^\vee$ by two applications of the universal property. Products of dual basis vectors evaluate as the coordinate functionals on the tensor basis, so the map sends a basis to the full dual basis and is an isomorphism. A bilinear form on $V\times W$ first corresponds uniquely to a functional on $V\otimes W$, then to this dual tensor product. The same proof gives the pairing of any finite number of finite-dimensional factors. No claim is made that all bilinear forms on unrestricted infinite-dimensional spaces arise from finite-rank tensors.
For finite-dimensional $V,W$, finite biduality (\kref{07.biddual}) and this perfect pairing also give the explicit isomorphism $V\otimes W\to\operatorname{Bil}(V^\vee,W^\vee;\K)$, $v\otimes w\mapsto[(\phi,\psi)\mapsto\phi(v)\psi(w)]$; the tensor basis maps to the bilinear coordinate basis, proving bijectivity.
The canonical map $V^\vee\otimes W^\vee\to(V\otimes W)^\vee$ is nevertheless injective for arbitrary spaces. Write a tensor as $\sum_{i=1}^r\phi_i\otimes\psi_i$ with the nonzero $\phi_i$ independent, by collecting a basis of their finite-dimensional span. The evaluation map $V\to\K^r$, $v\mapsto(\phi_i(v))_i$, is onto: otherwise a nonzero functional on $\K^r$ annihilating its proper image would give a nontrivial relation among the $\phi_i$. Choose $v_j$ mapping to the $j$th coordinate vector. If the tensor evaluates to zero on all $v\otimes w$, evaluation at $v_j\otimes w$ yields $\psi_j(w)=0$ for every $w$, hence every $\psi_j=0$. This proves injectivity but does not assert surjectivity in infinite dimensions.


\Kentry{15.multilinear_dimension}{Dimension of multilinear-map spaces}
Let $V_1,\ldots,V_p$ and $W$ all be finite-dimensional, with bases $e_{ji}$ and $f_l$ respectively. A multilinear map is determined by its values on basis tuples, by the expansion at \kref{15.multilinear}. Conversely, arbitrarily prescribe vectors $w_{i_1\ldots i_p}\in W$ and use that expansion to define a map; it is linear in every slot and has the prescribed values. Taking one prescribed value equal to one $f_l$ and all others zero gives a basis of the multilinear-map space. Its dimension is $\dim W\prod_j\dim V_j$. For $p=0$ the formula reads $\dim W$, in agreement with the empty-tuple convention. If $p>0$ and one input space is zero, every value is zero and the dimension formula vanishes.

\Kentry{15.inner_tensor}{Inner products on tensor products}
For finite-dimensional real or complex inner-product spaces, choose orthonormal bases $e_i,f_j$. Define
\[
 \left\langle\sum c_{ij}e_i\otimes f_j,\sum d_{ij}e_i\otimes f_j\right\rangle
 =\sum c_{ij}\overline{d_{ij}}.
\]
This is first-slot linear, conjugate symmetric, and positive definite because the tensor coefficients are unique. Expanding $v,w,v',w'$ in the orthonormal bases gives $\ip{v\otimes w}{v'\otimes w'}=\ip v{v'}\ip w{w'}$. A sesquilinear form satisfying this rule is unique since elementary tensors span in each slot. Thus the constructed form is independent of the chosen orthonormal bases and the product rule canonically defines it. The tensor basis of any two orthonormal bases is orthonormal by the same rule.


\paragraph{General bilinear pairings and metric tensors.}
Greub uses ``inner product'' also for a nondegenerate symmetric bilinear form, which need not be positive. More generally, for bilinear forms $b:V\times V'\to\K$ and $c:W\times W'\to\K$, there is a unique bilinear pairing
\[
 (b\otimes c)(v\otimes w,v'\otimes w')=b(v,v')c(w,w').
\]
Twice apply the tensor universal property to construct it; spanning proves uniqueness. If both pairings are perfect on finite-dimensional spaces, choose bases on one side and their pairing-dual bases on the other. The product bases again pair by Kronecker deltas, proving perfection of the tensor pairing. For finite-dimensional factor spaces, its left radical is
\[
 \operatorname{rad}_L(b)\otimes W+V\otimes\operatorname{rad}_L(c),
\]
by applying \kref{15.tensor_kernel} to the two maps into dual spaces and then the injective dual tensor map. The right-radical formula is obtained by interchanging the arguments. Hence the converse nondegeneracy assertion requires the factors on the relevant side to be nonzero; a zero tensor space cannot detect degeneracy of the other factor. Tensoring symmetric forms gives a symmetric form; positivity is an additional property already proved above, not part of the general definition here.

For a nondegenerate symmetric bilinear form $b$ on finite-dimensional $V$, let $G_{ij}=b(e_i,e_j)$. The covariant and contravariant metric tensors are
\[
 g_{\rm cov}=\sum_{i,j}G_{ij}e^i\otimes e^j,\qquad
 g_{\rm con}=\sum_{i,j}(G^{-1})_{ij}e_i\otimes e_j.
\]
The first tensor evaluates to $b(v,w)$, so is independent of basis by the perfect tensor pairing. The map $b^\flat:V\to V^\vee$, $v\mapsto b(v,\cdot)$, is an isomorphism with matrix $G^t=G$. Under the tensor--Hom correspondence, the second tensor is its inverse, since contracting its first slot with $\phi$ gives the vector with coordinate column $G^{-1}[\phi]$. This proves its basis independence as well. Contracting a covariant index of $g_{\rm cov}$ against a contravariant index of $g_{\rm con}$ gives the identity tensor: its remaining coefficients are $\sum_jG_{ij}(G^{-1})_{jk}=\delta_{ik}$. Thus the two coefficient arrays are $G$ and $G^{-1}$; they coincide with the identity only in an orthonormal basis. For a real Euclidean form the displayed sums reduce to the two sums in Greub's Section 3.17. Over a complex Hermitian space the Riesz map is conjugate-linear; one must use a conjugate vector-space factor to encode its metric this way. The bilinear metric formulas above are not silently asserted for Hermitian forms.

The same construction extends a nondegenerate symmetric bilinear form to exterior powers by
\[
 b_p(v_1\wedge\cdots\wedge v_p,w_1\wedge\cdots\wedge w_p)
   =\det(b(v_i,w_j)).
\]
Indeed compose $\Lambda^p(b^\flat)$ with the perfect exterior pairing of \kref{15.exterior_pairing}; both constituent maps are isomorphisms, proving nondegeneracy. In the real positive definite or complex Hermitian case, instead define the form on increasing wedges of an orthonormal basis by making them orthonormal. Expansion and the determinant formula give the same displayed Gram determinant, with the Hermitian slot convention in the complex case. Uniqueness on spanning wedges proves basis independence. This supplies the exterior metric used by the Euclidean Pfaffian bridge.

\Kentry{15.tensor_algebra}{The tensor algebra}
Set $V^{\otimes0}=\K$ and $T(V)=\bigoplus_{p\ge0}V^{\otimes p}$. Each element has finitely many nonzero homogeneous components. Multiplication concatenates elementary tensors and extends bilinearly. Associativity holds on elementary words and therefore on all finite sums; $1\in\K$ is the unit. If $f:V\to A$ is linear into a unital associative $\K$-algebra, send $v_1\otimes\cdots\otimes v_p$ to $f(v_1)\cdots f(v_p)$ and scalars to scalar multiples of the unit. Multilinearity gives each homogeneous linear map; their direct sum is an algebra homomorphism by concatenation. Any extending homomorphism must have that formula because $T(V)$ is generated by $V$ and $1$, proving uniqueness.

\Kentry{15.permutation_actions}{Permutation action with a fixed convention}
For $\sigma\in S_p$ set
\[
 P_\sigma(v_1\otimes\cdots\otimes v_p)
 =v_{\sigma^{-1}(1)}\otimes\cdots\otimes v_{\sigma^{-1}(p)}.
\]
The reordered map is multilinear, so gives a linear operator. With composition $(\sigma\tau)(i)=\sigma(\tau(i))$, the $j$th slot of $P_\sigma P_\tau$ is $v_{\tau^{-1}\sigma^{-1}(j)}=v_{(\sigma\tau)^{-1}(j)}$. Thus $P_\sigma P_\tau=P_{\sigma\tau}$ and $P_1=\id$. This inverse-index convention prevents reversal of the action law.

\Kentry{15.symmetrizers}{Averaging projectors}
Suppose $p!$ is invertible in $\K$. Define
\[
 S_p=\frac1{p!}\sum_{\sigma\in S_p}P_\sigma,\qquad
 A_p=\frac1{p!}\sum_{\sigma\in S_p}\operatorname{sgn}(\sigma)P_\sigma.
\]
Reindexing $\sigma\mapsto\tau\sigma$ gives $P_\tau S_p=S_p$ and $P_\tau A_p=\operatorname{sgn}(\tau)A_p$. On a tensor fixed by all $P_\tau$, the average $S_p$ is the identity. On a tensor satisfying $P_\tau t=\operatorname{sgn}(\tau)t$, each summand of $A_p t$ equals $t$, so $A_p t=t$. It follows that $S_p^2=S_p$ and $A_p^2=A_p$, with exactly these invariant and sign-invariant images. For $p=0,1$ both are identities. If $p!$ vanishes the displayed division is unavailable; the quotient constructions below remain valid and must not be silently replaced by these projectors.


For bilinear maps over a field with $2\ne0$, the symmetric and alternating parts are
$b_s(u,v)=(b(u,v)+b(v,u))/2$ and $b_a(u,v)=(b(u,v)-b(v,u))/2$. They sum to $b$, and a bilinear map that is both symmetric and alternating vanishes, since $2b(u,v)=0$. This proves uniqueness of the decomposition. More generally for $p\ge2$ a tensor that is simultaneously invariant and sign-invariant under all permutations is zero when $2\ne0$, by applying a transposition. This is an actual characteristic condition, not just a naming convention.

\Kentry{15.exterior}{Exterior algebra and alternating universality}
Let $J$ be the two-sided ideal of $T(V)$ generated by $v\otimes v$ for all $v\in V$. It is homogeneous. Define $\Lambda(V)=T(V)/J$, with product $\wedge$ and homogeneous parts $\Lambda^p V$. Expanding $(u+v)\wedge(u+v)=0$ gives $u\wedge v=-v\wedge u$. This identity holds also in characteristic two, but the stronger relation $v\wedge v=0$ remains imposed.

A multilinear $b:V^p\to H$ is called alternating if its value is zero whenever two inputs coincide. Substituting $u+v$ in two equal positions and expanding proves that exchanging those positions changes its sign. A linear map from $V^{\otimes p}$ associated to $b$ therefore annihilates the degree-$p$ part of $J$: generators have equal adjacent entries, and other entries may be fixed. Conversely a map annihilating that part is alternating, since any repeated entries can be moved next to one another by successive swaps. It consequently factors uniquely through $\Lambda^p V$. This proves the alternating universal property in every characteristic. The algebra quotient is associative and unital because $J$ is a two-sided ideal. A linear map $V\to A$ into an associative algebra with every image element squaring to zero extends uniquely to an algebra homomorphism $\Lambda(V)\to A$ by the tensor-algebra property and annihilation of $J$.

\Kentry{15.exterior_basis}{Exterior bases, including boundary degrees}
\begin{theorem}
For a basis $e_1,\ldots,e_n$, the increasing wedges $e_I=e_{i_1}\wedge\cdots\wedge e_{i_p}$, $i_1<\cdots<i_p$, form a basis of $\Lambda^pV$. Thus $\dim\Lambda^pV=\binom np$.
\end{theorem}
\begin{proof}
Elementary words span; reorder each word by adjacent swaps. A word with a repeated basis vector vanishes, and every other word is plus or minus an increasing word. For each increasing index set $I$, the function on $V^p$ sending $(v_1,\ldots,v_p)$ to the determinant of the $p$-by-$p$ array $(e^{i_a}(v_b))_{a,b}$ is alternating, by the determinant proof of \kref{05.det_uniqueness}. It factors through $\Lambda^pV$ and takes $e_J$ to $\delta_{IJ}$. Applying these functionals to a relation proves independence. The degree-zero element $1$ is a basis, and for $p>n$ there are no increasing words, so the space is zero. This proof includes characteristic two because it uses the defining repeated-input vanishing of determinant, not division by two.
\end{proof}

\Kentry{15.wedge_sign}{Graded signs and functorial maps}
For homogeneous $a,b$ of degrees $p,q$, move the $p$ factors of an elementary word past the $q$ factors of the other word. There are $pq$ swaps, giving $a\wedge b=(-1)^{pq}b\wedge a$; bilinearity extends this to arbitrary homogeneous elements. A linear $T:V\to W$ induces the tensor-algebra homomorphism applying $T$ to every slot. It sends $v\otimes v$ to $Tv\otimes Tv$, hence sends the defining ideal into the target ideal and descends to $\Lambda(T)$. Its degree-$p$ restriction is $\Lambda^pT$. Testing elementary wedges proves $\Lambda^p(ST)=\Lambda^pS\,\Lambda^pT$ and $\Lambda^p\id=\id$.


\paragraph{Images, kernels, quotients, and ranks.}
Let $f:V\to W$ be a linear map of finite-dimensional spaces and $U=\ker f$. Extend a basis of $U$ to a basis of $V$, thereby writing $V=U\oplus C$. The images of the chosen basis of $C$ form a basis of $\im f$. In the exterior monomial basis, every wedge containing a basis vector of $U$ maps to zero; the remaining wedges map bijectively to the increasing-wedge basis of $\Lambda(\im f)$. Hence
\[
 \im\Lambda(f)=\Lambda(\im f),\qquad
 \ker\Lambda(f)=\langle U\rangle_{\Lambda(V)},\qquad
 \rank\Lambda^p(f)=\binom{\rank f}{p}.
\]
Here $\langle U\rangle$ is the two-sided ideal generated by $U$: its monomial basis consists exactly of wedges containing at least one $U$-basis vector. Thus the kernel description is intrinsic even though the proof chose a complement. Applying it to the quotient map gives the canonical graded-algebra isomorphism
$\Lambda(V)/\langle U\rangle\simeq\Lambda(V/U)$.
Degree zero has rank one, even for a zero map, since the induced map fixes the scalar unit. For $p>\rank f$ the rank is zero. These statements hold in every characteristic. This elementary quotient compatibility is retained from Greub Section 5.24, separately from the more specialized ideal-classification results surrounding it.

\paragraph{Bivector decomposability and a corrected exercise.}
Let $z\in\Lambda^2V$ in finite dimension. The alternating form on $V^\vee$ defined by
$B_z(\phi,\psi)=\langle\phi\wedge\psi,z\rangle$
has matrix $Z$ with $Z_{ij}=z_{ij}$, $Z_{ji}=-z_{ij}$, $Z_{ii}=0$ when $z=\sum_{i<j}z_{ij}e_i\wedge e_j$. If $z=u\wedge v$, then
$B_z(\phi,\psi)=\phi(u)\psi(v)-\phi(v)\psi(u)$, whose rank is at most two. Conversely apply the alternating-form normal form of \kref{11.symplectic_blocks} to $B_z$. In the dual basis this gives $z=u_1\wedge v_1+\cdots+u_r\wedge v_r$ with all $2r$ vectors independent and $\rank Z=2r$; equality follows by evaluating both tensors on the exterior dual basis. Rank at most two therefore gives either zero or a single wedge. This proves, over every field, that a bivector is decomposable, with zero allowed, exactly when $\rank Z\le2$. If it is nonzero this is equivalent to rank two, not rank one as printed in Greub's problem 3 after Section 5.7.

If additionally $\operatorname{char}\K\ne2$, the same normal form gives
$z\wedge z=2\sum_{i<j}u_i\wedge v_i\wedge u_j\wedge v_j$.
Its nonzero summands are distinct exterior basis elements, so it vanishes exactly when $r\le1$. This proves the corresponding wedge-square criterion with its necessary characteristic restriction. In characteristic two every such square vanishes, even when $r\ge2$.


For a nondegenerate symmetric bilinear form $b$ on $V$ with $\operatorname{char}\K\ne2$, this also gives the skew-operator bridge
\[
 \Lambda^2V\longrightarrow\{T\in\End(V):b(Tx,y)=-b(x,Ty)\},\qquad
 \Phi(u\wedge v)x=b(u,x)v-b(v,x)u.
\]
The formula is bilinear and alternating in $u,v$, so induces a linear map. Its associated form is
$b(\Phi(u\wedge v)x,y)=b(u,x)b(v,y)-b(v,x)b(u,y)$,
which is alternating in $x,y$. Conversely every skew operator gives an alternating form by this identity, and every alternating form arises uniquely: use the isomorphism $b^\flat$ in each factor and the perfect second exterior pairing. Nondegeneracy of $b$ then recovers the operator uniquely from its form. Thus $\Phi$ is an isomorphism and the skew-operator space has dimension $\binom{\dim V}{2}$. This proves the general bilinear version of Greub Sections 8.1--8.2; it is not restricted to positive metrics. In characteristic two the target must instead be defined by the stronger condition $b(Tx,x)=0$ for all $x$, since skew-adjointness alone no longer implies alternation.

\Kentry{15.determinant_line}{The determinant line}
Let $\dim V=n$ and $\omega=e_1\wedge\cdots\wedge e_n$, a basis of $\Lambda^nV$. If $Te_j=\sum_iA_{ij}e_i$, expansion and reordering give
\[
 Te_1\wedge\cdots\wedge Te_n
 =\left(\sum_{\sigma\in S_n}\operatorname{sgn}(\sigma)\prod_jA_{\sigma(j),j}\right)\omega
 =(\det A)\omega.
\]
Thus $\Lambda^nT$ is scalar multiplication by the determinant already defined using alternating multilinearity. Replacing $\omega$ by $c\omega$, $c\ne0$, leaves that scalar unchanged. This proves the coordinate-independent bridge, rather than declaring the two definitions equal without a map. For $n=0$, $\Lambda^0T=\id_\K$ and its scalar is $1$.

\paragraph{Cauchy--Binet and the generalized Lagrange identity.}
Let $A\in\K^{m\times n}$, $B\in\K^{n\times r}$, and choose increasing index sets $I\subseteq\{1,\ldots,m\}$ and $J\subseteq\{1,\ldots,r\}$ of size $k$. Write $A[I,S]$ for the submatrix on the indicated ordered rows and columns. Then
\[
 \det(AB)[I,J]
 =\sum_{\substack{S\subseteq\{1,\ldots,n\}\\|S|=k}}
       \det A[I,S]\,\det B[S,J].
\]
Here $\det(AB)[I,J]$ means $\det((AB)[I,J])$. To prove the formula, expand $\Lambda^k A$ on the increasing wedge basis. In the image of $e_{s_1}\wedge\cdots\wedge e_{s_k}$, the coefficient of $e_{i_1}\wedge\cdots\wedge e_{i_k}$ is the signed sum over all ways to assign the $k$ chosen rows to the $k$ columns, precisely $\det A[I,S]$. The same statement holds for $B$. The coefficient of $e_I$ in $(\Lambda^kA)(\Lambda^kB)e_J$ is consequently the sum on the right. Exterior functoriality identifies this composite with $\Lambda^k(AB)$, proving the left side. This proof is valid over every field. If $n<k$, the intermediate exterior power is zero and the sum is empty. If $k=0$, every relevant minor is the empty determinant $1$ and the sum has one term.

Taking $X,Y\in\K^{n\times k}$, $A=X^t$ and $B=Y$ gives the generalized Lagrange identity
\[
 \det(X^tY)=\sum_{|S|=k}\det X[S,:]\,\det Y[S,:].
\]
Transpose does not change a determinant, which justifies the first minor. Over $\C$ use $A=\overline X^{\,t}$ instead, obtaining
\[
 \det(\overline X^{\,t}Y)=\sum_{|S|=k}\overline{\det X[S,:]}\,\det Y[S,:].
\]
In particular, for $Y=X$ the determinant is the sum of the squared absolute values of all $k$-by-$k$ row minors. Our first-slot-linear Gram matrix is the transpose of $\overline X^{\,t}X$, so its determinant is the same. The real Gram formula is the real specialization. These statements include $k>n$ (both sides zero) and $k=0$ (both sides one), and distinguish an algebraic transpose from a Hermitian conjugate.

\paragraph{General Laplace expansion.}
For $A\in\K^{n\times n}$ and a fixed increasing row set $I$ of size $k$, with complements also increasingly ordered,
\[
 \det A=\sum_{\substack{J\subseteq\{1,\ldots,n\}\\|J|=k}}
 (-1)^{\sum_{i\in I}i+\sum_{j\in J}j}
 \det A[I,J]\,\det A[I^c,J^c].
\]
Here indices start at $1$. Regard the rows $r_1,\ldots,r_n$ of $A$ as coordinate vectors in a fixed basis $e_1,\ldots,e_n$. Their full wedge is $\det(A)e_1\wedge\cdots\wedge e_n$, using transposition invariance of the determinant. Moving the rows in $I$ to the front, preserving both internal orders, has sign
$\epsilon_I=(-1)^{\sum_{i\in I}i-k(k+1)/2}$: the $a$th chosen row moves past $i_a-a$ unchosen rows. Expand $r_I$ and $r_{I^c}$ in increasing wedge bases. Their coefficients are the two minors by the coefficient calculation above. In their product only complementary column sets survive, and $e_J\wedge e_{J^c}=\epsilon_J e_1\wedge\cdots\wedge e_n$. Thus
$\epsilon_I\det A=\sum_J\epsilon_J\det A[I,J]\det A[I^c,J^c]$.
Multiplying by $\epsilon_I$ gives the displayed sign since $k(k+1)$ is even. This proof uses no invertibility and works in all characteristics. It includes $k=0$, $k=n$ and the empty $0$-by-$0$ matrix; singleton $I$ recovers the cofactor formula.

\Kentry{15.symmetric}{Symmetric powers and their universal properties}
Let $I$ be the homogeneous two-sided ideal of $T(V)$ generated by $u\otimes v-v\otimes u$. Define $\Sym(V)=T(V)/I$. Every two generators commute in this quotient, so all words commute by successive swaps; bilinearity proves the algebra is commutative. A linear map from $V$ to a commutative unital algebra extends to $T(V)$ and kills $I$, hence extends uniquely to $\Sym(V)$. More generally it is enough that the image elements commute with one another.

The degree-$p$ relations are generated by exchanging adjacent factors, with arbitrary fixed factors before and after them. A symmetric $p$-linear map annihilates these differences, so factors through $\Sym^p(V)$; conversely a map on this quotient gives a multilinear map invariant under adjacent swaps and hence under all permutations. This proves the degree-$p$ universal property over any field. It uses a quotient by differences, not an averaging operator, and remains valid when $p!$ is zero.


A linear $f:V\to W$ induces a unique unital graded-algebra map $\Sym(f)$ by applying $f$ to generators. On generators the composition and identity laws are immediate, and uniqueness extends them to the whole algebra. For finite dimensions choose $V=\ker f\oplus C$ and bases as in the exterior proof. The monomial basis proved next shows that the kernel of $\Sym(f)$ is the span of monomials containing a generator from $\ker f$, exactly its generated ideal, while the remaining monomials map to the basis of $\Sym(\im f)$. Consequently
\[
 \im\Sym(f)=\Sym(\im f),\quad
 \ker\Sym(f)=\langle\ker f\rangle,\quad
 \Sym(V)/\langle U\rangle\simeq\Sym(V/U).
\]
The isomorphism in the last formula is induced by the quotient map, so does not depend on the complement used in its proof. This proves the finite-dimensional quotient and kernel results in the exercises after Greub Section 9.12 over any field.

\Kentry{15.symmetric_basis}{Monomial bases}
For a basis $e_1,\ldots,e_n$ of $V$, the elements $e_1^{a_1}\cdots e_n^{a_n}$, $a_i\ge0$, $\sum_i a_i=p$, form a basis of $\Sym^p V$. Commutativity puts each elementary word into this form, so they span. The linear assignment $e_i\mapsto x_i$ into the formal polynomial algebra $\K[x_1,\ldots,x_n]$ extends by the universal property to an algebra map. It sends the claimed monomials to distinct formal monomials, which are independent by the definition of a polynomial as a finite coefficient array. They are consequently independent before mapping as well. This also constructs the inverse polynomial-algebra map $x_i\mapsto e_i$ and proves the graded isomorphism. Stars-and-bars gives dimension $\binom{n+p-1}{p}$ for $n\ge1$. For $n=0$, the only monomial is the empty product in degree zero; every positive degree vanishes.


\paragraph{Perfect symmetric-power pairing.}
Assume $p!$ is invertible in $\K$ and $V$ finite-dimensional. There is a pairing
\[
 \langle\phi_1\cdots\phi_p,v_1\cdots v_p\rangle_{\rm sym}
  =\operatorname{perm}(\phi_i(v_j))
  =\sum_{\sigma\in S_p}\prod_{i=1}^p\phi_i(v_{\sigma(i)}).
\]
The expression is multilinear and symmetric in both lists. The two symmetric universal properties therefore make it a well-defined bilinear pairing on $\Sym^p(V^\vee)\times\Sym^p(V)$. For multi-indices $a,b$ of degree $p$, it pairs $(e^1)^{a_1}\cdots(e^n)^{a_n}$ with $e_1^{b_1}\cdots e_n^{b_n}$ as zero unless $a=b$, and then as $\prod_i a_i!$: a nonzero permutation must match equal indices, and there are exactly this many such permutations. Every diagonal scalar is invertible because each $a_i!$ divides $p!$ as an integer. The pairing is therefore perfect. For $p=0$ it is scalar multiplication. All degrees are covered simultaneously in characteristic zero; without the factorial hypothesis the same formula exists but need not be perfect. For example in characteristic $\ell>0$ the one-dimensional degree-$\ell$ pairing is $\ell!=0$.

On averaged tensor representatives, ordinary tensor evaluation gives the permanent divided by $p!$: expanding two averages leaves $p!$ pairs for each relative permutation. Thus the quotient pairing used here is $p!$ times that average pairing, exactly as in the exterior determinant pairing. For $v\in V$, define $D_v(1)=0$ and
$D_v(\phi_1\cdots\phi_p)=\sum_j\phi_j(v)\phi_1\cdots\widehat{\phi_j}\cdots\phi_p$.
The formula is multilinear and invariant under permutations of the generators, so descends to each symmetric power; \kref{15.symmetric_derivation} proves its derivation rule. Under this convention multiplication by $v$ is dual to $D_v$:
\[
 \langle D_v\alpha,x\rangle_{\rm sym}=\langle\alpha,vx\rangle_{\rm sym}.
\]
To prove it on $\alpha=\phi_1\cdots\phi_p$, expand the permanent on the right according to the functional paired with $v$; its $j$th contribution is $\phi_j(v)$ times the permanent on the remaining factors, precisely the $j$th term of $D_v\alpha$. Elementary monomials span, proving the identity. Repeating the equality shows that multiplication by $v_1\cdots v_r$ is dual to $D_{v_r}\cdots D_{v_1}$. The contractions commute: double deletion in distinct slots gives equal terms in either order. This establishes the pairing and operator bridges in Greub Sections 4.16, 9.4, and 9.6 without transferring an unrecorded factorial.

\Kentry{15.polynomial_bridge}{Polynomial algebras versus polynomial functions}
Choose a finite basis of $V$ and its coordinate functionals $x_i\in V^\vee$. The previous construction identifies $\Sym(V^\vee)$ with the formal polynomial algebra in $x_1,\ldots,x_n$. Evaluation at vectors is an algebra map into functions $V\to\K$.

If $\K$ is infinite this map is injective. For one variable, a nonzero polynomial of degree $d$ has at most $d$ roots: divide by $x-a$ at each root and induct on degree. For $n$ variables, write $f=\sum_j f_j(x_1,\ldots,x_{n-1})x_n^j$. If its function is zero, fixing the first $n-1$ coordinates gives a univariate polynomial vanishing at every scalar; hence each coefficient function $f_j$ is zero. Induction on the number of variables makes every $f_j$ the zero formal polynomial. For $n=0$ evaluation of constants is already injective. Over a finite field with $q$ elements, $x^q-x$ is a nonzero formal polynomial with zero function: every nonzero scalar belongs to the multiplicative group of size $q-1$, and the group identity $a^{q-1}=1$ follows, for example, by multiplying all nonzero elements and cancelling after the permutation $u\mapsto au$. Zero also satisfies the equation. Thus the function identification cannot be used without the infinite-field condition.


The formal polynomial algebra has no zero divisors. Choose lexicographic order on the finite multi-indices appearing in two nonzero polynomials. If $a$ and $b$ are their largest indices, the largest index in the product is $a+b$, and its coefficient is the product of the two nonzero leading coefficients: replacing either summand by a strictly smaller index makes the sum strictly smaller. Hence the product is nonzero over a field. In zero variables this is just the field axiom. Through the already proved polynomial isomorphism, $\Sym(V)$ is therefore an integral domain for finite-dimensional $V$, in every characteristic. This is the no-zero-divisors result of Greub Section 9.7; no claim about injectivity of polynomial-function evaluation over finite fields is needed for it.

\Kentry{15.normalization}{Quotients, averaging, and divided powers}
Assume characteristic zero. Let $q_p:V^{\otimes p}\to\Sym^pV$ be the symmetric quotient. Since $q_pP_\sigma=q_p$, one has $q_pS_p=q_p$. The map $S_p$ annihilates permutation differences, so descends to $s_p:\Sym^pV\to V^{\otimes p}$ with $s_pq_p=S_p$. These identities imply $q_ps_p=\id$ and identify $\Sym^pV$ with $\im S_p$. For exterior powers, the alternating quotient satisfies $q_pP_\sigma=\operatorname{sgn}(\sigma)q_p$ and $q_pA_p=q_p$. The averaging $A_p$ kills tensors with two equal slots, by pairing each permutation with a transposition of those slots. It therefore descends to $a_p$ with $a_pq_p=A_p$ and $q_pa_p=\id$. This proves the isomorphism with $\im A_p$ rather than merely comparing dimensions.

Under either isomorphism the product of invariant tensors of degrees $p,q$ is $S_{p+q}(a\otimes b)$ or $A_{p+q}(a\otimes b)$ respectively. Indeed applying the quotient gives the ordinary quotient product, and the displayed inverse gives exactly that average. Its associativity follows because the quotient multiplication is associative and the map is bijective.

We use $u^k$ for an ordinary algebra power and $u^{[k]}=u^k/k!$ for a divided power. Greub's power notation in the indicated characteristic-zero sections denotes the latter. Therefore
\[
 u^{[r]}u^{[s]}=\binom{r+s}{r}u^{[r+s]},\qquad
 (u+v)^{[k]}=\sum_{r+s=k}u^{[r]}v^{[s]}
\]
when $u,v$ commute. Both follow by the ordinary binomial theorem, whose coefficients count choices of slots, and division by the indicated factorials. In an exterior algebra arbitrary odd-degree elements do not commute, so that binomial assertion is not transferred without a commuting hypothesis. These formulas and the explicit quotient--averaging maps fix the normalization used throughout this manuscript.

\Kentry{15.contraction}{Contraction of mixed tensors}
For $p,q\ge1$, choose one dual slot $r$ and one vector slot $s$ in $(V^\vee)^{\otimes p}\otimes V^{\otimes q}$. On an elementary tensor define
\[
 C_{r,s}(\phi_1\otimes\cdots\otimes\phi_p\otimes v_1\otimes\cdots\otimes v_q)
 =\phi_r(v_s)\,
 \phi_1\otimes\cdots\widehat{\phi_r}\cdots\otimes\phi_p
 \otimes v_1\otimes\cdots\widehat{v_s}\cdots\otimes v_q.
\]
Here a hat removes the indicated factor; when a side has no factors it is the scalar unit. The expression is linear in each input, including the paired ones because evaluation is bilinear, so the tensor universal property gives a unique linear map of bidegree $(-1,-1)$. No basis occurs in its construction. In dual product bases, $e^i(e_j)=\delta_{ij}$ shows that the coordinates are obtained by setting the contracted indices equal and summing them. In degree $(1,1)$ the tensor--Hom identification sends $\phi\otimes v$ to $x\mapsto\phi(x)v$, whose trace in a basis is $\sum_i\phi(e_i)e^i(v)=\phi(v)$. Thus contraction is trace in that case; the coordinate sum is a consequence of evaluation, not a separate definition.


\paragraph{Equivariance of contraction.}
For $g\in\operatorname{GL}(V)$ act on a mixed elementary tensor by replacing each $v$ by $gv$ and each $\phi$ by $\phi\circ g^{-1}$. This defines an automorphism because each replacement is linear; composition gives an action. Evaluation satisfies $(\phi\circ g^{-1})(gv)=\phi(v)$. Substituting this identity in the displayed contraction formula proves that contraction commutes with this action. Thus contracted coordinates obey the tensor transformation law automatically. It is the invariant evaluation pairing that proves this fact, rather than an assumption about coordinate arrays.

\Kentry{15.tensor_derivation}{Derivations induced by an endomorphism}
For $A\in\End(V)$ define $D_A(1)=0$ and
\[
 D_A(v_1\otimes\cdots\otimes v_p)
 =\sum_{j=1}^p v_1\otimes\cdots\otimes Av_j\otimes\cdots\otimes v_p.
\]
Each sum is multilinear and induces a linear map on the degree-$p$ tensor power. Splitting the summation at the boundary of two concatenated words proves $D_A(ab)=D_A(a)b+aD_A(b)$. It extends linearly to their finite sums and is degree preserving. Any derivation extending $A$ has this formula by repeated application of the Leibniz rule; also $D(1)=D(1\cdot1)=2D(1)$ forces $D(1)=0$. This proves uniqueness. The tensor algebra homomorphism $T(A)$ instead applies $A$ in every slot simultaneously and sends $1$ to $1$; it is a different construction.


The assignment $A\mapsto D_A$ is linear, by its slot formula. It also preserves the Lie bracket:
\[
 [D_A,D_B]=D_{AB-BA}.
\]
Expanding $D_AD_B(xy)-D_BD_A(xy)$ by Leibniz cancels the two mixed terms, so the commutator is a derivation. Its value on each generator $v$ is $(AB-BA)v$; uniqueness of the extending derivation proves the identity, in every characteristic. Equivalently the slot expansion cancels all terms acting on distinct slots and leaves the same-slot commutator. The latter argument shows simultaneously that this identity holds for the descended exterior and symmetric derivations below. Under the tensor evaluation pairing, $D_{A^\vee}$ is dual to $D_A$: evaluating the two slot sums gives term by term $\phi(Av)=(A^\vee\phi)(v)$. This also proves the corresponding exterior and symmetric duality identities using their determinant and permanent pairings.

\Kentry{15.exterior_pairing}{Perfect pairings of exterior powers}
For finite-dimensional $V$, define on elementary wedges
\[
 \langle\phi_1\wedge\cdots\wedge\phi_p,
          v_1\wedge\cdots\wedge v_p\rangle
 =\det(\phi_i(v_j))_{i,j=1}^p.
\]
The determinant is alternating in the list of functionals and in the list of vectors, so the two alternating universal properties produce a well-defined bilinear pairing. The increasing wedges in a dual basis and a primal basis pair by $\delta_{IJ}$, hence the pairing is perfect and identifies $\Lambda^p(V^\vee)$ with $(\Lambda^pV)^\vee$. For $T:V\to W$, evaluating elementary wedges gives
\[
 \langle\Lambda^p(T^\vee)\alpha,x\rangle
 =\langle\alpha,\Lambda^p(T)x\rangle,
\]
since both sides are $\det(\phi_i(Tv_j))$. Spanning proves the identity for all $\alpha,x$. The same proof works over every field; it supplies explicitly the extension beyond the characteristic-zero source route.

For the averaging model assume additionally that $p!$ is invertible in $\K$, as it is in characteristic zero. There is then a normalization to preserve: the ordinary tensor evaluation of $A_p(\phi_1\otimes\cdots\otimes\phi_p)$ on $A_p(v_1\otimes\cdots\otimes v_p)$ equals $\det(\phi_i(v_j))/p!$. To see this, expand the two averages. For each relative permutation there are $p!$ pairs of permutations, each with the corresponding determinant sign, while the denominator is $(p!)^2$. Thus the determinant pairing above is $p!$ times that inherited averaged-tensor pairing. This explains why unrecorded averaging identifications can lose factorials. The quotient pairing itself, proved in the preceding paragraph, does not require this extra invertibility assumption.


\paragraph{The exterior algebra of alternating forms.}
The determinant pairing identifies $\Lambda^p(V^\vee)$ with alternating $p$-linear scalar forms by
$\alpha\mapsto[(v_1,\ldots,v_p)\mapsto\langle\alpha,v_1\wedge\cdots\wedge v_p\rangle]$.
For forms $B,C$ of degrees $p,q$, put
\[
 (B\wedge C)(v_1,\ldots,v_{p+q})
 =\sum_{\sigma\in\operatorname{Sh}(p,q)}\operatorname{sgn}(\sigma)
 B(v_{\sigma(1)},\ldots,v_{\sigma(p)})
 C(v_{\sigma(p+1)},\ldots,v_{\sigma(p+q)}),
\]
where a shuffle preserves the order within the first $p$ and last $q$ chosen indices. For decomposable exterior covectors, expand the two determinants. Every permutation of $p+q$ indices factors uniquely as a shuffle followed by an internal permutation in each block, with the product of their signs. The sum is therefore the determinant defining the image of the exterior product. By spanning, this proves compatibility with multiplication for every pair of forms; it also proves associativity and the graded sign law by transporting them from $\Lambda(V^\vee)$. This shuffle formula works in every characteristic. In characteristic zero it equals $\frac{(p+q)!}{p!q!}\operatorname{Alt}(B\otimes C)$, whereas the image of a covector tensor is $p!$ times its averaged alternating form. These are distinct normalizations. They give the precise bridge to Greub Sections 5.30 and 5.33. Expanding the determinant in its first column proves that the interior operator associated with $v\in V$ corresponds to ordinary first-slot insertion $B\mapsto B(v,\cdot,\ldots,\cdot)$.

\Kentry{15.exterior_direct_sum}{Exterior algebra of a direct sum}
On $\Lambda(V)\otimes\Lambda(W)$ put the graded multiplication
\[
 (a\otimes b)(c\otimes d)=(-1)^{\deg b\deg c}(a\wedge c)\otimes(b\wedge d)
\]
for homogeneous factors, extended bilinearly. The two ways of multiplying three such tensors have sign exponents $|b||c|+(|b|+|d|)|e|$ and $|d||e|+|b|(|c|+|e|)$, which agree, so it is associative. Send $(v,w)\in V\oplus W$ to $v\otimes1+1\otimes w$. Its square is zero: the two pure squares vanish and the cross terms cancel by the displayed sign. The exterior universal property gives an algebra map from $\Lambda(V\oplus W)$. In the reverse direction send $a\otimes b$ to $\Lambda(\iota_V)(a)\wedge\Lambda(\iota_W)(b)$. Graded commutativity makes this multiplicative with precisely the same sign. Each composite fixes the generators and unit, hence is the identity. Taking degree $p$ gives
\[
 \Lambda^p(V\oplus W)\simeq\bigoplus_{i+j=p}\Lambda^iV\otimes\Lambda^jW.
\]
The proof works in all characteristics, though its characteristic-zero instance is the source convention. Omitting the graded sign would generally destroy multiplicativity.

\Kentry{15.interior_product}{Interior product and antiderivation identities}
For $\phi\in V^\vee$, define on tensor words the signed deletion operator
\[
 \iota_\phi(v_1\cdots v_p)=\sum_{j=1}^p(-1)^{j-1}\phi(v_j)
 v_1\cdots\widehat{v_j}\cdots v_p,\qquad\iota_\phi(1)=0.
\]
Multilinearity defines it on tensor powers. Splitting its sum at a concatenation gives $\iota_\phi(ab)=\iota_\phi(a)b+(-1)^{\deg a}a\iota_\phi(b)$ for homogeneous $a$. For a generator $vv$ of the exterior ideal its value is $\phi(v)v-v\phi(v)=0$. Applying the rule to $a(vv)b$ shows its image remains in the ideal, so it descends to $\Lambda(V)$ with the same formula and rule.

In a double deletion of different positions $i<j$, deleting $i$ first shifts the position $j$ down by one. The two orders consequently have opposite signs and identical scalar products. Pairing these terms proves $\iota_\phi^2=0$. In the expansion of $\iota_\phi\iota_\psi+\iota_\psi\iota_\phi$, pair a term deleting $i$ with $\phi$ and $j$ with $\psi$ with the term deleting those same slots in the reverse order. Again signs are opposite. Thus the two operators anticommute. Taking a first factor $v$ in the signed Leibniz rule gives
\[
 \iota_\phi(v\wedge a)+v\wedge\iota_\phi(a)=\phi(v)a.
\]
These identities prove the registered characteristic-zero assertions (and in fact their alternating-quotient versions in every characteristic). Contraction by a higher exterior element is a composition of these operators and is not generally itself a degree-minus-one antiderivation.


For $\alpha=\phi_1\wedge\cdots\wedge\phi_r$, define $\iota_\alpha=\iota_{\phi_r}\cdots\iota_{\phi_1}$. This depends only on $\alpha$: the composite is multilinear in the $\phi_i$, and its value vanishes when two coincide, by anticommutation and $\iota_\phi^2=0$. Thus it factors through the alternating quotient. It is dual to left multiplication by $\alpha$ under the determinant pairing, as expansion by a first row proves for $r=1$, and composition proves for general $r$. It satisfies $\iota_{\alpha\wedge\beta}=\iota_\beta\iota_\alpha$. Higher contractions do not satisfy the degree-minus-one antiderivation law in general: for dual pairs $e_1,e_2$ and $e^1,e^2$, $\iota_{e^1\wedge e^2}(e_1\wedge e_2)=1$ but its values on $e_1,e_2$ both vanish. This gives a concrete guard for the exercise after Greub Section 5.14, without the false assertion that every higher contraction, including the zero operator, fails every possible derivation identity.

\paragraph{Contraction detects coefficients and counts degree.}
For a finite basis $e_1,\ldots,e_n$ with dual basis $e^1,\ldots,e^n$, write $L_v(u)=v\wedge u$. On $\Lambda^pV$ one has
\[
 \sum_j L_{e_j}\iota_{e^j}=pI,
 \qquad \sum_j\iota_{e^j}L_{e_j}=(n-p)I.
\]
On an increasing basis wedge, the first sum returns that wedge once for each of its $p$ indices: deletion and reinsertion have the same sign. The second identity follows by summing $\iota_{e^j}L_{e_j}+L_{e_j}\iota_{e^j}=I$, already proved above. These formulas hold in every characteristic as identities with integer scalars; no division by $p$ is implicit. The same slot calculation gives $D_A=\sum_jL_{Ae_j}\iota_{e^j}$ for the induced derivation of an endomorphism $A$.

If $1\le k\le p$ and every $k$-fold contraction of $u\in\Lambda^pV$ vanishes, then $u=0$, even if the characteristic divides $p$. Indeed, choose any increasing index set $I$ of size $p$ and a subset $J\subseteq I$ of size $k$. In the contraction by the dual wedge on $J$, the coefficient of the complementary wedge on $I\setminus J$ is, up to its invertible sign, exactly the coefficient of $I$ in $u$: no other index set can have that same deletion and complement. Every coefficient therefore vanishes. This argument explains why the degree-counting identity alone must not be divided by $p$ in arbitrary characteristic.

\Kentry{15.exterior_derivation}{Induced exterior derivations and trace}
The tensor derivation $D_A$ sends $v\otimes v$ to $Av\otimes v+v\otimes Av$, an element of the exterior ideal because it is the difference between the squares of $v+Av$, $v$, and $Av$. The Leibniz rule then makes the entire ideal invariant, so $D_A$ descends to a degree-zero derivation of $\Lambda(V)$. Its formula is the sum of applying $A$ to one wedge factor at a time. The algebra is generated by $V$, so that rule gives uniqueness.

In dimension $n$, apply $D_A$ to $e_1\wedge\cdots\wedge e_n$. In the summand changing the $j$th factor, every off-diagonal component of $Ae_j$ repeats another basis vector and vanishes; its diagonal component contributes $A_{jj}$ times the top wedge. The total coefficient is $\sum_jA_{jj}=\tr A$. For $n=0$ it is zero on the unit. In contrast, $\Lambda^n(A)$ acts by $\det A$ and preserves the unit in degree zero. This establishes the trace/determinant distinction directly.
The induced derivations are natural for intertwining maps: if $f:V\to W$ satisfies $Bf=fA$, then $D_B\Lambda(f)=\Lambda(f)D_A$. On an elementary wedge, the two sides are the sums obtained by applying $Bf=fA$ in each selected slot; spanning proves equality. The identical slot argument applies to tensor and symmetric powers. This supplies the elementary naturality used from Greub's more general discussion of twisted antiderivations, without importing that full theory.


\Kentry{15.top_pairing}{Top-degree duality and the volume choice}
For $\dim V=n$ and $0\le p\le n$, multiplication gives
\[
 \Lambda^pV\times\Lambda^{n-p}V\longrightarrow\Lambda^nV.
\]
Fix $\omega=e_1\wedge\cdots\wedge e_n$. For increasing subsets $I,J$ of sizes $p,n-p$, $e_I\wedge e_J$ is zero unless $J=I^c$; in that case it is $\epsilon(I,I^c)\omega$, with $\epsilon=\pm1$ the reordering sign. The matrix of the resulting scalar pairing is therefore a permutation matrix with nonzero signed entries. It is invertible, proving perfection. In particular it identifies $\Lambda^{n-p}V$ with $(\Lambda^pV)^\vee$ after choosing $\omega$. Replacing $\omega$ by $c\omega$ divides the scalar pairing by $c$; without that choice the canonical target is the one-dimensional determinant line, not a preferred copy of $\K$. Degree zero and degree $n$ follow from the same complementary-subset argument.


The volume-dependent identifications have an explicit naturality and sign convention. Let $\omega^*$ be dual to $\omega$ and define $D_\omega(u)=\iota_u\omega^*$ and $E_\omega(\alpha)=\iota_\alpha\omega$. In complementary basis wedges,
\[
 D_\omega(e_I)=\epsilon(I,I^c)e^{I^c},\qquad
 E_\omega(e^J)=\epsilon(J,J^c)e_{J^c}.
\]
Both equations follow by applying the successive deletion formula for interior multiplication. Hence on degree $p$, $E_\omega D_\omega=(-1)^{p(n-p)}\id$, and the analogous reverse composition has the same sign. The case $p=0$ gives $D_\omega(1)=\omega^*$ and $E_\omega(1)=\omega$. If $g:V\to W$ is an isomorphism and $\eta=\Lambda^n(g)\omega$, then
\[
 D_\eta\Lambda^p(g)=\Lambda^{n-p}((g^{-1})^\vee)D_\omega.
\]
Pair either side with an arbitrary $(n-p)$-wedge $y$ in $W$: both are the coefficient of $\Lambda^p(g)u\wedge y$ relative to $\eta$, by functoriality of the wedge product. Perfection proves equality. These formulas supply the ordinary volume-choice Poincare bridge in Greub Sections 6.8--6.9; they do not assert that the additional mixed-exterior operator of Sections 6.10--6.11 has been developed here.

\Kentry{15.charcoefficients}{Characteristic coefficients as exterior traces}
\begin{theorem}
For $T\in\End(V)$, $\dim V=n$,
\[
 \det(tI-T)=\sum_{p=0}^n(-1)^p\tr(\Lambda^pT)t^{n-p}.
\]
\end{theorem}
\begin{proof}
Let $A$ be a matrix of $T$. Expanding $Te_{j_1}\wedge\cdots\wedge Te_{j_p}$ shows that the coefficient of $e_{i_1}\wedge\cdots\wedge e_{i_p}$ is the minor $\det A[I,J]$: all injections onto $I$ occur with the sign of their ordering. Thus $\tr(\Lambda^pT)=\sum_{|I|=p}\det A[I,I]$.

Expand $\det(tI-A)$ multilinearly in its columns, choosing either $te_j$ or the negative $A$-column in each column $j$. If the $A$-columns are chosen in a subset $I$ of size $p$, every nonzero determinant term fixes each complementary index, so the remaining permutation has precisely the determinant sign of its restriction to the ordered set $I$. The term is $(-1)^pt^{n-p}\det A[I,I]$. Summing all subsets and using the trace identity proves the formula. For $p=0$, $\Lambda^0T=\id_\K$ has trace $1$; for $n=0$ the formula reads $1=1$. No splitting or eigenvalue assumption is used, so this proof also proves the identity beyond characteristic zero.
\end{proof}

For endomorphisms $A,B$ of the same finite-dimensional space, the characteristic polynomials of $AB$ and $BA$ agree even when either map is singular. Indeed, for each $p$, exterior functoriality gives $\Lambda^p(AB)=(\Lambda^pA)(\Lambda^pB)$ and the analogous reversed product. For any two square matrices $X,Y$ of the same size, $\tr(XY)=\sum_{i,j}X_{ij}Y_{ji}=\sum_{j,i}Y_{ji}X_{ij}=\tr(YX)$. Thus all the exterior traces, hence all characteristic coefficients just computed, agree. This argument includes degree zero, zero-dimensional spaces, and every characteristic; it requires neither cancellation by a determinant nor density of invertible matrices.

\Kentry{15.symmetric_direct_sum}{Symmetric algebra of a direct sum}
The ordinary tensor product of commutative algebras $\Sym(V)\otimes\Sym(W)$ is commutative under $(a\otimes b)(c\otimes d)=ac\otimes bd$. The linear map $(v,w)\mapsto v\otimes1+1\otimes w$ induces an algebra homomorphism $\Sym(V\oplus W)\to\Sym(V)\otimes\Sym(W)$. Conversely multiply the images of the two inclusion-induced maps to send $a\otimes b$ to $\Sym(\iota_V)(a)\Sym(\iota_W)(b)$. This is bilinear and multiplicative because the target is commutative. Both composites fix all generators and the unit and hence all monomials and finite sums. They are inverse graded isomorphisms. Their degree-$p$ parts give $\Sym^p(V\oplus W)\simeq\bigoplus_{i+j=p}\Sym^iV\otimes\Sym^jW$. No graded minus sign belongs in this multiplication.

\Kentry{15.symmetric_derivation}{Symmetric derivations and contractions}
For $\phi\in V^\vee$, let $D_\phi(1)=0$ and
\[
 D_\phi(v_1\cdots v_p)=\sum_{j=1}^p\phi(v_j)v_1\cdots\widehat{v_j}\cdots v_p.
\]
This is symmetric and multilinear in the generators, so factors through the symmetric quotient. Splitting the sum proves the ordinary Leibniz rule. Its values on generators determine it uniquely by that rule; it lowers degree by one. For an endomorphism $A$ the analogous sum replacing $v_j$ by $Av_j$ is degree preserving and a derivation. It also follows by descending the tensor derivation: the derivative of $uv-vu$ is $(Au)v-v(Au)+u(Av)-(Av)u$, which lies in the commutator ideal.

For ordinary powers, $D_\phi(v^p)=p\phi(v)v^{p-1}$. In characteristic zero, division by $p!$ gives $D_\phi(v^{[p]})=\phi(v)v^{[p-1]}$ for $p\ge1$. This is the factorial translation required for the source divided-power convention; the two displayed formulas are not conflicting derivative rules.

\Kentry{15.hilbert_series}{Dimension generating functions}
Write formal power series with integer coefficients. The exterior basis selects each of $n$ basis vectors either zero or one times, so expanding $\prod_{i=1}^n(1+t)$ counts degree-$p$ increasing wedges. Hence $\sum_p\dim\Lambda^pV\,t^p=(1+t)^n$. A symmetric monomial chooses a nonnegative exponent independently for each basis vector, giving
\[
 \sum_{p\ge0}\dim\Sym^pV\,t^p
 =\prod_{i=1}^n\left(\sum_{a\ge0}t^a\right)=(1-t)^{-n}.
\]
Each coefficient of the product involves only finitely many exponent tuples, so formal multiplication is defined; $\sum_{a\ge0}t^a$ is the formal inverse of $1-t$ by coefficient cancellation. For $n=0$ both products are $1$. No analytic convergence or value of $t$ is assumed.


For the tensor algebra, $\dim V^{\otimes p}=n^p$ by the product basis, so its formal dimension series is $\sum_{p\ge0}n^pt^p=(1-nt)^{-1}$. More generally, if $V=\bigoplus_{j=1}^rV_j$ with positive weights $d_j$ and $\dim V_j=n_j$, the monomial bases give weighted series $\prod_j(1+t^{d_j})^{n_j}$ for the exterior algebra and $\prod_j(1-t^{d_j})^{-n_j}$ for the symmetric algebra. Tensor products of positively graded finite-dimensional pieces multiply their dimension series: the degree-$k$ piece is $\bigoplus_{a+b=k}U_a\otimes W_b$, whose dimension is the Cauchy-product coefficient. Positive weights ensure that each coefficient involves only finitely many choices. This explains the graded specializations of Greub Sections 2.5--2.6, 5.21, and 9.9--9.11.

\Kentry{15.polynomial_polarization}{Symmetric forms and homogeneous polynomials}
Assume characteristic zero and finite-dimensional $V$. For a symmetric $p$-linear form $B$, set $q(v)=B(v,\ldots,v)$. Expansion in coordinates makes $q$ a homogeneous polynomial of degree $p$. The inverse formula is
\[
 B(v_1,\ldots,v_p)=\frac1{p!}\sum_{S\subseteq\{1,\ldots,p\}}
 (-1)^{p-|S|}q\left(\sum_{i\in S}v_i\right).
\]
Indeed expansion of a term $q(\sum_{i\in S}v_i)$ is a sum over ordered $p$-tuples of indices drawn from $S$. If a tuple uses precisely the set $J$, its total coefficient in the subset sum is $\sum_{S\supseteq J}(-1)^{p-|S|}$, which is zero unless $J$ is the whole index set: choose an unused index and pair sets with and without it. If $J$ is the whole set, each index occurs once, so the surviving $p!$ terms are all $B(v_1,\ldots,v_p)$. This proves the formula and uniqueness.

Conversely a coordinate monomial $\phi_1(v)\cdots\phi_p(v)$ is the diagonal of the symmetric multilinear form
\[
 \frac1{p!}\sum_{\sigma\in S_p}\prod_{j=1}^p\phi_j(v_{\sigma(j)}).
\]
Linearly combine these forms for the monomials of an arbitrary homogeneous polynomial; this gives existence. Polynomial functions determine their formal polynomials because characteristic zero implies an infinite field, by \kref{15.polynomial_bridge}. For $p=0$, $B$ and $q$ are constants and the formula has one term, so the convention still works.

Homogeneity alone is insufficient. Over $\R^2$, set $q(x,y)=x^4/(x^2+y^2)$ away from the origin and $q(0,0)=0$. It satisfies $q(tv)=t^2q(v)$ for every real $t$. If it were a quadratic polynomial, its values at $(1,0)$ and $(0,1)$ would force coefficients of $x^2,y^2$ to be $1,0$; its equal values $1/2$ at $(1,1)$ and $(1,-1)$ would force the cross coefficient to be zero and simultaneously force those values to be $1$, a contradiction. The equivalence therefore requires a polynomial hypothesis, exactly as recorded.

\paragraph{Symmetric forms: two products and a necessary correction.}
In characteristic zero let $R_p:\Sym^p(V^\vee)\to\operatorname{SymMult}^p(V;\K)$ be the permanent map
\[
 R_p(\phi_1\cdots\phi_p)(v_1,\ldots,v_p)
 =\sum_{\sigma\in S_p}\prod_i\phi_i(v_{\sigma(i)}).
\]
Perfection of the symmetric pairing proves that this is an isomorphism. Give symmetric forms the shuffle product, the same shuffle sum as for alternating forms but without signs. Expanding on products of covectors shows that $R_{p+q}(uv)=R_p(u)\mathbin{\star}R_q(v)$: shuffling chooses the arguments for each factor, and the internal permutations order the chosen arguments, producing each permutation exactly once. Thus $R$ is a graded-algebra isomorphism. The shuffle product is $\frac{(p+q)!}{p!q!}\operatorname{Sym}(B\otimes C)$, the product printed in Greub Section 9.15. Its raw diagonal obeys
\[
 (B\mathbin{\star}C)(v,\ldots,v)=\binom{p+q}{p}B(v,\ldots,v)C(v,\ldots,v).
\]
Accordingly the algebra map from this convention to polynomial functions is the \emph{normalized} diagonal $B\mapsto B(v,\ldots,v)/p!$. The raw diagonal is instead multiplicative for the averaged product $\operatorname{Sym}(B\otimes C)$, under the isomorphisms $R_p/p!$. These alternatives cannot be combined: for two linear forms $B,C$ the shuffle diagonal is $2B(v)C(v)$. This directly corrects the incompatible raw-diagonal homomorphism claim in Greub Section 9.12 when read with his Section 9.15 product. It does not change the preceding polarization statement, which used the raw diagonal as a vector-space identification degree by degree.

Under $R_p$, the contraction $D_v$ corresponds to ordinary insertion in the first slot, because expanding the permanent according to the factor evaluated at $v$ gives $R_{p-1}(D_v u)=R_p(u)(v,\cdot,\ldots,\cdot)$. Under $R_p/p!$, the corresponding formula has the additional factor $p$. This explains the substitution-operator convention in Greub Section 9.16 and prevents a factorial error in passing from forms to polynomial derivatives. Degree zero in all these identifications is the scalar identity.

% END modules/M05.15.tex

\backmatter

% BEGIN source-register.tex
\chapter{Source register and provenance}

The frozen source copies below are reading sources. The manuscript is organized by mathematical content, not by reproducing their chapter text. A printed page and a physical PDF page are different locators. The entry headings give section starts; the source-led item index in the companion knowledge directory supplies finer numbered-item routes.

\begin{description}

\item[F04] Kenneth Hoffman, Ray Kunze. \emph{Linear Algebra}. Second edition, Prentice-Hall, 1971.

The registered PDF contains 416 physical pages.

\item[F05] Sheldon Axler. \emph{Linear Algebra Done Right}. Fourth edition; corrected author-hosted PDF.

The registered PDF contains 404 physical pages.

\item[F06] Werner Greub. \emph{Multilinear Algebra}. Second edition, Springer, 1978.

The registered PDF contains 303 physical pages.

\end{description}

\section*{What has been unified}

The same definitions of kernel, image, algebraic dual, tensor universal property, and determinant are used throughout. Chosen-coordinate identifications are given by actual maps. Quotient constructions of exterior and symmetric powers are valid over general fields; characteristic-zero averaging models are related to them by explicit inverse maps. Source hypotheses are retained or a wider statement is proved here.

In particular, semisimple and diagonalizable decompositions remain distinct; normal and self-adjoint operators are not conflated; positive-root uniqueness does not apply to arbitrary roots; and singular equality cases are separated from their invertible versions. The source-specific Pfaffian sign and volume-vector normalization and the divided-power factorial convention are made explicit.

\section*{Preserved companion data}

The companion knowledge directory retains the 248-entry original catalogue, the previous 199-entry snapshot, per-page search text, source hashes, section and numbered-item maps, and the original unresolved-work inventory. That earlier inventory is a historical snapshot. New stable entries from the final source-led review are recorded in the manuscript additions file; refinements of existing entries retain their IDs. The current manuscript has a separate coverage and review record. Original textbook PDFs have not been overwritten.

% END source-register.tex


% BEGIN verification-register.tex
\chapter{Coverage and mathematical review}
The manuscript contains 249 stable entries in the fifteen agreed D05 modules: 248 from the original registry and an additional real/complex representation bridge. Definitions, examples, convention bridges and entries with several proved subresults are counted together. Thus 249 is a coverage count, not a theorem count.

\section*{Independent review and source-led audit}
The initial manuscript was read in four groups by reviewers other than the respective authors. A second pass then worked backwards from the source books, checking the coverage registry itself against central results and mathematically significant exercises. It found and repaired substantive omissions. These are independent language-model reviews of mathematical text, not human peer review or Lean kernel verification.

\begin{center}
\begin{tabular}{lr}
\toprule
Module group & Current entries\\
\midrule
M05.01--M05.07: linear foundations & 89\\
M05.08--M05.11: operators and forms & 69\\
M05.12--M05.14: metric operators & 53\\
M05.15: multilinear constructions & 38\\
\midrule
Total & 249\\
\bottomrule
\end{tabular}
\end{center}
The source correspondence covers 65 Hoffman--Kunze sections, 35 Axler sections and 174 Greub routing units. These units have different sizes: Greub's final two advanced chapters are whole-chapter boundary screens. This is not an exhaustive count of all theorems or an assertion that every exercise is solved. Retained results, shared foundations, corrected source claims and excluded extensions are distinguished in the detailed records.

\section*{Substantive repairs and additions}
The final review supplied the rectangular rank--minor criterion previously used implicitly in scalar extension; relative cyclic decomposition and a commutant criterion; real/complex representation maps; general sesquilinear triangular reduction; spectral and projection refinements; and the missing tensor, exterior and symmetric bridges. Cauchy--Binet, the Gram--Lagrange identity and general Laplace expansion now have explicit coefficient proofs. Characteristic, finite-dimension, zero-factor and normalization conditions are retained. Several problematic printed source statements were replaced by precise corrected versions with counterexamples recorded in the source audits.

Every mathematical amendment was read independently of its author. The final multilinear supplements and added foundation definitions received separate independent checks. The companion \texttt{coverage-manifest.json} maps entries to sources and review routes. The \texttt{audit} directory preserves both earlier snapshots and a final amendment record with hashes of the exact reviewed files. No accepted D05 finding remains open in that record.

\section*{Closure boundary and validation}
Closure is relative to Chapter~\ref{ch:foundations}: classical set and field foundations, explicitly stated choice for unrestricted bases, the specified real-analysis and complex-polar facts, and measure foundations for volume. The required polynomial algebra, fundamental theorem of algebra and linear volume-change rule are proved from that interface. D05 does not reconstruct all of set theory, analysis or measure theory.

Structural checks cover unique entry IDs, resolved references, source-unit correspondence and unchanged source-PDF hashes. The modular project and standalone flattened input have the same mathematical content. Compilation, rendered-page review and delivery hashes are recorded separately. These checks establish document integrity; they do not replace mathematical review. No Lean proof for this manuscript has yet been generated or accepted.

The files \texttt{FINAL\_REVIEW.txt} and \texttt{SOURCE\_COVERAGE.txt} explain the audit and source boundaries; \texttt{source-coverage.json} gives a machine-readable counterpart. The earlier organization-stage catalogue remains preserved as historical data. A future formalization must retain the exact hypotheses and convention bridges and pass the normal proof-acceptance process.

% END verification-register.tex

\end{document}
